% Le moteur électrique — PCT 3ème — Cours signé OnBuch+
% Programme officiel MINESEC. Compiler avec : tectonic N20-moteur-electrique.tex
\newcommand{\DOCMATIERE}{PCT}
\newcommand{\DOCNIVEAU}{3ème}
\newcommand{\DOCMODULE}{Module 4 : Technologie et mécanique}
\newcommand{\DOCLECON}{N20}
\newcommand{\DOCTITRE}{Le moteur électrique}
% OnBuch+ — préambule « cours signés » ENRICHI (compilé avec tectonic / XeLaTeX)
% Système visuel « Pop » d'OnBuch+ : fond crème (jamais blanc), bordures encre
% épaisses, ombres « dures » décalées, titres Archivo Black en capitales.
% Version MATHÉMATIQUES : graphes (pgfplots/tikz), boîtes pédagogiques
% (définition, propriété, méthode, attention, le savais-tu, exercice résolu,
% exercice, TP, bilan), page de garde et signature OnBuch+ ancrée en bas.
%
% Macros attendues dans main.tex AVANT \input{preamble} :
%   \DOCMATIERE  \DOCNIVEAU  \DOCMODULE  \DOCLECON (numéro)  \DOCTITRE
\documentclass[11pt]{article}

\usepackage{fontspec}
\usepackage[french]{babel}
\usepackage[a4paper,margin=2cm,top=3.4cm,bottom=2.8cm,headheight=1.4cm,footskip=1.3cm]{geometry}
\usepackage{amsmath,amssymb,amsfonts}
\usepackage{mathtools} % \xmapsto, \coloneqq... souvent utilisés par les modèles
\usepackage{cancel} % simplifications barrées (bilans de forces)
% \abs{z} (module, valeur absolue) et \norm{u} (norme), utilisés par les modèles
\providecommand{\abs}{}\let\abs\relax\DeclarePairedDelimiter{\abs}{\lvert}{\rvert}
\providecommand{\norm}{}\let\norm\relax\DeclarePairedDelimiter{\norm}{\lVert}{\rVert}
\usepackage[table]{xcolor} % [table] : fournit aussi \rowcolors (alternance de lignes)
\usepackage{pagecolor}
\usepackage{tikz}
\usetikzlibrary{arrows.meta,positioning,calc,shapes.geometric,shapes.misc,shapes.arrows,decorations.pathmorphing,decorations.pathreplacing,decorations.markings,patterns,shadows,mindmap,trees,fit,backgrounds,matrix,intersections,angles,quotes}
\usepackage{pgfplots}
\usepackage[version=4]{mhchem} % équations nucléaires \ce{^{235}_{92}U}
\usepackage[european,siunitx]{circuitikz} % schémas électriques
\pgfplotsset{compat=1.18}
\usepgfplotslibrary{fillbetween,groupplots} % aires entre courbes (calcul intégral)
\usepackage{enumitem}
\usepackage{fancyhdr}
% [most] charge toutes les bibliothèques tcolorbox (listings, minted, raster,
% documentation...) et gonfle inutilement la mémoire TeX sur un document long
% (~300 boîtes) : on ne charge que ce qui est réellement utilisé.
\usepackage[skins,breakable]{tcolorbox}
\usepackage{graphicx}
\usepackage{colortbl}
\usepackage{booktabs}
\usepackage{tabularx}
\usepackage{multirow}
\usepackage{diagbox} % case d'en-tête diagonale des tableaux à double entrée
% Colonnes extensibles centrée / gauche / droite (C, L, R), souvent utilisées par les modèles
\newcolumntype{C}{>{\centering\arraybackslash}X}
\newcolumntype{L}{>{\raggedright\arraybackslash}X}
\newcolumntype{R}{>{\raggedleft\arraybackslash}X}
\usepackage{array}
\usepackage{multicol}
\usepackage{siunitx}
% parse-numbers=false : accepte aussi \SI{2,5 \times 10^{-3}}{...}
\sisetup{locale=FR,output-decimal-marker={,},group-minimum-digits=4,parse-numbers=false}
\DeclareSIUnit\ohm{\ensuremath{\symup{\Omega}}} % U+2126 absent de la police
\DeclareSIUnit\division{div} % réglages d'oscilloscope (ms/div, V/div)
\DeclareSIUnit\tr{tr} % tours (vitesse de rotation, ex. \SI{1500}{\tr\per\minute})
\DeclareSIUnit\molar{M} % concentration molaire (ex. \SI{3}{\molar}, \SI{1}{\micro\molar})
\DeclareSIUnit\an{an} % année (le modèle confond parfois avec \year, un compteur TeX) — ex. \SI{5}{\kilo\meter\per\an}
\DeclareSIUnit\FCFA{FCFA} % franc CFA (ex. \SI{500}{\FCFA\per\kilo\gram})
\DeclareSIUnit\degreeBrix{\textdegree Brix} % teneur en sucre d'un sirop/jus (réfractométrie)
\DeclareSIUnit\month{mois} % le modèle utilise parfois \month, un compteur TeX, comme unité de durée
\DeclareSIUnit\microliter{\micro\liter} % le modèle écrit parfois \microliter en un seul token
\DeclareSIUnit\million{millions} % dénombrement (ex. \SI{15}{\million\per\mL} spermatozoïdes)
\usepackage{qrcode}
% \degree : commande fréquemment utilisée par les modèles pour un angle ou une
% température en dehors de \SI{}{} (non fournie par les paquets chargés).
\providecommand{\degree}{^{\circ}}
% \qed (fin de démonstration), utilisé par les modèles sans amsthm
\providecommand{\qed}{\hfill$\square$}
% \barre : double barre des valeurs interdites dans un tableau de variations (tkz-tab)
\providecommand{\barre}{\ensuremath{\|}}
% environnement proof (amsthm non chargé) utilisé par les modèles
\ifdefined\proof\else\newenvironment{proof}[1][Démonstration]{\par\noindent\textit{#1.}\ }{\qed\par}\fi
\usepackage{lastpage}
\usepackage{hyperref}
\hypersetup{hidelinks}

% ============================================================
% Palette « Pop » OnBuch+ — mêmes hex que la classe P (plus_ui.dart)
% ============================================================
\definecolor{poporange}{HTML}{F59321}
\definecolor{popdark}{HTML}{E07A0C}
\definecolor{popcream}{HTML}{FFF6E8}
\definecolor{popink}{HTML}{1C1714}
\definecolor{poppurple}{HTML}{7A5AE0}
\definecolor{popgreen}{HTML}{2E9E6A}
\definecolor{poppink}{HTML}{E0457B}
\definecolor{popblue}{HTML}{2F7DE1}
\definecolor{popgold}{HTML}{C98A12}
\definecolor{popmuted}{HTML}{8A7F76}
\definecolor{popline}{HTML}{EADFD2}
\definecolor{popgray}{HTML}{8A7F76}
\colorlet{popgrey}{popgray}
% Alias tolérants (le modèle nomme parfois les couleurs différemment)
\colorlet{popwhite}{white}
\colorlet{popblack}{popink}
\colorlet{popred}{poppink}
\colorlet{popyellow}{popgold}
% Teintes claires (fonds de tableaux/encadrés) — le modèle nomme parfois
% les couleurs avec un suffixe L (ex. popblueL) sans les avoir définies.
\colorlet{poporangeL}{poporange!20}
\colorlet{popdarkL}{popdark!20}
\colorlet{poppurpleL}{poppurple!20}
\colorlet{popgreenL}{popgreen!20}
\colorlet{poppinkL}{poppink!20}
\colorlet{popblueL}{popblue!20}
\colorlet{popgoldL}{popgold!20}
\colorlet{popredL}{popred!20}
\colorlet{popgrayL}{popgray!20}
% Teintes claires pour fonds de tableaux / zones
\colorlet{popblueL}{popblue!10!white}
\colorlet{poporangeL}{poporange!14!white}
\colorlet{popgreenL}{popgreen!12!white}
\colorlet{poppurpleL}{poppurple!12!white}
\colorlet{poppinkL}{poppink!12!white}
\colorlet{popgoldL}{popgold!14!white}

\pagecolor{popcream}

% ============================================================
% Typographie
% ============================================================
\defaultfontfeatures{Ligatures=TeX}
\setmainfont{PlusJakartaSans-Medium.ttf}[Path=../../../fonts/,BoldFont=PlusJakartaSans-Bold.ttf]
% Maths en sans-serif (Fira Math) avec chiffres et lettres droites de Plus
% Jakarta, pour que les formules se fondent dans le texte.
\usepackage[math-style=ISO,bold-style=ISO]{unicode-math}
\setmathfont{FiraMath-Regular.otf}[Path=../../../fonts/]
\setmathfont{PlusJakartaSans-Medium.ttf}[Path=../../../fonts/,range={up/num,up/{Latin,latin},\mathup}]
\usepackage{icomma} % virgule décimale sans espace en mode math
% Caractères mathématiques Unicode absents des polices de texte (Plus Jakarta,
% Archivo Black) : rendus par leur équivalent mathématique, en texte comme en
% formule — sinon ils s'affichent en carré vide (ex. « Arithmétique dans ℤ »).
\usepackage{newunicodechar}
\newunicodechar{ℝ}{\ensuremath{\mathbb{R}}}
\newunicodechar{ℤ}{\ensuremath{\mathbb{Z}}}
\newunicodechar{ℂ}{\ensuremath{\mathbb{C}}}
\newunicodechar{ℕ}{\ensuremath{\mathbb{N}}}
\newunicodechar{ℚ}{\ensuremath{\mathbb{Q}}}
\newunicodechar{𝔻}{\ensuremath{\mathbb{D}}}
\newunicodechar{α}{\ensuremath{\alpha}}
\newunicodechar{β}{\ensuremath{\beta}}
\newunicodechar{γ}{\ensuremath{\gamma}}
\newunicodechar{δ}{\ensuremath{\delta}}
\newunicodechar{ε}{\ensuremath{\varepsilon}}
\newunicodechar{θ}{\ensuremath{\theta}}
\newunicodechar{λ}{\ensuremath{\lambda}}
\newunicodechar{μ}{\ensuremath{\mu}}
\newunicodechar{σ}{\ensuremath{\sigma}}
\newunicodechar{τ}{\ensuremath{\tau}}
\newunicodechar{φ}{\ensuremath{\varphi}}
\newunicodechar{ψ}{\ensuremath{\psi}}
\newunicodechar{ω}{\ensuremath{\omega}}
\newunicodechar{Σ}{\ensuremath{\Sigma}}
% majuscules produites par \MakeUppercase dans les titres (α-aminés -> Α)
\newunicodechar{Α}{\ensuremath{\alpha}}
\newunicodechar{Β}{\ensuremath{\beta}}
\newunicodechar{Γ}{\ensuremath{\Gamma}}
\newunicodechar{Δ}{\ensuremath{\Delta}}
\newunicodechar{Ε}{\ensuremath{\varepsilon}}
\newunicodechar{Θ}{\ensuremath{\Theta}}
\newunicodechar{Λ}{\ensuremath{\Lambda}}
\newunicodechar{Μ}{\ensuremath{\mu}}
\newunicodechar{Τ}{\ensuremath{\tau}}
\newunicodechar{Φ}{\ensuremath{\Phi}}
\newunicodechar{Ψ}{\ensuremath{\Psi}}
\newunicodechar{Ω}{\ensuremath{\Omega}}
\newunicodechar{↦}{\ensuremath{\mapsto}}
\newunicodechar{∈}{\ensuremath{\in}}
\newunicodechar{∉}{\ensuremath{\notin}}
\newunicodechar{∧}{\ensuremath{\wedge}}
\newunicodechar{⇔}{\ensuremath{\Leftrightarrow}}
\newunicodechar{⟺}{\ensuremath{\Longleftrightarrow}}
\newunicodechar{⇒}{\ensuremath{\Rightarrow}}
\newunicodechar{⟹}{\ensuremath{\Longrightarrow}}
\newunicodechar{∘}{\ensuremath{\circ}}
\newunicodechar{∪}{\ensuremath{\cup}}
\newunicodechar{∩}{\ensuremath{\cap}}
\newunicodechar{≡}{\ensuremath{\equiv}}
\newunicodechar{⊂}{\ensuremath{\subset}}
\newunicodechar{⊥}{\ensuremath{\perp}}
\newunicodechar{∃}{\ensuremath{\exists}}
\newunicodechar{∀}{\ensuremath{\forall}}
\newunicodechar{∛}{\ensuremath{\sqrt[3]{\,}}}
\newunicodechar{∜}{\ensuremath{\sqrt[4]{\,}}}
\newunicodechar{⃗}{\ensuremath{{}^{\to}}}
\newunicodechar{ⁿ}{\ifmmode^{n}\else\textsuperscript{\ensuremath{n}}\fi}
\newunicodechar{ˣ}{\ifmmode^{x}\else\textsuperscript{\ensuremath{x}}\fi}
\newunicodechar{ᵇ}{\ifmmode^{b}\else\textsuperscript{\ensuremath{b}}\fi}
\newunicodechar{ʳ}{\ifmmode^{r}\else\textsuperscript{\ensuremath{r}}\fi}
\newunicodechar{ᵘ}{\ifmmode^{u}\else\textsuperscript{\ensuremath{u}}\fi}
\newunicodechar{ʸ}{\ifmmode^{y}\else\textsuperscript{\ensuremath{y}}\fi}
\newunicodechar{ᵃ}{\ifmmode^{a}\else\textsuperscript{\ensuremath{a}}\fi}
\newunicodechar{ᵉ}{\ifmmode^{e}\else\textsuperscript{\ensuremath{e}}\fi}
\newunicodechar{ᵏ}{\ifmmode^{k}\else\textsuperscript{\ensuremath{k}}\fi}
\newunicodechar{ᵖ}{\ifmmode^{p}\else\textsuperscript{\ensuremath{p}}\fi}
\newunicodechar{ᴺ}{\ifmmode^{N}\else\textsuperscript{\ensuremath{N}}\fi}
\newunicodechar{ᶜ}{\ifmmode^{c}\else\textsuperscript{\ensuremath{c}}\fi}
\newunicodechar{ʰ}{\ifmmode^{h}\else\textsuperscript{\ensuremath{h}}\fi}
\newunicodechar{ᵠ}{\ifmmode^{q}\else\textsuperscript{\ensuremath{q}}\fi}
\newunicodechar{ₙ}{\ifmmode_{n}\else\textsubscript{\ensuremath{n}}\fi}
\newunicodechar{ₐ}{\ifmmode_{a}\else\textsubscript{\ensuremath{a}}\fi}
\newunicodechar{ᵢ}{\ifmmode_{i}\else\textsubscript{\ensuremath{i}}\fi}
\newunicodechar{ᵣ}{\ifmmode_{r}\else\textsubscript{\ensuremath{r}}\fi}
\newunicodechar{ₖ}{\ifmmode_{k}\else\textsubscript{\ensuremath{k}}\fi}
\newunicodechar{ᵦ}{\ifmmode_{\beta}\else\textsubscript{\ensuremath{\beta}}\fi}
\newunicodechar{ₓ}{\ifmmode_{x}\else\textsubscript{\ensuremath{x}}\fi}
\newfontfamily\popdisplay{ArchivoBlack-Regular.ttf}[Path=../../../fonts/]
\newfontfamily\poplogo{SpaceGrotesk-Bold.ttf}[Path=../../../fonts/]
\newfontfamily\popsemi{PlusJakartaSans-SemiBold.ttf}[Path=../../../fonts/]
\newfontfamily\popxbold{PlusJakartaSans-ExtraBold.ttf}[Path=../../../fonts/]
\newfontfamily\popxboldls{PlusJakartaSans-ExtraBold.ttf}[Path=../../../fonts/,LetterSpace=3.0]

\setlist[enumerate]{leftmargin=1.7em,itemsep=5pt,topsep=4pt}
\setlist[itemize]{leftmargin=1.7em,itemsep=4pt,topsep=4pt,label={\color{poporange}\textbullet}}
\setlength{\parindent}{0pt}
\setlength{\parskip}{5pt}
\renewcommand{\arraystretch}{1.25}

% pgfplots : style Pop par défaut
\pgfplotsset{
  every axis/.append style={axis line style={popink,line width=1pt},tick style={popink},
    grid style={popline,line width=0.5pt},label style={font=\small},tick label style={font=\footnotesize}},
  popcurve/.style={poporange,line width=1.8pt,no markers,smooth},
}
% Même style disponible dans un \draw TikZ simple (hors pgfplots)
\tikzset{popcurve/.style={poporange,line width=1.8pt}}
\tikzset{popgrid/.style={popline,very thin}}
\colorlet{popcurve}{poporange} % le nom du style est parfois utilisé comme couleur

% ============================================================
% Pill / squiggle
% ============================================================
\newcommand{\poppill}[3]{%
  \tikz[baseline=-0.55ex]\node[draw=popink,line width=1.2pt,rounded corners=2.6pt,%
    fill=#2,inner xsep=6.5pt,inner ysep=3pt]{%
    \popxboldls\fontsize{7}{7}\selectfont\color{#3}\MakeUppercase{#1}};%
}
\newcommand{\popsquiggle}[1]{%
  \tikz[baseline=-0.4ex]{
    \draw[#1,line width=2.6pt,line cap=round]
      plot [smooth] coordinates {(0,0.09) (0.34,0.01) (0.68,0.21) (1.02,0.01) (1.36,0.10)};
  }%
}
% Mot-clé surligné (marqueur orange sous le texte)
\newcommand{\cle}[1]{{\popxbold\color{popdark}#1}}

% ============================================================
% En-tête / pied
% ============================================================
\pagestyle{fancy}
\fancyhf{}
\renewcommand{\headrulewidth}{0pt}
\renewcommand{\footrulewidth}{0pt}
\lhead{\raisebox{-0.15ex}{\poplogo\fontsize{13}{13}\selectfont\color{popink}OnBuch\color{poporange}+}}
\rhead{\poppill{\DOCMATIERE\ \textbullet\ \DOCNIVEAU\ \textbullet\ Leçon \DOCLECON}{white}{popink}}
\lfoot{\popxbold\fontsize{7.6}{7.6}\selectfont\color{popmuted}\MakeUppercase{Cours signé OnBuch+}\ \textbullet\ {\popsemi\DOCTITRE}}
\rfoot{\tikz[baseline=-0.6ex]\node[rounded corners=8.5pt,fill=popink,minimum height=17pt,minimum width=17pt,inner xsep=5pt,inner sep=0]{\popxbold\fontsize{8}{8}\selectfont\color{popcream}\thepage\,/\,\pageref*{LastPage}};}

% ============================================================
% Titres
% ============================================================
\newcommand{\chaptitre}[2]{%
  \begin{tcolorbox}[enhanced,colback=poporange,colframe=popink,boxrule=2.6pt,arc=11pt,
      drop shadow=popink,left=15pt,right=15pt,top=12pt,bottom=11pt,before skip=3pt,after skip=18pt]
    \poppill{#2}{popink}{popcream}\par\vspace{9pt}
    {\raggedright\hyphenpenalty=10000\exhyphenpenalty=10000\popdisplay\fontsize{22}{24}\selectfont\color{popcream}\MakeUppercase{#1}\par}
  \end{tcolorbox}
}
% Section numérotée : pastille orange avec le numéro + titre Archivo
\newcounter{popsec}
\newcommand{\coursec}[1]{%
  \stepcounter{popsec}\setcounter{popsub}{0}%
  \par\vspace{16pt}\noindent
  \begin{minipage}{\linewidth}
  \tikz[baseline=-0.7ex]\node[draw=popink,line width=1.6pt,fill=poporange,rounded corners=4pt,minimum size=22pt,inner sep=0]{\popdisplay\fontsize{12}{12}\selectfont\color{popcream}\thepopsec};%
  \hspace{8pt}{\popdisplay\fontsize{15}{17}\selectfont\color{popink}\MakeUppercase{#1}}\par
  \vspace{4pt}\noindent{\color{poporange}\rule{36pt}{3.6pt}}
  \end{minipage}\par\nopagebreak\vspace{8pt}
}
\newcounter{popsub}
\newcommand{\courssub}[1]{%
  \stepcounter{popsub}%
  \par\vspace{9pt}\noindent{\popxbold\fontsize{11.5}{13}\selectfont\color{popink}{\color{poporange}\thepopsec.\thepopsub}\hspace{6pt}#1}\par\nopagebreak\vspace{4pt}
}

% ============================================================
% Boîtes pédagogiques — même recette : fond blanc, bordure épaisse colorée,
% ombre dure encre, badge plein. Titre optionnel [#1].
% ============================================================
\tcbset{popbase/.style={breakable,enhanced,colback=white,boxrule=2.3pt,arc=9pt,
  shadow={3pt}{-3pt}{0pt}{popink},left=14pt,right=14pt,top=11pt,bottom=11pt,before skip=11pt,after skip=15pt,
  fontupper=\popsemi\fontsize{10.6}{14.5}\selectfont\color{popink},
  fontlower=\popsemi\fontsize{10.6}{14.5}\selectfont\color{popink}}}
% Coche dessinée (le glyphe ✓ n'existe pas dans les polices du document)
\newcommand{\popcheck}[1]{\tikz[baseline=-0.5ex]\draw[#1,line width=1.6pt,line cap=round,line join=round](-0.13,0.0)--(-0.04,-0.09)--(0.13,0.1);}
\providecommand{\coche}{\popcheck{popgreen}}
\providecommand{\resultat}[1]{\par\smallskip\noindent\fcolorbox{popgreen}{popgreenL}{\parbox{\dimexpr\linewidth-2\fboxsep-2\fboxrule\relax}{\textbf{Résultat.} #1}}\par\smallskip}
\newcommand{\popboxhead}[3]{\poppill{#1}{#2}{white}\ifx\relax#3\relax\else\hspace{6pt}{\popxbold\fontsize{10.6}{12}\selectfont\color{popink}#3}\fi\par\vspace{7pt}}

\newtcolorbox{aretenir}[1][]{popbase,colframe=popgreen,before upper={\popboxhead{À retenir}{popgreen}{#1}}}
\newtcolorbox{exemplebox}[1][]{popbase,colframe=poporange,before upper={\popboxhead{Exemple}{poporange}{#1}}}
\newtcolorbox{definition}[1][]{popbase,colframe=popblue,before upper={\popboxhead{Définition}{popblue}{#1}}}
\newtcolorbox{propriete}[1][]{popbase,colframe=poppurple,before upper={\popboxhead{Propriété}{poppurple}{#1}}}
\newtcolorbox{methode}[1][]{popbase,colframe=popgold,before upper={\popboxhead{Méthode}{popgold}{#1}}}
\newtcolorbox{attention}[1][]{popbase,colframe=poppink,before upper={\popboxhead{Attention}{poppink}{#1}}}
\newtcolorbox{experience}[1][]{popbase,colframe=popblue,colback=popblueL,before upper={\popboxhead{Expérience}{popblue}{#1}}}
% « Le savais-tu ? » : carte violette pleine (contexte, culture, Cameroun)
\newtcolorbox{savaistu}[1][]{popbase,colback=poppurple,colframe=popink,
  fontupper=\popsemi\fontsize{10.4}{14.2}\selectfont\color{white},
  before upper={\poppill{Le savais-tu\,?}{white}{poppurple}\ifx\relax#1\relax\else\hspace{6pt}{\popxbold\color{white}#1}\fi\par\vspace{7pt}}}
% Exercice résolu : énoncé en haut, \tcblower puis la solution sur fond crème
\newcounter{popexo}\newcounter{popexr}
\newtcolorbox{exoresolu}[1][]{popbase,colframe=poporange,colbacklower=popcream,
  segmentation style={popink,line width=1.2pt,dashed},
  before upper={\stepcounter{popexr}\popboxhead{Exercice résolu \thepopexr}{poporange}{#1}},
  before lower={\poppill{Solution}{popink}{popcream}\par\vspace{6pt}}}
% Exercice (non résolu en place) — la correction est dans la partie Corrigés
\newtcolorbox{exercice}[1][]{popbase,colframe=popink,
  before upper={\stepcounter{popexo}\popboxhead{Exercice \thepopexo}{popink}{#1}}}
% Corrigé
\newtcolorbox{corrige}[1][]{popbase,colframe=popgreen,colback=popgreenL,
  before upper={\popboxhead{Corrigé}{popgreen}{#1}}}
% Objectifs / prérequis / bilan
\newtcolorbox{objectifs}{popbase,colframe=popink,colback=poporangeL,
  before upper={\popboxhead{Objectifs de la leçon}{popink}{}}}
\newtcolorbox{prerequis}{popbase,colframe=popink,colback=popblueL,
  before upper={\popboxhead{Prérequis}{popblue}{}}}
\newtcolorbox{bilan}{popbase,colframe=popink,colback=white,boxrule=2.8pt,
  before upper={\popboxhead{Fiche bilan}{poporange}{L'essentiel à connaître pour l'examen}}}
% Figure encadrée avec légende
\newtcolorbox{popfigure}[1][]{enhanced,breakable=false,colback=white,colframe=popline,boxrule=1.4pt,arc=8pt,
  left=10pt,right=10pt,top=10pt,bottom=8pt,before skip=10pt,after skip=12pt,halign=center,
  lower separated=false,halign lower=center,
  fontlower=\popsemi\fontsize{9}{11.5}\selectfont\color{popmuted},
  #1}
\newcommand{\legende}[1]{\par\vspace{6pt}{\popsemi\fontsize{9}{11.5}\selectfont\color{popmuted}#1}}

% ============================================================
% Signature OnBuch+
% ============================================================
\newcommand{\poplogoinline}[1]{{\poplogo\fontsize{#1}{#1}\selectfont\color{popink}OnBuch\color{poporange}+}}
\newcommand{\signatureonbuch}{%
  \begin{tcolorbox}[enhanced,colback=popink,colframe=popink,boxrule=2.2pt,arc=10pt,
      drop shadow=poporange,left=14pt,right=14pt,top=10pt,bottom=10pt,before skip=0pt,after skip=0pt]
    \tikz[baseline=-0.7ex]\node[circle,fill=popgreen,draw=popcream,line width=1.3pt,minimum size=22pt,inner sep=0]{\popcheck{white}};%
    \hspace{9pt}\begin{minipage}[c]{0.62\linewidth}
      {\popxbold\fontsize{12}{13}\selectfont\color{popcream}Signé\ \textbullet\ {\poplogo OnBuch\color{poporange}+}}\par\vspace{2pt}
      {\raggedright\popsemi\fontsize{8}{10.5}\selectfont\color{popline}\MakeUppercase{Équipe pédagogique OnBuch+}\\\MakeUppercase{Conforme au programme officiel MINESEC}\par}
    \end{minipage}\hfill
    {\poplogo\fontsize{22}{22}\selectfont\color{popcream}OnBuch\color{poporange}+}
  \end{tcolorbox}%
}

% ============================================================
% Page de garde
% #1 titre · #2 matière · #3 niveau · #4 module · #5 n° leçon · #6 durée ·
% #7 description / objectifs (texte ou itemize)
% ============================================================
\newcommand{\pagegarde}[7]{%
  \thispagestyle{empty}
  \newgeometry{a4paper,margin=1.7cm,top=1.6cm,bottom=1.4cm}%
  \begin{tikzpicture}[remember picture,overlay]
    \fill[poporange!18] ([xshift=-3.2cm,yshift=-3.6cm]current page.north east) circle (4.6cm);
    \fill[poppurple!14] ([xshift=2.4cm,yshift=7.6cm]current page.south west) circle (3.8cm);
  \end{tikzpicture}%
  {\poplogo\fontsize{30}{30}\selectfont\color{popink}OnBuch\color{poporange}+}\hfill
  \raisebox{0.6ex}{\poppill{Cours signé \textbullet\ Premium}{popink}{popcream}}\par
  \vspace{1pt}\popsquiggle{poppurple}\par
  \vspace{8pt}
  \begin{tcolorbox}[enhanced,colback=poporange,colframe=popink,boxrule=2.8pt,arc=13pt,
      drop shadow=popink,left=18pt,right=18pt,top=12pt,bottom=13pt,after skip=10pt]
    \poppill{#2 \textbullet\ #3}{popink}{popcream}\hspace{5pt}\poppill{#4}{white}{popink}\par\vspace{8pt}
    {\popdisplay\fontsize{14}{14}\selectfont\color{popink}LEÇON #5}\par\vspace{4pt}
    {\raggedright\hyphenpenalty=10000\exhyphenpenalty=10000\popdisplay\fontsize{25}{27}\selectfont\color{popcream}\MakeUppercase{#1}\par}
    \vspace{8pt}
    {\popxbold\fontsize{9.5}{9.5}\selectfont\color{popink}\MakeUppercase{Durée conseillée : #6}}
  \end{tcolorbox}
  \begin{tcolorbox}[enhanced,colback=white,colframe=popink,boxrule=2.2pt,arc=10pt,
      drop shadow=popink,left=16pt,right=16pt,top=10pt,bottom=10pt,after skip=8pt]
    \poppill{Dans cette leçon}{poppurple}{white}\par\vspace{5pt}
    \popsemi\fontsize{9.3}{12.6}\selectfont\color{popink}#7
  \end{tcolorbox}
  \noindent\begin{minipage}[t]{0.487\linewidth}
    \begin{tcolorbox}[enhanced,colback=popgreen,colframe=popink,boxrule=2.2pt,arc=10pt,
        drop shadow=popink,left=11pt,right=11pt,top=10pt,bottom=10pt,height=3.1cm,valign=center]
      \tcbox[colback=white,colframe=popink,boxrule=1.3pt,arc=5pt,boxsep=0pt,nobeforeafter,
        left=3pt,right=3pt,top=3pt,bottom=3pt,valign=center]{\qrcode[height=1.9cm]{https://play.google.com/store/apps/details?id=com.luvvix.onbuch&hl=fr}}%
      \hspace{8pt}\begin{minipage}[c]{0.5\linewidth}
        {\popdisplay\fontsize{11}{12.5}\selectfont\color{white}\MakeUppercase{Télécharge OnBuch}}\par\vspace{4pt}
        {\raggedright\popsemi\fontsize{8.4}{11}\selectfont\color{white}Gratuit \textbullet\ Google Play\\Tuteur IA, annales, concours, résultats\par}
      \end{minipage}
    \end{tcolorbox}
  \end{minipage}\hfill
  \begin{minipage}[t]{0.487\linewidth}
    \begin{tcolorbox}[enhanced,colback=poppurple,colframe=popink,boxrule=2.2pt,arc=10pt,
        drop shadow=popink,left=11pt,right=11pt,top=10pt,bottom=10pt,height=3.1cm,valign=center]
      \tikz[baseline=-0.7ex]\node[circle,fill=white,draw=popink,line width=1.3pt,minimum size=1.6cm,inner sep=0]{\popdisplay\fontsize{24}{24}\selectfont\color{poppurple}+};%
      \hspace{8pt}\begin{minipage}[c]{0.6\linewidth}
        {\popdisplay\fontsize{11}{12.5}\selectfont\color{white}\MakeUppercase{Passe à OnBuch+}}\par\vspace{4pt}
        {\raggedright\popsemi\fontsize{8.4}{11}\selectfont\color{white}Tous les cours signés, Tuteur IA illimité, fiches PDF — onglet OnBuch+ de l'app\par}
      \end{minipage}
    \end{tcolorbox}
  \end{minipage}
  \vspace{8pt}
  \popcontenu
  \vfill
  \signatureonbuch
  \restoregeometry
  \newpage
}

% Rangée « Ce cours contient » (page de garde)
\newcommand{\poptile}[3]{% #1 couleur #2 gros texte #3 légende
  \begin{tcolorbox}[enhanced,colback=#1,colframe=popink,boxrule=2pt,arc=9pt,shadow={2.5pt}{-2.5pt}{0pt}{popink},
      left=6pt,right=6pt,top=9pt,bottom=9pt,halign=center,valign=center,height=1.75cm,width=\linewidth,nobeforeafter]
    {\popdisplay\fontsize{12.5}{13}\selectfont\color{white}\MakeUppercase{#2}}\par\vspace{3pt}
    {\centering\popsemi\fontsize{7.8}{9.5}\selectfont\color{white}#3\par}
  \end{tcolorbox}}
\newcommand{\popcontenu}{%
  \par\noindent{\popxboldls\fontsize{7.5}{7.5}\selectfont\color{popmuted}CE COURS SIGNÉ CONTIENT}\par\vspace{7pt}
  \noindent\begin{minipage}[t]{0.235\linewidth}\poptile{popblue}{Cours}{complet, illustré, pas à pas}\end{minipage}\hfill
  \begin{minipage}[t]{0.235\linewidth}\poptile{poporange}{Méthodes}{et exercices résolus}\end{minipage}\hfill
  \begin{minipage}[t]{0.235\linewidth}\poptile{poppink}{Exercices}{type examen + corrigés}\end{minipage}\hfill
  \begin{minipage}[t]{0.235\linewidth}\poptile{popgreen}{Fiche bilan}{l'essentiel à retenir}\end{minipage}\par}

% Sommaire
\newtcolorbox{sommaire}{enhanced,colback=white,colframe=popink,boxrule=2.2pt,arc=10pt,
  drop shadow=popink,left=18pt,right=18pt,top=13pt,bottom=13pt,before skip=6pt,after skip=20pt,
  before upper={\poppill{Sommaire}{poppurple}{white}\par\vspace{9pt}\popsemi\fontsize{10.6}{14.5}\selectfont\color{popink}}}

% Signature de fin de cours — poussée tout en bas de la dernière page
\newcommand{\findecours}{%
  \par\vfill\nopagebreak
  \begin{center}\popsquiggle{poporange}\end{center}\vspace{4pt}
  \signatureonbuch
}
\AtBeginDocument{\def\llbracket{\mathopen{[\mkern-3.5mu[}}\def\rrbracket{\mathclose{]\mkern-3.5mu]}}} % intervalles d'entiers (glyphe ⟦ absent de la police)
% Glyphes absents de Fira Math : redéfinis à partir de glyphes disponibles
\AtBeginDocument{%
  \def\setminus{\mathbin{\backslash}}\let\smallsetminus\setminus
  \def\vdots{\mathord{\vbox{\baselineskip3.2pt\lineskiplimit0pt\kern4pt\hbox{.}\hbox{.}\hbox{.}}}}%
  \def\ddots{\mathinner{\mkern1mu\raise6.5pt\hbox{.}\mkern2mu\raise3.6pt\hbox{.}\mkern2mu\raise0.7pt\hbox{.}\mkern1mu}}%
  \def\triangle{\mathord{\scalebox{0.9}{$\symup{\Delta}$}}}%
  \def\nabla{\mathord{\raisebox{\depth}{\scalebox{1}[-1]{$\symup{\Delta}$}}}}%
  \def\checkmark{\popcheck{popdark}}%
  \def\star{\mathbin{\ast}}\def\bigstar{\mathord{\ast}}%
  \def\diamond{\mathbin{\scalebox{0.6}{$\Diamond$}}}%
  \def\bigcup{\mathop{\mathchoice{\raisebox{-0.3em}{\scalebox{1.7}{$\cup$}}}{\scalebox{1.25}{$\cup$}}{\cup}{\cup}}\displaylimits}%
  \def\bigcap{\mathop{\mathchoice{\raisebox{-0.3em}{\scalebox{1.7}{$\cap$}}}{\scalebox{1.25}{$\cap$}}{\cap}{\cap}}\displaylimits}%
  \def\lesseqqgtr{\mathrel{\lessgtr}}\def\bigoplus{\mathop{\oplus}\displaylimits}%
}
\providecommand{\ssi}{si et seulement si} % abréviation fréquente
\AtBeginDocument{\providecommand{\wideparen}{\overparen}} % arc de cercle

%% FIGFIX-BEGIN v5 — correctifs globaux des figures (docs/onbuch-generation/tools/figfix)
\usepackage{adjustbox}\usepackage{etoolbox}\tcbuselibrary{raster}
% toute figure TikZ/pgfplots plus large que la ligne est réduite à la largeur disponible
\BeforeBeginEnvironment{tikzpicture}{\begin{adjustbox}{max width=\linewidth}}
\AfterEndEnvironment{tikzpicture}{\end{adjustbox}}
% légendes pgfplots lisibles par-dessus les courbes
\pgfplotsset{every axis/.append style={legend style={fill=white,fill opacity=0.92,text opacity=1,draw=black!15,inner sep=3pt}}}
% étiquettes d'arêtes (syntaxe edge["..."]) avec fond blanc
\tikzset{every edge quotes/.append style={fill=white,inner sep=1.2pt}}
%% FIGFIX-END
\begin{document}

\pagegarde{Le moteur électrique}{PCT}{3ème}{Module 4 : Technologie et mécanique}{N20}{1 séance (1 h chacune)}{Cette leçon présente le moteur électrique, ses constituants essentiels (stator, rotor, aimant, bobine) et son principe de fonctionnement fondé sur les forces électromagnétiques. À travers des expériences simples et des situations de la vie courante camerounaise, l'élève apprend à identifier, décrire et utiliser un moteur en respectant les consignes de sécurité.
\vspace{4pt}
\begin{multicols}{2}\small
\begin{itemize}
\item À la fin de cette leçon, je sais définir un moteur électrique et citer des exemples d'appareils à moteur utilisés au quotidien
\item À la fin de cette leçon, je sais identifier les constituants essentiels d'un moteur électrique à courant continu (stator, rotor, bobine, aimant, collecteur, balais)
\item À la fin de cette leçon, je sais décrire expérimentalement l'action d'un aimant sur une bobine parcourue par un courant
\item À la fin de cette leçon, je sais expliquer le principe de fonctionnement du moteur électrique (conversion d'énergie électrique en énergie mécanique)
\item À la fin de cette leçon, je sais réaliser et interpréter l'expérience du moteur élémentaire (spire mobile dans un champ magnétique)
\item À la fin de cette leçon, je sais établir la chaîne énergétique d'un appareil à moteur (ventilateur, pompe, perceuse)
\end{itemize}\end{multicols}}

\begin{sommaire}
\begin{multicols}{2}
\begin{enumerate}
\item Le moteur électrique dans la vie courante
\item Étude technique : constitution d'un moteur électrique à courant continu
\item Étude théorique : action d'un aimant sur une bobine parcourue par un courant
\item Principe de fonctionnement du moteur électrique
\item Utilisation, plaque signalétique et sécurité
\item Activité d'intégration
\item Méthodes et exercices résolus
\item Exercices
\item Corrigés des exercices
\item Fiche bilan
\end{enumerate}
\end{multicols}
\end{sommaire}

\begin{objectifs}
\begin{itemize}
\item À la fin de cette leçon, je sais définir un moteur électrique et citer des exemples d'appareils à moteur utilisés au quotidien
\item À la fin de cette leçon, je sais identifier les constituants essentiels d'un moteur électrique à courant continu (stator, rotor, bobine, aimant, collecteur, balais)
\item À la fin de cette leçon, je sais décrire expérimentalement l'action d'un aimant sur une bobine parcourue par un courant
\item À la fin de cette leçon, je sais expliquer le principe de fonctionnement du moteur électrique (conversion d'énergie électrique en énergie mécanique)
\item À la fin de cette leçon, je sais réaliser et interpréter l'expérience du moteur élémentaire (spire mobile dans un champ magnétique)
\item À la fin de cette leçon, je sais établir la chaîne énergétique d'un appareil à moteur (ventilateur, pompe, perceuse)
\item À la fin de cette leçon, je sais appliquer les consignes de sécurité liées à l'utilisation des moteurs électriques
\item À la fin de cette leçon, je sais interpréter la plaque signalétique d'un moteur (tension, puissance)
\end{itemize}
\end{objectifs}

\begin{prerequis}
\begin{itemize}
\item Le courant électrique continu : intensité, tension, circuit simple (leçons d'électricité de 3ème)
\item Les aimants et les actions magnétiques (pôles, attraction, répulsion)
\item Les formes d'énergie et les conversions d'énergie (notion vue en 4ème et rappelée en début d'année)
\item Les forces et leurs effets : mise en mouvement, rotation (module mécanique)
\item La puissance électrique P = U × I (notion d'électricité)
\end{itemize}
\end{prerequis}

\newpage

\coursec{Le moteur électrique dans la vie courante}

À Yaoundé comme à Douala, le ventilateur devient indispensable pendant la saison chaude, le moulin à maïs du quartier tourne toute la journée pour moudre les grains, et la pompe à eau remplit les réservoirs des familles. Tous ces appareils ont un point commun : un moteur électrique qui transforme l'énergie électrique en mouvement. Mais comment un simple courant électrique peut-il faire tourner un axe ? Pourquoi le moteur chauffe-t-il et peut-il griller si on le bloque ? Cette leçon répond à ces questions en dévoilant le principe de fonctionnement du moteur électrique.

\courssub{Recensement des appareils à moteur dans la vie courante}

\textbf{Observation autour de nous.}\par

Faisons ensemble une enquête dans notre environnement quotidien. À la maison, au marché, à l'atelier du quartier, de nombreux appareils produisent un mouvement dès qu'on les branche sur une prise de courant ou qu'on les alimente en électricité. Voici ce que l'on observe :

\begin{itemize}
\item le \cle{ventilateur} : ses pales tournent et brassent l'air ;
\item la \cle{pompe à eau} : elle aspire l'eau du forage ou du puits et la refoule vers le réservoir ;
\item le \cle{moulin à maïs} du quartier : il écrase les grains pour produire la farine ;
\item la \cle{perceuse} de l'atelier de menuiserie : son foret tourne pour percer le bois ou le métal ;
\item le \cle{mixeur} de la cuisine : ses couteaux tournent très vite pour broyer les fruits et légumes ;
\item la \cle{moto-pompe} : elle puise l'eau pour irriguer les champs ;
\item la \cle{machine à coudre électrique} du tailleur : son aiguille monte et descend grâce à un axe en rotation.
\end{itemize}

\begin{experience}[Enquête : les appareils du quartier]
\textbf{Matériel :} un cahier, un crayon, et l'observation des appareils de la maison et du quartier (avec l'accord des adultes).

\textbf{Protocole :}
\begin{enumerate}
\item Dresser la liste des appareils électriques utilisés à la maison et dans le quartier.
\item Pour chaque appareil, noter s'il produit un mouvement (rotation d'un axe, va-et-vient) ou non.
\item Classer les appareils en deux colonnes : « avec moteur » et « sans moteur ».
\end{enumerate}

\textbf{Observations :} le ventilateur, le mixeur, la perceuse, la pompe à eau et la machine à coudre produisent un mouvement de rotation. En revanche, la lampe, le fer à repasser, la bouilloire et le chargeur de téléphone ne produisent aucun mouvement : ils produisent de la lumière ou de la chaleur.

\textbf{Interprétation :} les appareils qui produisent un mouvement contiennent tous un organe commun : le \cle{moteur électrique}. C'est lui qui met en rotation l'axe de l'appareil.

\textbf{Conclusion :} un appareil électrique qui met quelque chose en mouvement possède un moteur électrique.
\end{experience}

\begin{center}
\begin{tabularx}{\linewidth}{|X|X|}
\hline
\rowcolor{popblueL} \textbf{Appareils AVEC moteur (mouvement)} & \textbf{Appareils SANS moteur (pas de mouvement)} \\
\hline
Ventilateur, pompe à eau, moulin à maïs, perceuse, mixeur, moto-pompe, machine à coudre électrique, réfrigérateur (compresseur), essoreuse & Lampe, fer à repasser, bouilloire, grille-pain, chargeur de téléphone, téléviseur, radio, cuisinière électrique \\
\hline
\end{tabularx}
\end{center}

\courssub{Le constat commun : de l'électricité au mouvement}

\textbf{Ce que tous ces appareils partagent.}\par

Reprenons notre liste. Le ventilateur, la pompe, le moulin, la perceuse, le mixeur : tous sont \textbf{branchés sur une source électrique} (le réseau ENEO, un groupe électrogène, une batterie ou un panneau solaire) et tous \textbf{produisent un mouvement}, le plus souvent la \textbf{rotation d'un axe}.

\begin{aretenir}[Constat fondamental]
Un appareil contenant un moteur électrique, lorsqu'il est alimenté par une source d'électricité, produit un mouvement, généralement la \cle{rotation d'un axe}. Le moteur électrique est donc un \textbf{convertisseur d'énergie} : il reçoit de l'énergie électrique et la transforme en énergie mécanique.
\end{aretenir}

\begin{definition}[Moteur électrique]
Un \cle{moteur électrique} est un appareil qui \textbf{convertit l'énergie électrique en énergie mécanique}, c'est-à-dire en mouvement (le plus souvent un mouvement de rotation d'un axe).
\end{definition}

On peut schématiser cette conversion par une \textbf{chaîne énergétique} :

\begin{popfigure}
\begin{tikzpicture}[font=\small]
% Source électrique
\draw[fill=popgoldL, draw=popdark, rounded corners=3pt] (0,0) rectangle (2.6,1.4);
\node[align=center] at (1.3,0.7) {\textbf{Source}\\ \textbf{électrique}\\ (ENEO, pile,\\ batterie)};
% Flèche 1
\draw[-{Stealth[length=3mm]}, very thick, poporange] (2.7,0.7) -- (4.1,0.7);
\node[above, poporange] at (3.4,0.7) {\footnotesize énergie électrique};
% Moteur
\draw[fill=popblueL, draw=popdark, rounded corners=3pt] (4.2,0) rectangle (6.4,1.4);
\node[align=center] at (5.3,0.7) {\textbf{Moteur}\\ \textbf{électrique}};
% Flèche 2
\draw[-{Stealth[length=3mm]}, very thick, popgreen] (6.5,0.7) -- (7.9,0.7);
\node[above, popgreen!60!black] at (7.2,0.7) {\footnotesize énergie mécanique};
% Axe en rotation
\draw[fill=popgreenL, draw=popdark, rounded corners=3pt] (8,0) rectangle (10.4,1.4);
\node[align=center] at (9.2,0.7) {\textbf{Axe en}\\ \textbf{rotation}\\ (pales du\\ ventilateur)};
% Petite flèche de rotation
\draw[-{Stealth[length=2mm]}, thick, popdark] (9.9,1.6) arc (60:-60:0.35);
\end{tikzpicture}
\legende{Chaîne énergétique simplifiée d'un ventilateur : la source électrique fournit l'énergie électrique, le moteur la convertit en énergie mécanique qui fait tourner l'axe et les pales.}
\end{popfigure}

\begin{methode}[Calculer le courant consommé par un moteur]
\textbf{Données :} Puissance $P$ (en watts, \si{W}) et tension $U$ (en volts, \si{V}) lues sur la plaque signalétique.
\begin{enumerate}
\item Écrire la relation : $P = U \times I$.
\item Isoler l'intensité : $I = \dfrac{P}{U}$.
\item Effectuer le calcul et exprimer le résultat en ampères (\si{A}).
\end{enumerate}
\end{methode}

\begin{exemplebox}[Un calcul simple : le ventilateur de la maison]
Un ventilateur porte sur sa plaque l'indication « \SI{45}{W} -- \SI{220}{V} ». Quel courant consomme-t-il ?

On utilise la relation de la puissance électrique : $P = U \times I$, donc :
\[ I = \frac{P}{U} = \frac{45}{220} \approx \SI{0,2}{A} \]

Le ventilateur consomme donc un courant d'environ \SI{0,2}{A}. C'est un courant faible : c'est pourquoi un ventilateur peut fonctionner sur une rallonge ordinaire sans danger, contrairement à un fer à repasser de \SI{1000}{W} qui demande environ \SI{4,5}{A}.
\end{exemplebox}

\begin{exoresolu}[Identifier le convertisseur d'énergie]
\textbf{Énoncé :} Mama Brigitte branche son mixeur pour préparer la sauce de folong. a) Quelle forme d'énergie le mixeur reçoit-il ? b) Quelle forme d'énergie produit-il ? c) Comment appelle-t-on l'organe qui réalise cette conversion ?
\tcblower
\textbf{Solution détaillée :}

a) Le mixeur est branché sur la prise de courant : il reçoit de l'\textbf{énergie électrique}.

b) Les couteaux du mixeur se mettent à tourner très vite : le mixeur produit de l'\textbf{énergie mécanique} (mouvement de rotation). On entend aussi du bruit et le moteur chauffe légèrement : une petite partie de l'énergie est perdue en chaleur (rendement inférieur à 100\,\%).

c) L'organe qui convertit l'énergie électrique en énergie mécanique est le \textbf{moteur électrique}.
\end{exoresolu}

\courssub{Moteur ou générateur : ne pas confondre !}

\textbf{Deux machines, deux sens de conversion.}\par

Le moteur électrique a un « cousin » qui fait exactement le travail inverse : la \cle{dynamo} (ou \cle{générateur}, ou \cle{alternateur}). Sur une bicyclette, la dynamo éclaire la lampe quand la roue tourne : c'est le mouvement qui produit l'électricité. Le groupe électrogène du quartier fonctionne sur le même principe : un moteur à essence fait tourner un alternateur qui produit le courant.

\begin{center}
\begin{tabularx}{\linewidth}{|l|X|X|}
\hline
\rowcolor{popblueL} \textbf{Machine} & \textbf{Énergie reçue} & \textbf{Énergie produite} \\
\hline
Moteur électrique & Énergie électrique & Énergie mécanique (mouvement) \\
\hline
Dynamo / générateur / alternateur & Énergie mécanique (mouvement) & Énergie électrique \\
\hline
\end{tabularx}
\end{center}

\begin{attention}[Erreur fréquente : moteur ou générateur ?]
Ne confonds jamais les deux sens de conversion :
\begin{itemize}
\item \textbf{Moteur} : électricité $\rightarrow$ mouvement (il \emph{consomme} de l'électricité). Exemple : le ventilateur, la pompe, le mixeur.
\item \textbf{Générateur (dynamo, alternateur)} : mouvement $\rightarrow$ électricité (il \emph{produit} de l'électricité). Exemple : la dynamo de vélo, l'alternateur du groupe électrogène, l'alternateur de voiture.
\end{itemize}
Au BEPC, une question classique demande : « Le ventilateur contient-il un générateur ? » La réponse est \textbf{non} : il contient un moteur, car il consomme l'électricité pour produire un mouvement. (Ce résultat est admis à ce niveau ; la conversion inverse sera étudiée plus tard.)
\end{attention}

\begin{attention}[Consignes de sécurité liées au moteur]
\begin{itemize}
\item \textbf{Ne jamais bloquer l'axe d'un moteur en marche} (ex. : coincer les pales du ventilateur, forcer sur la perceuse) : le courant augmente brutalement, le moteur surchauffe et peut griller ou provoquer un incendie.
\item \textbf{Ne pas insérer les doigts ni d'objets} près des parties tournantes (pales, courroies, engrenages) : risque de sectionnement ou d'arrachement.
\item \textbf{Vérifier la tension d'alimentation} (souvent \SI{220}{V} au Cameroun) avant de brancher : un moteur \SI{110}{V} branché sur \SI{220}{V} grille instantanément.
\item \textbf{Utiliser des câbles et prises adaptés} à la puissance (section du fil, disjoncteur) pour éviter l'échauffement des fils.
\end{itemize}
\end{attention}

\courssub{Le rôle technologique du moteur électrique dans le développement}

\textbf{Une machine au service du quotidien camerounais.}\par

Le moteur électrique est l'une des inventions qui ont le plus transformé la vie des sociétés. Avant lui, beaucoup de travaux se faisaient à la force des bras ou des animaux. Aujourd'hui :

\begin{itemize}
\item \textbf{Dans l'artisanat :} la perceuse, la scie électrique, la machine à coudre et la meuleuse permettent au menuisier, au tailleur et au soudeur de travailler plus vite et avec plus de précision.
\item \textbf{Dans l'agriculture :} le moulin à maïs, à manioc ou à mil transforme en quelques minutes ce qui demandait des heures de pilage au mortier ; la moto-pompe irrigue les champs et maraîchages pendant la saison sèche.
\item \textbf{Dans l'approvisionnement en eau :} les pompes électriques remplissent les châteaux d'eau et les réservoirs des maisons, là où il faudrait autrement tirer l'eau du puits seau par seau.
\item \textbf{Dans le transport :} les ascenseurs, les trains électriques et maintenant les voitures et motos électriques utilisent des moteurs électriques.
\item \textbf{Dans le confort domestique :} ventilateur, réfrigérateur, machine à laver, mixeur rendent de grands services à la famille.
\end{itemize}

\begin{savaistu}[Le moulin à maïs du quartier]
Le moulin à maïs du quartier utilise un moteur électrique de \textbf{plusieurs kilowatts} (souvent entre \SI{3}{kW} et \SI{10}{kW}), soit une puissance environ \textbf{100 fois plus grande} que celle d'un ventilateur de \SI{45}{W} ! C'est pourquoi le meunier doit utiliser des câbles épais et un disjoncteur adapté : un tel moteur consomme un courant de 15 à \SI{45}{A}, qui ferait fondre une simple rallonge domestique. Le moteur électrique existe dans toutes les tailles : du micromoteur d'un jouet (moins d'un watt) aux moteurs géants des usines de la SONARA ou des barrages hydroélectriques.
\end{savaistu}

\begin{exercice}[À toi de jouer]
\begin{enumerate}
\item Recopie et complète : un moteur électrique convertit l'énergie \ldots\ en énergie \ldots
\item Parmi les appareils suivants, entoure ceux qui contiennent un moteur électrique : bouilloire, perceuse, lampe, pompe à eau, fer à repasser, mixeur.
\item Une moto-pompe porte l'indication « \SI{1100}{W} -- \SI{220}{V} ». Calcule l'intensité du courant qu'elle consomme.
\item La dynamo d'une bicyclette est-elle un moteur ? Justifie ta réponse.
\item Pourquoi ne doit-on pas bloquer l'axe d'un moteur en marche ?
\end{enumerate}
\end{exercice}

\begin{corrige}[Exercice : À toi de jouer]
\begin{enumerate}
\item Un moteur électrique convertit l'énergie \textbf{électrique} en énergie \textbf{mécanique}.
\item Appareils contenant un moteur : la \textbf{perceuse}, la \textbf{pompe à eau} et le \textbf{mixeur} (ils produisent un mouvement). La bouilloire, la lampe et le fer à repasser produisent de la chaleur ou de la lumière, sans mouvement.
\item $I = \dfrac{P}{U} = \dfrac{1100}{220} = \SI{5}{A}$. La moto-pompe consomme un courant de \SI{5}{A}.
\item \textbf{Non.} La dynamo reçoit un mouvement (la rotation de la roue) et produit de l'électricité : c'est un \textbf{générateur} (ou \cle{dynamo}), c'est-à-dire l'inverse d'un moteur.
\item Bloquer l'axe empêche la rotation, ce qui supprime la force contre-électromotrice (phénomène vu plus tard) : l'intensité du courant augmente fortement, le moteur surchauffe (effet Joule) et risque de griller ou de provoquer un incendie.
\end{enumerate}
\end{corrige}

\begin{aretenir}[L'essentiel de la section]
\begin{itemize}
\item De nombreux appareils du quotidien (ventilateur, pompe, moulin, perceuse, mixeur, machine à coudre) contiennent un \cle{moteur électrique}.
\item Le moteur électrique \textbf{convertit l'énergie électrique en énergie mécanique} : alimenté par une source électrique, il produit la rotation d'un axe.
\item Le \cle{générateur} (dynamo, alternateur) fait la conversion inverse : mouvement $\rightarrow$ électricité.
\item Le moteur électrique joue un rôle majeur dans le développement : artisanat, agriculture, eau, transport et confort domestique.
\item \textbf{Sécurité :} ne pas bloquer l'axe, ne pas toucher les parties tournantes, respecter la tension et la section des câbles.
\end{itemize}
\end{aretenir}

\coursec{Étude technique : constitution d'un moteur électrique à courant continu}

Dans la section précédente, nous avons découvert que le moteur électrique est présent partout autour de nous : dans le ventilateur qui nous rafraîchit pendant la saison sèche, dans le mixeur de la cuisine, dans la perceuse du menuisier, dans le démarreur de la moto-taxi. Nous avons aussi constaté son rôle général : transformer l'\cle{énergie électrique} en \cle{énergie mécanique} de rotation. Mais une question reste posée : \emph{que trouve-t-on à l'intérieur de ce petit cylindre mystérieux ?} Pour répondre à cette question, les techniciens procèdent à une \cle{étude technique} : ils démontent l'objet, observent ses pièces, les nomment et cherchent la fonction de chacune. C'est exactement ce que nous allons faire maintenant, comme de vrais ingénieurs.

\courssub{Observation d'un petit moteur démonté : la démarche d'étude technique}

\begin{experience}[Démontage d'un moteur de jouet ou de ventilateur hors d'usage]
\textbf{Matériel :} un petit moteur à courant continu (récupéré sur un jouet cassé, un ventilateur de table hors d'usage ou un lecteur DVD abandonné), un tournevis, une pince, une loupe, une pile plate de \SI{4,5}{V} et deux fils de connexion.

\textbf{Protocole :}
\begin{enumerate}
\item Observer le moteur fermé : noter sa forme, repérer les deux bornes électriques et l'axe qui dépasse.
\item Brancher le moteur sur la pile : constater que l'axe tourne. Inverser les bornes : constater que le sens de rotation change.
\item Débrancher, puis ouvrir délicatement le carter avec le tournevis.
\item Retirer la partie mobile et étaler les pièces sur la table dans l'ordre du démontage.
\item Identifier chaque pièce et noter ses observations dans un tableau.
\end{enumerate}

\textbf{Observations :}
\begin{itemize}
\item À l'intérieur du carter, collées contre la paroi, on trouve des pièces qui attirent un trombone : ce sont des \textbf{aimants}. Ils ne bougent pas.
\item La partie qui tournait est un cylindre autour duquel est enroulé du \textbf{fil de cuivre} verni, en plusieurs bobines.
\item À l'extrémité de la partie mobile, on voit un petit cylindre de cuivre fendu en plusieurs secteurs : c'est le \textbf{collecteur}.
\item Contre le collecteur frottent deux petits blocs noirs reliés aux bornes : ce sont les \textbf{balais} (ou charbons).
\item L'axe central, appelé \textbf{arbre}, dépasse du carter : c'est lui qui entraîne l'hélice ou le galet.
\end{itemize}

\textbf{Interprétation :} le moteur est organisé en deux grandes parties : une partie \emph{fixe} (le carter avec les aimants) et une partie \emph{mobile} (l'ensemble bobines + collecteur + arbre). Le courant arrive aux bobines mobiles grâce au frottement des balais sur le collecteur.

\textbf{Conclusion :} un moteur électrique à courant continu est constitué d'un \cle{stator} (partie fixe), d'un \cle{rotor} (partie mobile), d'un système d'alimentation de la partie tournante (collecteur et balais), d'un arbre de sortie et d'un carter de protection.
\end{experience}

Cette démarche — \emph{observer, démonter méthodiquement, nommer, attribuer une fonction à chaque pièce} — est la démarche d'étude technique. Elle est utilisée par tous les techniciens, du réparateur de motos de quartier à l'ingénieur de la SONARA.

\courssub{Le stator : la partie fixe qui crée le champ magnétique}

\begin{definition}[Stator]
Le \cle{stator} est la partie \textbf{fixe} du moteur électrique. Dans un petit moteur à courant continu, il porte un \textbf{aimant permanent} (ou plusieurs) appelé \cle{inducteur}, qui crée un champ magnétique dans l'espace où tourne le rotor.
\end{definition}

Retiens bien l'intuition : le stator est « l'assise » du moteur, comme le socle d'une statue. Il est solidaire du carter et ne bouge jamais. Son rôle est capital : sans aimant, pas de champ magnétique, et donc — nous le verrons dans l'étude théorique — pas de force sur la bobine, donc pas de rotation.

Dans les petits moteurs de jouets, l'inducteur est un aimant permanent en forme de tuile, collé contre la paroi intérieure du carter, avec un pôle Nord et un pôle Sud face à face. Dans les moteurs plus puissants (démarreur de voiture, moteur industriel), l'aimant est remplacé par un \textbf{électroaimant} : une bobine fixe parcourue par un courant, qui produit un champ magnétique plus intense et réglable.

\begin{attention}[Qui tourne, qui ne tourne pas ?]
Erreur très fréquente au BEPC : croire que \textbf{l'aimant tourne} et que la bobine est fixe. C'est l'inverse ! Dans le moteur à courant continu étudié en classe de 3ème :
\begin{itemize}
\item l'\textbf{aimant} est fixe : il fait partie du \textbf{stator} ;
\item la \textbf{bobine} tourne : elle fait partie du \textbf{rotor}.
\end{itemize}
Moyen mnémotechnique : \emph{« L'aimant est assis (stator), la bobine court (rotor) »}. Attention toutefois : dans certains appareils (alternateur de vélo, éolienne), c'est l'aimant qui tourne devant une bobine fixe. Ce n'est pas le cas de notre moteur : lis toujours bien l'énoncé.
\end{attention}

\courssub{Le rotor : la partie mobile qui tourne}

\begin{definition}[Rotor]
Le \cle{rotor} est la partie \textbf{mobile} du moteur électrique. Il est constitué d'une ou plusieurs \textbf{bobines} de fil de cuivre émaillé, enroulées sur un \textbf{noyau} de fer doux, et montées sur l'arbre de rotation. C'est lui qui tourne lorsque le moteur fonctionne.
\end{definition}

Détaillons les trois constituants du rotor :

\begin{itemize}
\item \textbf{Les bobines.} Elles sont faites de fil de cuivre recouvert d'un vernis isolant (fil « émaillé »), enroulé en de nombreuses spires. Le cuivre est choisi parce qu'il est excellent conducteur ; le vernis empêche les spires de se toucher et de créer un court-circuit. Les petits moteurs de jouets possèdent généralement \textbf{trois bobines} disposées à \ang{120} l'une de l'autre, ce qui rend la rotation régulière.
\item \textbf{Le noyau de fer doux.} C'est un cylindre de tôles de fer empilées sur lequel les bobines sont enroulées. Le fer doux \emph{concentre} le champ magnétique : il rend les forces magnétiques sur la bobine beaucoup plus intenses, donc le moteur plus puissant.
\item \textbf{L'arbre.} C'est la tige d'acier central, solidaire du rotor. Il tourne avec lui et traverse le carter grâce à des paliers (ou roulements) qui réduisent les frottements.
\end{itemize}

\courssub{Le collecteur et les balais : alimenter une pièce qui tourne}

Voici le problème le plus délicat que les inventeurs du moteur ont dû résoudre : \emph{comment amener le courant électrique dans une bobine qui tourne sans cesse ?} Si l'on branchait directement des fils sur la bobine, ils s'enrouleraient et casseraient au bout de quelques tours ! La solution est un système ingénieux : le couple \cle{collecteur}-\cle{balais}.

\begin{definition}[Collecteur et balais]
Le \cle{collecteur} est un petit cylindre de cuivre, solidaire du rotor, divisé en \textbf{lamelles} isolées les unes des autres (deux lamelles pour une bobine). Chaque lamelle est reliée à une extrémité de la bobine. Les \cle{balais} (ou \textbf{charbons}) sont deux petits blocs de graphite, fixes, reliés aux bornes du générateur, qui \textbf{frottent} sur le collecteur et assurent le contact électrique avec la partie tournante.
\end{definition}

Le collecteur a une deuxième fonction, plus subtile et absolument essentielle : il \textbf{inverse le sens du courant dans la bobine à chaque demi-tour}. En effet, quand le rotor a fait un demi-tour, chaque lamelle quitte son balais et touche l'autre balais : le courant change alors de sens dans la bobine. Nous verrons dans l'étude théorique que sans cette inversion, la bobine s'arrêterait en position d'équilibre au lieu de tourner continûment. Le collecteur est donc un \emph{inverseur automatique de courant}, commandé par la rotation elle-même.

\begin{savaistu}[Pourquoi appelle-t-on les balais des « charbons » ?]
Les balais ne sont pas en métal mais en \textbf{graphite}, une forme du carbone — d'où leur nom de « charbons ». Le graphite possède trois qualités remarquables : il \textbf{conduit bien le courant électrique}, il est \textbf{tendre} (il frotte sur le cuivre sans rayer le collecteur) et il est \textbf{autolubrifiant} (il glisse presque sans frottement). C'est le même graphite que celui de la mine de ton crayon ! Les charbons s'usent lentement : c'est pourquoi, sur une perceuse ou une meuleuse âgée, le réparateur doit parfois « changer les charbons ». Quand tu vois des étincelles bleues à l'intérieur d'une perceuse en marche, ce sont les charbons qui travaillent sur le collecteur.
\end{savaistu}

\courssub{L'arbre de sortie et le carter}

\begin{itemize}
\item \textbf{L'arbre de sortie} est l'extrémité de l'axe du rotor qui dépasse du moteur. Il \textbf{transmet le mouvement de rotation} à l'appareil entraîné : hélice du ventilateur, meule de la meuleuse, tambour de la machine à laver (par l'intermédiaire d'une courroie), galet du magnétophone. On peut y fixer une poulie, un pignon ou une manivelle.
\item \textbf{Le carter} est l'enveloppe protectrice du moteur, en métal ou en plastique dur. Il protège les pièces internes de la poussière et des chocs, protège l'utilisateur des pièces tournantes et des parties sous tension, et sert de support au stator et aux paliers de l'arbre.
\end{itemize}

\courssub{Schéma en coupe d'un moteur à courant continu}

\begin{popfigure}
\begin{tikzpicture}[scale=1.0, every node/.style={font=\small}]
% Carter
\draw[fill=popmuted!25, thick, popdark] (-4.6,-2.6) rectangle (4.6,2.6);
\node[popdark, font=\small\bfseries] at (-3.4,2.25) {6. Carter};
% Aimants stator (tuiles)
\draw[fill=popblue!60, thick] (-4.0,0.9) arc[start angle=120, end angle=60, radius=2.3] -- (2.0,1.55) arc[start angle=60, end angle=120, radius=1.55] -- cycle;
\draw[fill=poporange!70, thick] (-4.0,-0.9) arc[start angle=240, end angle=300, radius=2.3] -- (2.0,-1.55) arc[start angle=300, end angle=240, radius=1.55] -- cycle;
\node[white, font=\small\bfseries] at (0,1.85) {N};
\node[white, font=\small\bfseries] at (0,-1.85) {S};
\node[popdark, font=\small] at (-3.1,1.5) {1. Stator};
\node[popdark, font=\small] at (-3.1,1.15) {(aimant)};
% Noyau rotor
\draw[fill=poppurple!35, thick] (0,0) circle (1.05);
\node[popdark, font=\small] at (0,-0.35) {2. Rotor};
\node[popdark, font=\small] at (0,-0.68) {(bobines + noyau)};
% Spires bobine (hachures)
\foreach \a in {30,60,120,150,210,240,300,330}{
  \draw[poppurple, thick] (\a:0.95) circle (0.10);
}
% Arbre
\draw[fill=popdark!70] (-4.6,-0.12) rectangle (4.6,0.12);
\node[popdark, font=\small] at (3.6,0.42) {5. Arbre};
% Collecteur
\draw[fill=popgold!80, thick] (2.6,-0.45) rectangle (3.2,0.45);
\draw[popdark, thick] (2.9,-0.45) -- (2.9,0.45);
\node[popdark, font=\small] at (2.9,-0.85) {3. Collecteur};
% Balais
\draw[fill=popgreen!60!black, thick] (2.55,0.45) rectangle (3.25,0.85);
\draw[fill=popgreen!60!black, thick] (2.55,-0.85) rectangle (3.25,-0.45);
\node[popdark, font=\small] at (4.15,1.15) {4. Balais};
\draw[->, popdark] (3.9,1.0) -- (3.1,0.75);
% Bornes
\draw[thick, popdark] (2.75,0.85) -- (2.75,2.6);
\draw[thick, popdark] (3.05,-0.85) -- (3.05,-2.6);
\node[popdark, font=\small] at (2.75,2.85) {$+$};
\node[popdark, font=\small] at (3.05,-2.85) {$-$};
\end{tikzpicture}
\legende{Coupe schématique d'un moteur à courant continu. 1 : stator (aimant inducteur, pôles N et S) ; 2 : rotor (bobines de cuivre sur noyau de fer doux) ; 3 : collecteur à lamelles ; 4 : balais (charbons) reliés aux bornes $+$ et $-$ du générateur ; 5 : arbre de sortie ; 6 : carter protecteur.}
\end{popfigure}

\begin{popfigure}
\begin{tikzpicture}[scale=1.0, every node/.style={font=\small}]
% ===== Position 1 =====
\begin{scope}[xshift=-3.6cm]
\node[popdark, font=\small\bfseries] at (0,2.3) {Position 1};
% bobine (cadre)
\draw[thick, popblue, fill=popblue!15] (-1.1,-0.7) rectangle (1.1,0.7);
\node[popblue, font=\small] at (0,0) {bobine};
% collecteur deux lamelles
\draw[fill=popgold!80, thick] (1.5,-0.45) arc[start angle=-90, end angle=90, radius=0.45] -- cycle;
\draw[fill=poporange!80, thick] (2.1,0.45) arc[start angle=90, end angle=270, radius=0.45] -- cycle;
\draw[thick, popblue] (1.1,0.3) -- (1.55,0.35);
\draw[thick, popblue] (1.1,-0.3) -- (1.55,-0.35);
% balais
\draw[fill=popgreen!60!black] (1.35,0.55) rectangle (1.75,0.95);
\draw[fill=popgreen!60!black] (1.35,-0.95) rectangle (1.75,-0.55);
\node[popdark] at (1.1,1.15) {$+$};
\node[popdark] at (1.1,-1.15) {$-$};
\draw[thick, popdark] (1.55,0.95) -- (1.55,1.35);
\draw[thick, popdark] (1.55,-0.95) -- (1.55,-1.35);
% courant
\draw[->, very thick, poppink] (1.55,0.45) -- (1.55,0.05);
\node[poppink, font=\small] at (2.05,0.25) {$I$};
\node[popdark, font=\small] at (0,-1.7) {Le courant entre par la lamelle A};
\end{scope}
% ===== Position 2 : demi-tour =====
\begin{scope}[xshift=3.6cm]
\node[popdark, font=\small\bfseries] at (0,2.3) {Après un demi-tour};
\draw[thick, popblue, fill=popblue!15] (-1.1,-0.7) rectangle (1.1,0.7);
\node[popblue, font=\small] at (0,0) {bobine};
% lamelles inversées
\draw[fill=poporange!80, thick] (1.5,-0.45) arc[start angle=-90, end angle=90, radius=0.45] -- cycle;
\draw[fill=popgold!80, thick] (2.1,0.45) arc[start angle=90, end angle=270, radius=0.45] -- cycle;
\draw[thick, popblue] (1.1,0.3) -- (1.55,0.35);
\draw[thick, popblue] (1.1,-0.3) -- (1.55,-0.35);
\draw[fill=popgreen!60!black] (1.35,0.55) rectangle (1.75,0.95);
\draw[fill=popgreen!60!black] (1.35,-0.95) rectangle (1.75,-0.55);
\node[popdark] at (1.1,1.15) {$+$};
\node[popdark] at (1.1,-1.15) {$-$};
\draw[thick, popdark] (1.55,0.95) -- (1.55,1.35);
\draw[thick, popdark] (1.55,-0.95) -- (1.55,-1.35);
\draw[->, very thick, poppink] (1.55,-0.05) -- (1.55,-0.45);
\node[poppink, font=\small] at (2.15,-0.25) {$I$};
\node[popdark, font=\small] at (0,-1.7) {Le courant a changé de sens};
\end{scope}
\end{tikzpicture}
\legende{Rôle inverseur du collecteur à deux lamelles. À gauche : la lamelle A (dorée) touche le balais $+$, le courant descend dans la bobine. À droite : après un demi-tour, la lamelle B (orange) touche le balais $+$ : le courant circule en sens inverse dans la bobine. Les balais restent fixes, les lamelles tournent avec le rotor.}
\end{popfigure}

\courssub{Bilan : rôle de chaque élément}

\begin{aretenir}[Les six constituants du moteur à courant continu]
\begin{center}
\begin{tabularx}{\linewidth}{|l|c|X|}
\hline
\rowcolor{popblueL} \textbf{Élément} & \textbf{Nature} & \textbf{Fonction} \\
\hline
Stator (aimant inducteur) & Fixe & Crée le champ magnétique dans lequel tourne le rotor. \\
\hline
Rotor (bobines + noyau) & Mobile & Reçoit le courant et tourne sous l'action des forces magnétiques. \\
\hline
Collecteur & Mobile & Amène le courant dans la bobine et inverse son sens à chaque demi-tour. \\
\hline
Balais (charbons) & Fixes & Assurent le contact électrique (par frottement) entre les bornes et le collecteur. \\
\hline
Arbre de sortie & Mobile & Transmet le mouvement de rotation à l'appareil entraîné (hélice, meule, tambour). \\
\hline
Carter & Fixe & Protège les pièces internes et l'utilisateur ; supporte le stator et les paliers. \\
\hline
\end{tabularx}
\end{center}
\end{aretenir}

\begin{methode}[Identifier les constituants d'un moteur sur un schéma ou un objet réel]
Face à un schéma de moteur (ou à un vrai moteur démonté), procède toujours dans le même ordre :
\begin{enumerate}
\item \textbf{Repère d'abord l'aimant} (ou l'électroaimant) : pièce fixe, souvent dessinée contre la paroi avec les lettres N et S. Tu as trouvé le \textbf{stator}.
\item \textbf{Cherche ensuite la bobine} de fil de cuivre enroulée sur un noyau, au centre : tu as trouvé le \textbf{rotor}.
\item \textbf{Suis le fil de la bobine} jusqu'à un petit cylindre fendu en lamelles : c'est le \textbf{collecteur}. Les deux petits blocs qui le touchent sont les \textbf{balais}.
\item \textbf{Repère l'axe} qui traverse le rotor et dépasse du moteur : c'est l'\textbf{arbre de sortie}.
\item \textbf{Termine par l'enveloppe extérieure} : c'est le \textbf{carter}.
\item Vérifie ta réponse en classant chaque pièce : \emph{fixe} (stator, balais, carter) ou \emph{mobile} (rotor, collecteur, arbre).
\end{enumerate}
\end{methode}

\begin{exoresolu}[Identifier les éléments d'un moteur de ventilateur]
Énoncé. En démontant le moteur d'un vieux ventilateur, un élève de 3ème observe : (a) un aimant en forme d'anneau fixé contre la paroi ; (b) trois bobines de fil de cuivre verni montées sur un axe ; (c) un petit cylindre de cuivre fendu en trois secteurs solidaire de l'axe ; (d) deux petits blocs noirs qui frottent sur ce cylindre ; (e) l'axe qui porte l'hélice. Nomme chaque élément, précise s'il est fixe ou mobile, et donne sa fonction.
\tcblower
Solution détaillée.
\begin{itemize}
\item (a) L'aimant fixé contre la paroi est le \textbf{stator} (inducteur) : il est \textbf{fixe} et crée le champ magnétique.
\item (b) Les trois bobines forment le \textbf{rotor} : elles sont \textbf{mobiles} ; parcourues par le courant, elles subissent les forces magnétiques qui font tourner le moteur.
\item (c) Le cylindre fendu est le \textbf{collecteur} : il est \textbf{mobile} ; il alimente les bobines et inverse le courant à chaque demi-tour.
\item (d) Les blocs noirs sont les \textbf{balais} (charbons) : ils sont \textbf{fixes} et assurent le contact électrique avec le collecteur par frottement.
\item (e) L'axe est l'\textbf{arbre de sortie} : il est \textbf{mobile} et transmet la rotation à l'hélice du ventilateur.
\end{itemize}
Vérification : trois éléments fixes (aimant, balais, carter non cité) et trois éléments mobiles (bobines, collecteur, arbre) : la classification est cohérente.
\end{exoresolu}

\begin{exercice}[À toi de jouer]
Un schéma de moteur à courant continu porte les légendes suivantes, dans le désordre : arbre, balais, carter, collecteur, rotor, stator.
\begin{enumerate}
\item Classe ces éléments en deux colonnes : partie fixe / partie mobile.
\item Explique pourquoi on ne peut pas relier la bobine directement aux bornes du générateur par des fils soudés.
\item Que se passerait-il si le collecteur n'inversait pas le courant dans la bobine ?
\end{enumerate}
\end{exercice}

\begin{corrige}[Exercice : À toi de jouer]
\begin{enumerate}
\item Partie fixe : stator (aimant), balais, carter. Partie mobile : rotor (bobines), collecteur, arbre.
\item Parce que la bobine tourne en permanence : des fils soudés s'enrouleraient autour de l'arbre et se rompraient après quelques tours. Le contact glissant balais-collecteur résout ce problème.
\item Sans inversion du courant à chaque demi-tour, les forces magnétiques changeraient de sens après un demi-tour et ramèneraient la bobine en arrière : le rotor oscillerait puis s'immobiliserait en position d'équilibre au lieu de tourner continûment.
\end{enumerate}
\end{corrige}

Nous connaissons maintenant toutes les pièces du moteur et leur rôle. Mais une question fondamentale demeure : \emph{pourquoi} une bobine parcourue par un courant se met-elle à tourner entre les pôles d'un aimant ? C'est l'objet de l'étude théorique de la section suivante.

\coursec{Étude théorique : action d'un aimant sur une bobine parcourue par un courant}

Dans la section précédente, nous avons démonté (en imagination !) un moteur électrique à courant continu : nous y avons trouvé un aimant fixe (le \cle{stator}), une bobine qui tourne (le \cle{rotor}), un collecteur et des balais. Mais une question reste entière : \emph{pourquoi} la bobine se met-elle à tourner quand on branche le moteur ? Quelle force mystérieuse la fait pivoter ? C'est à cette question que répond l'étude théorique. Nous allons procéder comme de vrais scientifiques : réaliser des expériences simples, observer, puis conclure. À la clé : le principe physique qui fait tourner les ventilateurs pendant les délestages, les perceuses du menuisier et le mixer de maman.

\courssub{Situation-problème : le fil qui bouge tout seul}

\begin{exemplebox}[Une observation intrigante]
Au laboratoire du collège, le professeur suspend un fil de cuivre souple entre les deux branches d'un aimant en U (l'aimant « en fer à cheval »). Le fil est immobile. Puis il relie les deux extrémités du fil à une pile. Aussitôt, le fil \textbf{saute} et se déplace latéralement ! Il n'y a pourtant aucun contact, aucun souffle, aucune secousse. D'où vient cette force invisible ?
\end{exemplebox}

Pour répondre, menons l'enquête par une démarche d'investigation complète : hypothèse, expériences, observations, interprétation, conclusion.

\textbf{Hypothèse de travail.} Le déplacement du fil pourrait être dû à la présence simultanée de deux ingrédients : le \cle{courant électrique} qui circule dans le fil et l'\cle{aimant} qui crée un champ magnétique autour du fil. Vérifions cela en faisant varier un seul paramètre à la fois — c'est la règle d'or de toute expérience scientifique.

\courssub{Les quatre expériences fondamentales}

\begin{experience}[Expérience 1 : la spire parcourue par un courant se déplace]
\textbf{Matériel :} un aimant en U, une spire de fil de cuivre isolé (ou un cadre rectangulaire souple), un générateur de courant continu (pile ou alimentation \SI{6}{V}), un interrupteur, des fils de connexion, un support avec fil de suspension.

\textbf{Protocole :}
\begin{enumerate}
\item Suspendre la spire entre les deux pôles de l'aimant en U, de façon qu'elle puisse pivoter librement.
\item Relier la spire au générateur par l'interrupteur ouvert.
\item Fermer l'interrupteur et observer.
\end{enumerate}

\textbf{Observation :} dès que le circuit est fermé, la spire \textbf{se déplace} (elle pivote ou s'écarte de sa position d'équilibre). Elle subit donc une \textbf{force}.

\textbf{Interprétation :} le courant seul (sans aimant) ne produit rien de visible ; l'aimant seul (sans courant) non plus. C'est la \emph{réunion} des deux qui fait naître une force.
\end{experience}

\begin{experience}[Expérience 2 : on inverse le sens du courant]
\textbf{Protocole :} reprendre le montage de l'expérience 1, mais permuter les deux bornes du générateur (le pôle $+$ à la place du pôle $-$ et inversement).

\textbf{Observation :} la spire se déplace \textbf{dans le sens opposé} à celui observé dans l'expérience 1.

\textbf{Interprétation :} le sens de la force électromagnétique \textbf{dépend du sens du courant}. Inverser le courant inverse le mouvement.
\end{experience}

\begin{experience}[Expérience 3 : on inverse les pôles de l'aimant]
\textbf{Protocole :} remettre le générateur comme dans l'expérience 1, mais retourner l'aimant en U : le pôle Nord prend la place du pôle Sud et inversement.

\textbf{Observation :} la spire se déplace à nouveau \textbf{dans le sens opposé} à celui de l'expérience 1.

\textbf{Interprétation :} le sens de la force dépend aussi du \textbf{sens du champ magnétique} (c'est-à-dire de la position des pôles). Inverser les pôles inverse le mouvement.
\end{experience}

\begin{experience}[Expérience 4 : on coupe le courant]
\textbf{Protocole :} ouvrir l'interrupteur pendant que la spire est entre les pôles de l'aimant.

\textbf{Observation :} la spire revient à sa position initiale et reste immobile. \textbf{Aucun déplacement} ne se produit, même si l'aimant est très puissant.

\textbf{Interprétation :} la force électromagnétique \textbf{n'existe que si le courant circule}. Pas de courant, pas de force, pas de mouvement.
\end{experience}

\begin{popfigure}
\begin{center}
\begin{tikzpicture}[scale=1.0,>=Stealth]
% Aimant en U
\fill[popblue!80] (-2.6,-0.2) rectangle (-1.6,2.2);
\fill[poporange!85] (1.6,-0.2) rectangle (2.6,2.2);
\fill[popmuted!40] (-2.6,-1.4) rectangle (2.6,-0.2);
\fill[popblue!80] (-2.6,-0.2) rectangle (2.6,0.0);
\node[white,font=\bfseries] at (-2.1,1.0) {N};
\node[white,font=\bfseries] at (2.1,1.0) {S};
% Suspension
\draw[thick,popdark] (-0.5,3.4) -- (-0.5,2.6);
\draw[thick,popdark] (0.5,3.4) -- (0.5,2.6);
\draw[thick,popdark] (-1.2,3.4) -- (1.2,3.4);
\fill[poppurple] (-1.2,3.4) circle (0.06);
\fill[poppurple] (1.2,3.4) circle (0.06);
\node[font=\small,popdark] at (0,3.75) {fils de suspension};
% Spire (cadre)
\draw[very thick,poppurple] (-0.5,2.6) -- (-0.5,0.4) -- (0.5,0.4) -- (0.5,2.6);
\node[font=\small,poppurple] at (0,-0.05) {spire (fil de cuivre)};
% Courant dans le brin gauche
\draw[->,very thick,popgreen] (-0.5,1.9) -- (-0.5,1.0);
\node[font=\small,popgreen,left] at (-0.55,1.45) {$I$};
% Flèche de déplacement
\draw[->,ultra thick,poporange] (0.7,1.5) -- (1.5,1.5);
\node[font=\small,poporange] at (1.1,1.85) {déplacement};
% Générateur
\draw[thick,popdark] (-0.5,2.6) -- (-0.5,4.4) -- (-4.0,4.4) -- (-4.0,-2.0);
\draw[thick,popdark] (0.5,2.6) -- (0.5,4.6) -- (4.4,4.6) -- (4.4,-2.0);
% Interrupteur
\draw[thick,popdark] (-4.0,-2.0) -- (-3.2,-2.0);
\draw[thick,popdark] (-3.2,-2.0) -- (-2.6,-1.6);
\fill[popdark] (-3.2,-2.0) circle (0.05);
\fill[popdark] (-2.5,-2.0) circle (0.05);
\draw[thick,popdark] (-2.5,-2.0) -- (-1.4,-2.0);
\node[font=\small,popdark] at (-2.85,-2.45) {interrupteur};
% Pile
\draw[very thick,popdark] (-1.0,-2.15) -- (-1.0,-1.85);
\draw[very thick,popdark] (-0.6,-2.35) -- (-0.6,-1.65);
\draw[thick,popdark] (-1.4,-2.0) -- (-1.0,-2.0);
\draw[thick,popdark] (-0.6,-2.0) -- (4.4,-2.0);
\node[font=\small,popdark] at (-0.8,-2.7) {générateur};
\node[font=\small,popdark] at (-1.15,-1.55) {$+$};
\end{tikzpicture}
\end{center}
\legende{Schéma de l'expérience : une spire de fil de cuivre, suspendue entre les pôles d'un aimant en U, est reliée à un générateur. Quand le courant $I$ circule, la spire se déplace dans le sens de la flèche orange : elle subit une force électromagnétique.}
\end{popfigure}

\courssub{Tableau récapitulatif des résultats expérimentaux}

Rassemblons nos quatre observations dans un tableau. C'est une habitude à prendre pour le BEPC : un tableau clair vaut souvent de longues phrases.

\begin{popfigure}
\begin{center}
\begin{tabular}{|c|c|c|c|}
\hline
\rowcolor{popblueL} \textbf{Expérience} & \textbf{Sens du courant} & \textbf{Position des pôles} & \textbf{Déplacement} \\
\hline
1 & sens initial & N à gauche, S à droite & vers la droite \\
\hline
2 & \textbf{inversé} & N à gauche, S à droite & vers la \textbf{gauche} \\
\hline
3 & sens initial & \textbf{S à gauche, N à droite} & vers la \textbf{gauche} \\
\hline
4 & \textbf{courant coupé} & N à gauche, S à droite & \textbf{aucun} \\
\hline
\end{tabular}
\end{center}
\legende{Bilan des quatre expériences : inverser le courant \emph{ou} les pôles inverse le sens du déplacement ; sans courant, il n'y a aucun déplacement.}
\end{popfigure}

\courssub{Conclusion : la force électromagnétique}

\begin{propriete}[Force électromagnétique (établie expérimentalement)]
Un conducteur \textbf{parcouru par un courant électrique} et placé dans un \textbf{champ magnétique} (par exemple entre les pôles d'un aimant) subit une \cle{force électromagnétique} qui peut le mettre en mouvement.
\begin{itemize}
\item Le \textbf{sens} de cette force dépend du sens du courant et du sens du champ magnétique (position des pôles).
\item Cette force \textbf{n'existe que si le courant circule} : dès qu'on coupe le courant, elle disparaît.
\end{itemize}
Cette propriété est établie par l'expérience ; aucune démonstration mathématique n'est exigible au BEPC, mais tu dois savoir la formuler et l'illustrer par le schéma de l'expérience.
\end{propriete}

\begin{attention}[Erreur fréquente : « l'aimant attire le fil »]
Beaucoup d'élèves écrivent : « le fil bouge parce que l'aimant l'attire ». C'est \textbf{faux} ! L'expérience 4 le prouve : quand le courant est coupé, l'aimant est toujours là, et pourtant le fil de cuivre ne bouge pas d'un millimètre. D'ailleurs, un aimant n'attire pas le cuivre (essaie avec une pièce de monnaie jaune ou un fil électrique dénudé !). La force n'apparaît que lorsque \textbf{deux conditions sont réunies} : le courant circule \emph{et} le conducteur est dans un champ magnétique. Au BEPC, écris toujours les deux conditions.
\end{attention}

\courssub{Prévoir le sens du déplacement}

\begin{methode}[Prévoir le sens du mouvement]
Pour prévoir le sens du déplacement d'un conducteur placé entre les pôles d'un aimant :
\begin{enumerate}
\item Repérer le \textbf{sens du courant} dans le conducteur (du $+$ du générateur vers le $-$, à l'extérieur).
\item Repérer la \textbf{position des pôles} de l'aimant (où est le Nord, où est le Sud).
\item Appliquer les règles issues des expériences :
\begin{itemize}
\item inverser le courant \textbf{seul} $\Rightarrow$ le mouvement s'inverse ;
\item inverser les pôles \textbf{seuls} $\Rightarrow$ le mouvement s'inverse ;
\item inverser le courant \textbf{et} les pôles en même temps $\Rightarrow$ le mouvement \textbf{ne change pas} (deux inversions s'annulent).
\end{itemize}
\end{enumerate}
\end{methode}

\begin{exoresolu}[Prévision du sens de déplacement]
\textbf{Énoncé.} Une spire suspendue entre les pôles d'un aimant en U se déplace vers la droite quand le courant circule dans un certain sens (expérience 1). Que se passe-t-il si l'on inverse à la fois les bornes du générateur \emph{et} les pôles de l'aimant ? Justifier.

\tcblower
\textbf{Solution détaillée.}
\begin{itemize}
\item Inverser les bornes du générateur inverse le sens du courant : d'après l'expérience 2, cela inverse le sens du déplacement (la spire irait vers la gauche).
\item Inverser les pôles de l'aimant : d'après l'expérience 3, cela inverse aussi le sens du déplacement.
\item Les deux inversions se compensent : $-1 \times (-1) = +1$. Le déplacement redevient identique à celui de l'expérience 1.
\end{itemize}
\textbf{Conclusion :} la spire se déplace toujours \textbf{vers la droite}. Inverser le courant et les pôles en même temps laisse le mouvement inchangé.
\end{exoresolu}

\courssub{De la force au mouvement de rotation : l'idée de Faraday}

Nos expériences montrent qu'une spire parcourue par un courant et placée dans un champ magnétique \emph{pivote}. Mais une spire suspendue ne tourne que d'un petit angle, puis s'arrête. Comment obtenir une \textbf{rotation continue}, capable d'entraîner un ventilateur ou une perceuse ?

C'est exactement le problème que s'est posé le physicien anglais Michael Faraday. En 1821, il réalisa un dispositif étonnant : un fil conducteur vertical, plongé dans un bain de mercure (métal liquide qui conduit le courant), avec un aimant au fond. Quand le courant traversait le fil, celui-ci se mettait à \textbf{tourner sans fin autour de l'aimant}. Pour la première fois dans l'histoire, l'électricité produisait un mouvement de rotation continu : le moteur électrique était né.

\begin{savaistu}[1821 : la naissance du moteur électrique]
Michael Faraday (1791-1867), fils de forgeron devenu l'un des plus grands expérimentateurs de tous les temps, construisit en 1821 le premier dispositif transformant l'énergie électrique en mouvement de rotation continu. On l'appelle parfois le « moteur homopolaire ». Dix ans plus tard, en 1831, il découvrit le phénomène inverse — un aimant en mouvement crée un courant dans une bobine — qui est à la base des générateurs et des alternateurs de nos groupes électrogènes et des barrages hydroélectriques comme ceux de la Sanaga. Aujourd'hui, des milliards de moteurs électriques fonctionnent sur le principe que tu viens d'étudier : du moteur de la moto-taxi (démarreur électrique) au ventilateur du salon.
\end{savaistu}

\courssub{Bilan de l'étude théorique}

\begin{aretenir}[La chaîne énergétique du moteur électrique]
Le phénomène à la base du moteur électrique se résume en une chaîne simple :
\[
\text{courant électrique} + \text{champ magnétique} \;\longrightarrow\; \text{force électromagnétique} \;\longrightarrow\; \text{mouvement}
\]
\begin{itemize}
\item Un conducteur parcouru par un courant et placé dans un champ magnétique subit une \textbf{force électromagnétique}.
\item Cette force ne s'exerce \textbf{que si le courant circule}.
\item Son sens dépend du \textbf{sens du courant} et de la \textbf{position des pôles}.
\item Dans un moteur, la bobine du rotor est parcourue par un courant et placée dans le champ magnétique du stator : les forces électromagnétiques la font \textbf{tourner}.
\end{itemize}
Le moteur électrique convertit ainsi l'\textbf{énergie électrique} en \textbf{énergie mécanique} (mouvement).
\end{aretenir}

Tu comprends maintenant \emph{pourquoi} la bobine du rotor tourne : chaque portion de fil parcourue par le courant, plongée dans le champ de l'aimant, est poussée par une force électromagnétique. Reste un détail technique : pour que la rotation ne s'arrête pas après un demi-tour, il faut inverser le courant dans la bobine au bon moment. C'est précisément le rôle du collecteur et des balais que nous avons découverts dans l'étude technique — et que nous relierons au principe de fonctionnement dans la prochaine section.

\begin{exercice}[Pour t'entraîner]
\textbf{Exercice 1.} Un fil de cuivre rectiligne est placé entre les pôles d'un aimant en U. On observe qu'il ne bouge pas. Donne deux explications possibles et l'expérience qui permettrait de trancher.

\textbf{Exercice 2.} Dans l'expérience de la spire suspendue, le déplacement s'effectue vers la gauche. On souhaite obtenir un déplacement vers la droite. Cite deux modifications indépendantes du montage qui le permettent, puis une modification double qui laisserait le déplacement vers la gauche.
\end{exercice}

\begin{corrige}[Exercice 1]
Deux explications possibles : (a) le circuit est ouvert ou le générateur est débranché, donc aucun courant ne circule dans le fil ; (b) le fil n'est pas correctement placé dans le champ magnétique (par exemple trop loin des pôles). Pour trancher, on ferme le circuit et on vérifie avec une lampe témoin (ou un ampèremètre) que le courant circule : si le courant passe et que le fil ne bouge toujours pas, c'est la position dans le champ qui est en cause. On retient surtout : sans courant, aucune force électromagnétique n'existe.
\end{corrige}

\begin{corrige}[Exercice 2]
Pour inverser le déplacement (de gauche vers droite), on peut : inverser les bornes du générateur (sens du courant inversé) \emph{ou} retourner l'aimant (pôles inversés). Une modification double qui laisse le déplacement inchangé (toujours vers la gauche) : inverser \textbf{à la fois} le sens du courant et les pôles de l'aimant, car les deux inversions se compensent.
\end{corrige}

\coursec{Le moteur électrique dans la vie courante}

Partout autour de nous, des dispositifs transforment l'énergie électrique en mouvement de rotation : le ventilateur qui brasse l'air dans la chambre, la pompe à eau qui remonte l'eau du forage, le lave-linge qui essore le linge, le démarreur de la moto-taxi ou de la voiture, le groupe électrogène qui démarre seul, la meuleuse du soudeur, l'ascenseur d'un immeuble. Tous ces appareils contiennent un \cle{moteur électrique}.

\begin{definition}[Moteur électrique]
Un \cle{moteur électrique} est un \cle{convertisseur d'énergie} qui transforme l'\textbf{énergie électrique} (fourni par une pile, une batterie ou le réseau ENEO) en \textbf{énergie mécanique} de rotation (arbre tournant) pour entraîner une machine (hélice, pompe, tambour, meule, roue).
\end{definition}

\begin{savaistu}[Le démarreur de la moto-taxi]
Quand tu tournes la clé ou appuies sur le bouton de démarrage d'une moto-taxi (type Yamaha XTZ 125 ou Bajaj Boxer), un petit moteur électrique puissant — le \emph{démarreur} — s'enclenche. Il puise un fort courant dans la batterie (souvent \SI{12}{V}, \SI{50}{A} au pic) pour faire tourner le moteur thermique le temps qu'il démarre tout seul. Sans ce moteur électrique, il faudrait une pédale de kick robuste et beaucoup de force physique !
\end{savaistu}

\coursec{Étude technique : constitution d'un moteur électrique à courant continu}

Avant de comprendre \emph{comment} il tourne, ouvrons un moteur à courant continu (celui des jouets, des perceuses portatives, des essuie-glaces, des vitres électriques de voiture). Il comporte deux ensembles principaux :

\begin{itemize}
\item Le \cle{stator} (partie \textbf{fixe}) : il crée le champ magnétique. C'est soit un \textbf{aimant permanent} (petits moteurs), soit un \textbf{électroaimant} (gros moteurs : bobines sur un noyau en fer doux alimentées par le même courant que le rotor).
\item Le \cle{rotor} (partie \textbf{mobile}, aussi appelé \cle{armature}) : il tourne. Il est constitué d'un \textbf{noyau en fer doux} lamellé (pour limiter les courants de Foucault), sur lequel sont enroulées une ou plusieurs \textbf{bobines} (enroulements de fil de cuivre émaillé). Les extrémités des bobines sont soudées sur un \cle{collecteur} (ou \cle{collecteur à bagues} pour le courant continu).
\item Les \cle{balais} (ou \cle{charbons}) : pièces fixes en \textbf{graphite} (bon conducteur, auto-lubrifiant) qui pressent sur le collecteur tournant et assurent le contact électrique entre le circuit extérieur (pile, ENEO) et les bobines du rotor.
\item La \textbf{carcasse} (châssis), les \textbf{roulements} (pour que l'axe tourne sans frottement excessif) et le \textbf{ventilateur} de refroidissement (souvent fixé sur l'axe côté opposé à la sortie utile).
\end{itemize}

\begin{popfigure}
\begin{center}
\begin{tikzpicture}[scale=0.85, every node/.style={font=\small}]
% Carcasse
\draw[line width=2pt, popdark] (-3.5,-2.5) rectangle (3.5,2.5);
\draw[line width=1pt, popmuted] (-3.3,-2.3) rectangle (3.3,2.3);
% Stator : aimants permanents (simplifié)
\fill[popred!70] (-3.3,-2) rectangle (-2.7,2);
\fill[popblue!70] (2.7,-2) rectangle (3.3,2);
\node[white, font=\bfseries, rotate=90] at (-3.0,0) {N};
\node[white, font=\bfseries, rotate=-90] at (3.0,0) {S};
% Rotor : noyau + enroulements (schématisés)
\draw[fill=popgray!30] (-1.2,-1.2) rectangle (1.2,1.2); % noyau
\foreach \y in {-0.8, 0, 0.8} {
  \draw[line width=1.5pt, poporange] (-1.2,\y) -- (1.2,\y); % spires
}
% Axe
\draw[line width=3pt, popdark] (-1.2,0) -- (-2.5,0);
\draw[line width=3pt, popdark] (1.2,0) -- (2.5,0);
% Collecteur (côté droit)
\draw[fill=popgold!80!black] (1.3,-0.4) rectangle (1.7,0.4);
% Balais
\draw[fill=popdark] (2.0,-0.5) rectangle (2.5,-0.3);
\draw[fill=popdark] (2.0,0.3) rectangle (2.5,0.5);
% Fils d'alimentation
\draw[line width=1.5pt, popblue] (2.5,-0.4) -- (3.3,-0.4);
\draw[line width=1.5pt, popred] (2.5,0.4) -- (3.3,0.4);
% Légendes par flèches
\draw[-{Stealth}, popblue, line width=1pt] (-3.8,1.5) -- (-3.0,1.0) node[left, align=right] {Stator\\(aimants N/S)};
\draw[-{Stealth}, poporange, line width=1pt] (-3.8,0) -- (-1.5,0.2) node[left] {Rotor (noyau + bobines)};
\draw[-{Stealth}, popgold!80!black, line width=1pt] (3.8,0) -- (1.8,0) node[right] {Collecteur};
\draw[-{Stealth}, popdark, line width=1pt] (3.8,-1.5) -- (2.5,-0.5) node[right] {Balais (charbons)};
\draw[-{Stealth}, popmuted, line width=1pt] (0,-3.0) -- (0,-2.5) node[below] {Axe de rotation};
\end{tikzpicture}
\end{center}
\legende{Coupe schématique d'un moteur à courant continu à aimants permanents (vue de face). Le stator crée le champ B (N vers S). Le rotor porte les bobines (orange) connectées au collecteur (or). Les balais (noirs) frottent sur le collecteur.}
\end{popfigure}

\begin{attention}[Ne pas confondre stator et rotor]
\emph{Stator} = \textbf{Stat}ionnaire (Fixe). \emph{Rotor} = \textbf{Rot}ation (Mobile). Astuce mnémotechnique : le \textbf{St}ator \textbf{St}agne, le \textbf{Ro}tor \textbf{Ro}ule.
\end{attention}

\coursec{Étude théorique : action d'un aimant sur une bobine parcourue par un courant}

\courssub{Rappel : la force électromagnétique (force de Laplace)}

En classe de 4ème, nous avons mis en évidence qu'un \textbf{conducteur rectiligne} parcouru par un courant \textbf{I} et placé dans un \textbf{champ magnétique B} subit une force \textbf{F} (force de Laplace) si le courant n'est pas parallèle au champ.

\begin{propriete}[Force électromagnétique sur un conducteur — Savoir admis]
\begin{itemize}
\item \textbf{Condition :} conducteur parcouru par un courant $I$, immergé dans un champ magnétique $\vec{B}$, non parallèle à $\vec{B}$.
\item \textbf{Direction :} perpendiculaire au plan formé par le conducteur et $\vec{B}$.
\item \textbf{Sens :} donné par la \textbf{règle des trois doigts de la main droite} (pouce = courant $I$, index = champ $\vec{B}$, majeur = force $\vec{F}$).
\item \textbf{Valeur :} $F = B \cdot I \cdot L \cdot \sin\alpha$ (avec $L$ longueur du conducteur dans le champ, $\alpha$ angle entre $I$ et $B$). Au BEPC, on retient que la force est \textbf{maximale} quand le fil est \textbf{perpendiculaire} au champ ($\alpha = 90^\circ$) et \textbf{nulle} quand il est \textbf{parallèle}.
\end{itemize}
\end{propriete}

\courssub{Application à une bobine rectangulaire : naissance du couple}

Plaçons une \textbf{bobine rectangulaire} (une spire ou plusieurs spires en série) dans un champ magnétique uniforme (entre les pôles N et S du stator). Le courant entre par un balai, parcourt la bobine, et sort par l'autre balai.

\begin{itemize}
\item Les deux \textbf{côtés actifs} (perpendiculaires à $\vec{B}$) sont parcourus par des courants de \textbf{sens opposés} (le courant descend d'un côté, remonte de l'autre).
\item D'après la règle des trois doigts, les forces $\vec{F}_1$ et $\vec{F}_2$ sur ces deux côtés ont \textbf{des sens opposés}.
\item Ces deux forces sont \textbf{parallèles}, de \textbf{même valeur}, mais \textbf{non confondues} (elles ne s'exercent pas sur la même droite).
\end{itemize}

\begin{definition}[Couple de rotation]
Un \cle{couple de rotation} est un système de deux forces \textbf{parallèles}, de \textbf{sens opposés}, de \textbf{même grandeur}, dont les \textbf{lignes d'action sont distinctes}. Il provoque une \textbf{rotation} autour d'un axe sans translation. Son moment (en N$\cdot$m) vaut $M = F \times d$ ($d$ = distance entre les lignes d'action).
\end{definition}

\begin{exemplebox}[Image concrète : la molette de la bonbonne de gaz]
Pour ouvrir le robinet d'une bonbonne de gaz butane, tu places deux doigts de part et d'autre de la molette. Tu pousses l'un vers l'avant, l'autre vers l'arrière. Tes deux forces créent un couple qui fait tourner la molette. Dans le moteur, les forces sont \emph{électromagnétiques} (sans contact mécanique) mais l'effet est identique : la bobine pivote.
\end{exemplebox}

\courssub{La position d'équilibre (le point mort)}

La bobine tourne sous l'effet du couple. Mais quand elle arrive en position \textbf{perpendiculaire au champ magnétique} (plan de la bobine parallèle à $\vec{B}$), les deux côtés actifs deviennent \textbf{parallèles à $\vec{B}$}.
\begin{itemize}
\item L'angle $\alpha$ entre $I$ et $B$ devient $0^\circ$ (ou $180^\circ$).
\item $\sin\alpha = 0 \Rightarrow F = 0$.
\item Les forces s'annulent : le couple est \textbf{nul}.
\end{itemize}
C'est la \cle{position d'équilibre} (ou \cle{point mort}). Sans dispositif additionnel, la bobine s'arrête là (ou oscille autour si elle a de l'inertie).

\begin{attention}[Piège BEPC : « La bobine s'arrête car les forces s'annulent »]
Oui, mais \textbf{précise} : « Les forces deviennent \textbf{nulles} car les conducteurs sont \textbf{parallèles au champ magnétique} ($\vec{I} \parallel \vec{B}$), donc $\sin\alpha = 0$. » Ne dis pas « les forces s'annulent » (ce qui sous-entend qu'elles existent encore mais s'additionnent à zéro), dis « les forces \textbf{deviennent nulles} ».
\end{attention}

\coursec{Principe de fonctionnement du moteur électrique : le rôle du collecteur}

\courssub{L'ingéniosité du collecteur (commutateur)}

Pour que la rotation continue \textbf{après} le passage du point mort, il faut que les forces \textbf{changent de sens} au bon moment, c'est-à-dire \textbf{juste après} la position d'équilibre. Or, le sens de la force dépend du sens du courant. Il faut donc \textbf{inverser le courant dans la bobine à chaque demi-tour}.

C'est le rôle du \cle{collecteur} (aussi appelé \cle{commutateur}) : c'est un \textbf{cylindre fendu en deux demi-anneaux isolés l'un de l'autre} (souvent en cuivre), solidaire de l'axe du rotor. Chaque demi-anneau est relié à une extrémité de la bobine. Les \textbf{balais (charbons)} fixes frottent sur ces demi-anneaux.

\begin{propriete}[Fonction du collecteur — Savoir admis, explication exigible]
Dans un moteur à courant continu, le \cle{collecteur} associé aux \cle{balais} \textbf{inverse automatiquement le sens du courant dans la bobine du rotor à chaque demi-tour}.
\begin{itemize}
\item Avant le point mort : le balai $+$ est sur le demi-anneau $A$, le balai $-$ sur $B$. Courant dans un sens.
\item Au point mort : les balais frôlent l'isolant entre les deux demi-anneaux (coupure brève).
\item Après le point mort : le balai $+$ passe sur le demi-anneau $B$, le balai $-$ sur $A$. Le courant dans la bobine \textbf{s'inverse}.
\end{itemize}
Les forces électromagnétiques s'inversent donc aussi, et poussent la bobine \textbf{toujours dans le même sens de rotation}.
\end{propriete}

\begin{popfigure}
\begin{center}
\begin{tikzpicture}[scale=0.95, every node/.style={font=\small}]
% --- ÉTAPE 1 : Position horizontale, courant initial ---
\begin{scope}[shift={(0,0)}]
\node[font=\small\bfseries, popblue] at (0,2.8) {Étape 1 : Bobine horizontale, courant initial};
% Aimant stator
\fill[popred!70] (-2.5,-1.5) rectangle (-1.7,1.5);
\fill[popblue!70] (1.7,-1.5) rectangle (2.5,1.5);
\node[white, font=\bfseries] at (-2.1,0) {N};
\node[white, font=\bfseries] at (2.1,0) {S};
% Champ B (flèches horizontales)
\foreach \y in {-1.0, 0, 1.0} {
  \draw[-{Stealth}, popmuted, line width=0.8pt] (-1.5,\y) -- (1.5,\y);
}
\node[popmuted] at (0,1.6) {$\vec{B}$};
% Bobine (rectangle horizontal)
\draw[line width=2.5pt, poporange] (-1.2,-0.5) rectangle (1.2,0.5);
% Courants sur faces actives (gauche et droite)
\draw[-{Stealth}, line width=1.2pt, popgreen!60!black] (-1.2,0.8) -- (-1.2,1.3) node[above]{\scriptsize $I$ monte};
\draw[-{Stealth}, line width=1.2pt, popgreen!60!black] (1.2,1.3) -- (1.2,0.8) node[above]{\scriptsize $I$ descend};
% Collecteur (deux demi-anneaux)
\draw[line width=2.5pt, popgold!80!black] (0,0.85) arc (90:270:0.35);
\draw[line width=2.5pt, popgold!80!black] (0,0.35) arc (-90:90:0.35);
% Balais
\draw[fill=popdark] (-0.6,1.3) rectangle (-0.2,1.5) node[above, font=\tiny] {Balai $+$};
\draw[fill=popdark] (0.2,1.3) rectangle (0.6,1.5) node[above, font=\tiny] {Balai $-$};
% Forces
\draw[-{Stealth}, line width=1.8pt, poppurple] (-1.2,0.5) -- (-1.2,1.7) node[above]{\scriptsize $\vec{F}_1$ (avant)};
\draw[-{Stealth}, line width=1.8pt, poppurple] (1.2,-0.5) -- (1.2,-1.7) node[below]{\scriptsize $\vec{F}_2$ (arrière)};
% Sens rotation
\draw[-{Stealth}, poppurple, dashed, line width=1pt] (1.8,0.8) arc (20:-160:2.0);
\node[poppurple] at (0,-2.2) {\scriptsize Rotation horaire};
\end{scope}

% Flèche transition
\draw[-{Stealth}, line width=2pt, popmuted] (3.2,0) -- (4.0,0) node[midway, above, font=\tiny] {1/4 tour};

% --- ÉTAPE 2 : Position verticale (point mort) ---
\begin{scope}[shift={(5.0,0)}]
\node[font=\small\bfseries, popblue] at (0,2.8) {Étape 2 : Point mort (bobine verticale)};
\fill[popred!70] (-2.5,-1.5) rectangle (-1.7,1.5);
\fill[popblue!70] (1.7,-1.5) rectangle (2.5,1.5);
\node[white, font=\bfseries] at (-2.1,0) {N};
\node[white, font=\bfseries] at (2.1,0) {S};
\foreach \y in {-1.0, 0, 1.0} {
  \draw[-{Stealth}, popmuted, line width=0.8pt] (-1.5,\y) -- (1.5,\y);
}
\node[popmuted] at (0,1.6) {$\vec{B}$};
% Bobine verticale
\draw[line width=2.5pt, poporange] (-0.5,-1.2) rectangle (0.5,1.2);
% Courants : faces actives maintenant haut/bas, // à B
\draw[-{Stealth}, line width=1.2pt, popgreen!60!black] (-0.5,1.5) -- (-0.5,1.2);
\draw[-{Stealth}, line width=1.2pt, popgreen!60!black] (0.5,1.2) -- (0.5,1.5);
% Collecteur : balais sur l'isolant (zone de coupure)
\draw[line width=2.5pt, popgold!80!black] (0,0.85) arc (90:270:0.35);
\draw[line width=2.5pt, popgold!80!black] (0,0.35) arc (-90:90:0.35);
\draw[fill=popdark] (-0.6,1.3) rectangle (-0.2,1.5);
\draw[fill=popdark] (0.2,1.3) rectangle (0.6,1.5);
% Pas de forces (ou nulles)
\node[poppurple, align=center] at (0,-2.0) {\scriptsize $\vec{I} \parallel \vec{B}$\scriptsize $\Rightarrow$ $\vec{F}=0$\scriptsize Couple nul};
\end{scope}

% Flèche transition
\draw[-{Stealth}, line width=2pt, popmuted] (8.2,0) -- (9.0,0) node[midway, above, font=\tiny] {1/4 tour};

% --- ÉTAPE 3 : Après point mort, courant inversé ---
\begin{scope}[shift={(10.0,0)}]
\node[font=\small\bfseries, popblue] at (0,2.8) {Étape 3 : Après point mort, courant inversé};
\fill[popred!70] (-2.5,-1.5) rectangle (-1.7,1.5);
\fill[popblue!70] (1.7,-1.5) rectangle (2.5,1.5);
\node[white, font=\bfseries] at (-2.1,0) {N};
\node[white, font=\bfseries] at (2.1,0) {S};
\foreach \y in {-1.0, 0, 1.0} {
  \draw[-{Stealth}, popmuted, line width=0.8pt] (-1.5,\y) -- (1.5,\y);
}
\node[popmuted] at (0,1.6) {$\vec{B}$};
% Bobine horizontale (autre face visible)
\draw[line width=2.5pt, poporange] (-1.2,-0.5) rectangle (1.2,0.5);
% Courants INVERSÉS
\draw[-{Stealth}, line width=1.2pt, popgreen!60!black] (-1.2,1.3) -- (-1.2,0.8) node[above]{\scriptsize $I$ descend};
\draw[-{Stealth}, line width=1.2pt, popgreen!60!black] (1.2,0.8) -- (1.2,1.3) node[above]{\scriptsize $I$ monte};
% Collecteur : balais ont changé de demi-anneau
\draw[line width=2.5pt, popgold!80!black] (0,0.85) arc (90:270:0.35);
\draw[line width=2.5pt, popgold!80!black] (0,0.35) arc (-90:90:0.35);
\draw[fill=popdark] (-0.6,1.3) rectangle (-0.2,1.5) node[above, font=\tiny] {Balai $+$};
\draw[fill=popdark] (0.2,1.3) rectangle (0.6,1.5) node[above, font=\tiny] {Balai $-$};
% Forces INVERSÉES
\draw[-{Stealth}, line width=1.8pt, poppurple] (-1.2,-0.5) -- (-1.2,-1.7) node[below]{\scriptsize $\vec{F}_1$ (arrière)};
\draw[-{Stealth}, line width=1.8pt, poppurple] (1.2,0.5) -- (1.2,1.7) node[above]{\scriptsize $\vec{F}_2$ (avant)};
% Sens rotation maintenu
\draw[-{Stealth}, poppurple, dashed, line width=1pt] (1.8,0.8) arc (20:-160:2.0);
\node[poppurple] at (0,-2.2) {\scriptsize Rotation horaire (même sens)};
\end{scope}
\end{tikzpicture}
\end{center}
\legende{Fonctionnement pas à pas. 1) Bobine horizontale : forces max, couple horaire. 2) Bobine verticale (point mort) : forces nulles, le collecteur coupe le courant. 3) Bobine horizontale inversée : le collecteur a inversé le courant, les forces s'inversent, le couple reste horaire. La rotation est continue.}
\end{popfigure}

\begin{methode}[Rédaction type BEPC : Expliquer le fonctionnement du moteur à courant continu]
\begin{enumerate}
\item Le stator (aimant ou électroaimant) crée un champ magnétique $\vec{B}$ entre ses pôles N et S.
\item Le rotor (bobine sur noyau fer) est alimenté en courant continu via les balais et le collecteur. Ses deux faces actives, perpendiculaires à $\vec{B}$, sont parcourues par des courants de sens opposés.
\item Elles subissent deux forces de Laplace de sens opposés (règle des trois doigts) qui forment un \textbf{couple de rotation} : le rotor tourne.
\item Au passage par la position d'équilibre (bobine parallèle à $\vec{B}$), les forces deviennent nulles. Le collecteur (demi-anneaux + balais) \textbf{inverse le sens du courant} dans la bobine.
\item Les forces s'inversent au bon moment : le couple garde le même sens, la rotation devient \textbf{continue}.
\end{enumerate}
\end{methode}

\coursec{Bilan énergétique, rendement et sécurité}

\courssub{Chaîne énergétique et rendement}

Le moteur est un convertisseur. L'énergie électrique absorbée ($E_e$ ou puissance $P_e$) se répartit en énergie mécanique utile ($E_m$, $P_m$) et en énergie thermique dissipée ($E_{th}$, $P_{th}$) par effet Joule dans les enroulements ($R \cdot I^2$) et frottements mécaniques.

\begin{popfigure}
\begin{center}
\begin{tikzpicture}[every node/.style={font=\small}]
% Source
\draw[fill=popblueL, rounded corners=4pt, line width=1pt, popblue] (0,-0.8) rectangle (3.8,0.8);
\node[align=center] at (1.9,0) {\textbf{Énergie électrique}\\ \scriptsize Réseau ENEO / Batterie / Pile\\ \scriptsize Puissance absorbée $P_e$};
% Flèche entrée
\draw[-{Stealth}, line width=2.5pt, popblue] (3.8,0) -- (5.0,0);
% Moteur
\draw[fill=poporangeL, rounded corners=4pt, line width=1.5pt, poporange!70!black] (5.0,-1.2) rectangle (8.5,1.2);
\node[align=center] at (6.75,0) {\textbf{MOTEUR ÉLECTRIQUE}\\ \scriptsize Conversion};
% Sortie Mécanique
\draw[-{Stealth}, line width=2.5pt, popgreen!60!black] (8.5,0.5) -- (10.5,0.5);
\draw[fill=popgreenL, rounded corners=4pt, line width=1pt, popgreen!70!black] (10.5,-0.3) rectangle (14.5,1.3);
\node[align=center] at (12.5,0.5) {\textbf{Énergie mécanique utile}\\ \scriptsize Rotation arbre (ventilateur, pompe...)\\ \scriptsize Puissance $P_m$};
% Sortie Thermique
\draw[-{Stealth}, line width=2.5pt, popred!80!black] (8.5,-0.5) -- (10.5,-0.5);
\draw[fill=poppinkL, rounded corners=4pt, line width=1pt, popred!70!black] (10.5,-1.3) rectangle (14.5,0.3);
\node[align=center] at (12.5,-0.5) {\textbf{Énergie thermique (pertes)}\\ \scriptsize Échauffement (Joule + frottements)\\ \scriptsize Puissance $P_{th}$};
\end{tikzpicture}
\end{center}
\legende{Chaîne énergétique du moteur électrique. La conservation de l'énergie impose $P_e = P_m + P_{th}$.}
\end{popfigure}

\begin{definition}[Rendement d'un moteur]
Le \cle{rendement} $\eta$ (lettre grecqueêta) est le rapport de la puissance mécanique utile fournie à la puissance électrique absorbée :
\[ \eta = \frac{P_m}{P_e} = \frac{P_e - P_{th}}{P_e} = 1 - \frac{P_{th}}{P_e} \]
Il est sans unité, souvent exprimé en \textbf{pourcentage} ($\eta \times 100$). Un bon moteur industriel a $\eta \approx 80\% \text{ à } 95\%$. Un petit moteur de jouet peut avoir $\eta < 50\%$.
\end{definition}

\begin{exemplebox}[Calcul de rendement — Pompe de forage]
Une pompe de forage immergée (type \emph{Grundfos} ou \emph{Pedrollo}) est alimentée par le réseau ENEO (\SI{230}{V}). Sa plaque signalétique indique : $P_e = \SI{750}{W}$ (puissance absorbée), $P_m = \SI{600}{W}$ (puissance mécanique à l'arbre).
\begin{align*}
P_{th} &= P_e - P_m = 750 - 600 = \SI{150}{W} \\
\eta &= \frac{600}{750} = 0,80 = 80\,\%
\end{align*}
\textbf{Interprétation :} 80\,\% de l'électricité payée à ENEO sert à pomper l'eau, 20\,\% chauffe le moteur (pertes).
\end{exemplebox}

\courssub{Pourquoi un moteur bloqué grille (surchauffe)}

\begin{experience}[Démonstration : ventilateur bloqué (à faire en démo prof, pas élèves)]
\textbf{Matériel :} petit moteur \SI{12}{V} (ventilateur PC), alimentation \SI{12}{V} \SI{2}{A}, thermomètre infrarouge ou main (prudence), objet pour bloquer l'hélice (bâton isolant).
\textbf{Protocole :} 1) Moteur libre : mesure $I$, température stable. 2) Bloquer l'hélice fermement 10~s : observer $I$ qui monte fort, température qui grimpe vite, odeur de vernis. 3) Débrancher immédiatement.
\textbf{Interprétation :} Rotor bloqué $\Rightarrow$ pas de rotation $\Rightarrow$ pas d'énergie mécanique produite ($P_m = 0$). Toute la puissance $P_e = U \times I$ devient chaleur $P_{th}$ dans les fils de cuivre (résistance $R$). $I$ augmente car la \emph{force contre-électromotrice} (FCEM) disparaît (voir Terminale). L'isolant fond $\rightarrow$ court-circuit $\rightarrow$ moteur grillé.
\end{experience}

\begin{savaistu}[Sécurité au quotidien : « Mon ventilateur sent le brûlé ! »]
Si un ventilateur, une pompe, une meuleuse ou une perceuse force (charge trop lourde, corps étranger coincé), \textbf{le courant augmente brutalement} et l'échauffement est exponentiel.
\begin{itemize}
\item \textbf{Réflexe vital :} \textbf{Couper l'alimentation immédiatement} (interrupteur, prise, disjoncteur).
\item \textbf{Ne jamais} maintenir un appareil bloqué en espérant qu'il « passe ».
\item Les moteurs modernes intègrent souvent une \textbf{sonde thermique} (Klixon) qui coupe le courant si $T > \SI{130}{\celsius}$ environ. Mais ne compte pas uniquement dessus !
\end{itemize}
\end{savaistu}

\coursec{Utilisation, plaque signalétique et sécurité}

\courssub{Lire une plaque signalétique (obligatoire sur tout moteur)}

La plaque rivetée sur la carcasse donne les conditions d'utilisation nominales (point de fonctionnement optimal). Exemple pour un moteur triphasé asynchrone de pompe \SI{1,5}{kW} (courant au Cameroun) :

\begin{center}
\begin{tabularx}{\linewidth}{|l|X|l|X|}\hline
\rowcolor{popblueL}
\textbf{Paramètre} & \textbf{Signification} & \textbf{Exemple} & \textbf{Unité / Remarque} \\ \hline
Tension $U_n$ & Tension d'alimentation nominale & 380 (ou 400) & V (Volts) — Triphasé : 380~V entre phases \\ \hline
Courant $I_n$ & Courant absorbé en charge nominale & 3,2 & A (Ampères) — Prévoir disjoncteur $> I_n$ \\ \hline
Puissance $P_n$ & Puissance \textbf{mécanique} utile à l'arbre & 1,5 & kW (kilowatts) — \textbf{Pas} la puissance électrique ! \\ \hline
Vitesse $n_n$ & Vitesse de rotation nominale & 2880 & tr/min (tours par minute) \\ \hline
Fréquence $f$ & Fréquence du réseau & 50 & Hz (Hertz) — Cameroun = 50~Hz \\ \hline
$\cos \varphi$ (cos phi) & Facteur de puissance & 0,82 & Sans unité — Qualité du couplage réseau \\ \hline
Rendement $\eta$ & Rendement à charge nominale & 84,5 & \% — Ici $\eta = 84,5\,\%$ \\ \hline
Classe Isolation & Résistance thermique de l'isolant & F (\SI{155}{\celsius}) ou B (\SI{130}{\celsius}) & Lettre code (A, E, B, F, H) \\ \hline
Indice IP & Protection corps solides / liquides & IP55 & IP55 = protégé poussières + jets d'eau \\ \hline
Connexion & Schéma de raccordement & $\triangle$ / Y & Triangle (basse tension) / Étoile (haute tension) \\ \hline
\end{tabularx}
\end{center}

\begin{attention}[Erreur fréquente : Confondre $P_n$ (mécanique) et $P_e$ (électrique)]
Sur la plaque, \textbf{$P_n$ est la puissance MÉCANIQUE} (celle qui sort de l'arbre). La puissance électrique absorbée $P_e$ est \textbf{plus grande} : $P_e = P_n / \eta$. Pour le dimensionnement du disjoncteur et du câble, on utilise le \textbf{courant nominal $I_n$} (donné sur la plaque), pas un calcul approximatif.
\end{attention}

\courssub{Consignes de sécurité associées}

\begin{itemize}
\item \textbf{Mise à la terre (prise 3 broches / borne PE) :} Obligatoire pour tout moteur métallique. Si un fil nu touche la carcasse, le courant fuit vers la terre $\rightarrow$ disjoncteur différentiel (30~mA) coupe instantanément. \textbf{Vital au Cameroun où le neutre n'est pas toujours fiable.}
\item \textbf{Protection magnéto-thermique (disjoncteur) :} Calibré sur $I_n$ (ou légèrement au-dessus pour le démarrage). Protège contre surcharge (thermique) et court-circuit (magnétique).
\item \textbf{Ventilation :} Ne jamais boucher les grilles d'aération du moteur (poussières, sacs, murs collés). La surchauffe détruit l'isolant.
\item \textbf{Couple de démarrage :} Un moteur démarre avec un courant \textbf{5 à 7 fois $I_n$} pendant quelques dixièmes de seconde. Le disjoncteur doit supporter ce pic (courbe D ou K pour moteurs).
\item \textbf{Sens de rotation :} Vérifier le sens (flèche sur la carcasse) avant de coupler la pompe/ventilateur. Inverser deux phases (triphasé) ou bobinage auxiliaire (monophasé) pour inverser le sens.
\end{itemize}

\begin{savaistu}[Le démarreur « étoile-triangle » pour gros moteurs]
Pour les gros moteurs (ex. compresseur usine, pompe irrigation \SI{> 7,5}{kW}), le courant de démarrage direct (5-7 $I_n$) ferait disjoncter le réseau ou ferait chuter la tension (les lampes clignotent dans le quartier). On utilise un démarreur \textbf{étoile-triangle} : au démarrage, les bobines sont couplées en \textbf{étoile} (courant divisé par 3, couple divisé par 3), puis après quelques secondes, on bascule en \textbf{triangle} (pleine puissance). C'est de la technologie pure, très courant dans l'industrie camerounaise (SONARA, brasseries, usines à bois).
\end{savaistu}

\coursec{Complément culturel : La réversibilité moteur / générateur}

\begin{savaistu}[Hors programme strict : Le moteur devient générateur]
Le dispositif est \textbf{réversible}. Si l'on \textbf{entraîne mécaniquement} l'axe du rotor (manivelle, chute d'eau, vent, moteur thermique), la bobine tourne dans le champ magnétique $\vec{B}$ : une \textbf{tension alternative} apparaît à ses bornes (loi de Faraday-Lenz). C'est le principe de :
\begin{itemize}
\item La \cle{dynamo} de vélo (alternateur à aimant permanent) : éclairage sans pile.
\item L'\cle{alternateur} de voiture/moto : recharge la batterie moteur tournant.
\item Le \cle{groupe électrogène} : moteur thermique (essence/gasoil SONARA) + alternateur $\rightarrow$ courant alternatif \SI{230}{V} / \SI{380}{V}.
\item Le \cle{freinage régénératif} (voitures électriques, trains) : le moteur devient générateur pour freiner et recharger la batterie.
\end{itemize}
\textbf{Au BEPC 3ème :} Seul le fonctionnement \textbf{MOTEUR} (électrique $\rightarrow$ mécanique) est exigible. La réversibilité est une ouverture culturelle pour comprendre le monde énergétique.
\end{savaistu}

\begin{exoresolu}[Analyse complète : Moteur de ventilateur de plafond]
Un ventilateur de plafond (marque \emph{KDK} ou \emph{Panasonic}, courant au Cameroun) porte la plaque : $U = \SI{230}{V}$~; $P_n = \SI{75}{W}$~; $\eta = 72\,\%$~; $n = \SI{300}{tr/min}$~; $\cos\varphi = 0,95$.
\begin{enumerate}
\item Quelle est la puissance électrique absorbée $P_e$ ?
\item Calculer le courant nominal $I_n$ (monophasé : $P_e = U \cdot I_n \cdot \cos\varphi$).
\item Quelle est la puissance thermique dissipée $P_{th}$ ?
\item Si le ventilateur tourne dans le vide (pas de pale), que se passe-t-il pour $P_m$, $P_{th}$, $I$ ?
\end{enumerate}
\tcblower
\textbf{Solution détaillée.}
\begin{enumerate}
\item $P_n$ est la puissance mécanique. $\eta = P_n / P_e \Rightarrow P_e = P_n / \eta = 75 / 0,72 = \SI{104,2}{W}$.
\item $P_e = U \cdot I_n \cdot \cos\varphi \Rightarrow I_n = P_e / (U \cdot \cos\varphi) = 104,2 / (230 \times 0,95) = \SI{0,476}{A} \approx \SI{0,48}{A}$.
\item $P_{th} = P_e - P_n = 104,2 - 75 = \SI{29,2}{W}$ (pertes Joule + frottements roulements/air).
\item Sans pale, la charge mécanique est quasi nulle (juste frottements). $P_m \approx 0$. Le moteur tourne plus vite (proche de la vitesse de synchronisme), le courant $I$ \textbf{diminue} (contrairement au blocage rotor). $P_e$ diminue, $P_{th}$ diminue. \textbf{Attention :} ne pas confondre « rotor bloqué » (surchauffe destructive) et « moteur à vide » (surchauffe modérée, mais pas de ventilation du moteur lui-même $\rightarrow$ risque à long terme).
\end{enumerate}
\end{exoresolu}

\begin{aretenir}[L'essentiel de la leçon — Moteur électrique (3ème PCT)]
\begin{itemize}
\item \textbf{Rôle :} Convertisseur \textbf{Énergie électrique $\rightarrow$ Énergie mécanique (rotation) + Énergie thermique (pertes)}.
\item \textbf{Constitution :} Stator (champ B, aimant ou électroaimant) + Rotor (noyau fer + bobines + collecteur) + Balais (charbons graphite).
\item \textbf{Principe :} Bobine dans champ B $\rightarrow$ forces de Laplace sur faces actives (courants opposés) $\rightarrow$ \textbf{Couple de rotation}.
\item \textbf{Point mort :} Bobine $\parallel$ B $\rightarrow$ Forces nulles. Sans collecteur, le moteur s'arrête.
\item \textbf{Collecteur (commutateur) :} Inverse le courant dans la bobine \textbf{à chaque demi-tour} $\rightarrow$ forces s'inversent au bon moment $\rightarrow$ \textbf{Rotation continue unidirectionnelle}.
\item \textbf{Bilan énergie :} $P_e = P_m + P_{th}$. Rendement $\eta = P_m / P_e < 1$ (jamais 100\,\%).
\item \textbf{Danger :} Rotor bloqué $\Rightarrow P_m = 0 \Rightarrow P_e = P_{th}$ $\rightarrow$ surchauffe rapide $\rightarrow$ destruction (grillage). \textbf{Couper le courant immédiatement.}
\item \textbf{Plaque signalétique :} Donne $U_n, I_n, P_n (\text{mécanique}), n_n, \cos\varphi, \eta, \text{Classe IP/Isolation}$. Indispensable pour choisir protection (disjoncteur, section câble) et raccordement ($\triangle$/Y).
\item \textbf{Sécurité :} Mise à la terre obligatoire, ventilation libre, protection différentielle 30~mA, disjoncteur calibré $I_n$.
\item \textbf{(Culture) Réversibilité :} Entraîné mécaniquement $\rightarrow$ Générateur (dynamo, alternateur, groupe électrogène).
\end{itemize}
\end{aretenir}

\coursec{Utilisation, plaque signalétique et sécurité}

Nous savons maintenant comment un moteur électrique transforme l'énergie électrique en énergie mécanique : un courant parcourt une bobine placée dans le champ d'un aimant, et les forces électromagnétiques font tourner le rotor. Mais dans la vie courante, il ne suffit pas de savoir \emph{comment} un moteur tourne : il faut savoir le \emph{choisir}, le \emph{brancher correctement} et l'\emph{entretenir}. Au marché de Douala ou dans un atelier de réparation de Yaoundé, le réparateur lit d'abord la petite étiquette métallique rivetée sur le moteur avant de décider quoi faire. Cette étiquette, c'est la \cle{plaque signalétique}. Cette dernière partie de la leçon t'apprend à la lire, à choisir un moteur adapté à une source, et à respecter les consignes de sécurité qui protègent la machine... et surtout les personnes.

\courssub{La plaque signalétique : la carte d'identité du moteur}

\begin{definition}[Plaque signalétique]
La \cle{plaque signalétique} d'un moteur électrique est une petite plaque (ou étiquette) fixée sur le carter du moteur par le constructeur. Elle indique les caractéristiques nominales de fonctionnement, c'est-à-dire les valeurs pour lesquelles le moteur a été conçu :
\begin{itemize}
\item la \cle{tension d'alimentation} $U$, exprimée en volts (\si{\volt}) ;
\item la \cle{puissance} $P$, exprimée en watts (\si{\watt}) ;
\item la \cle{vitesse de rotation} de l'arbre, exprimée en tours par minute (tr/min) ;
\item parfois aussi : l'intensité nominale $I$ en ampères (\si{\ampere}), le type de courant (continu $\mathrm{=}$ ou alternatif $\sim$), la fréquence (\SI{50}{Hz} au Cameroun) et l'indice de protection contre l'eau et la poussière (code IP).
\end{itemize}
\end{definition}

\begin{popfigure}
\begin{tikzpicture}
\draw[fill=popgoldL,draw=popdark,thick,rounded corners=2pt] (0,0) rectangle (8.6,4.4);
\foreach \x in {0.25,8.35} \foreach \y in {0.25,4.15} \fill[popmuted] (\x,\y) circle (0.09);
\draw[popdark,thick] (0.6,3.7) -- (8.0,3.7);
\node[font=\bfseries\small,text=popdark] at (4.3,4.0) {MOTEUR ÉLECTRIQUE --- Type MEC 220};
\draw[popdark] (0.6,0.6) rectangle (8.0,3.5);
\draw[popdark] (4.3,0.6) -- (4.3,3.5);
\node[anchor=west,font=\small] at (0.8,3.05) {Tension $U$};
\node[anchor=west,font=\small\bfseries,text=popblue] at (4.6,3.05) {\SI{220}{V} $\sim$ \SI{50}{Hz}};
\node[anchor=west,font=\small] at (0.8,2.35) {Puissance $P$};
\node[anchor=west,font=\small\bfseries,text=popblue] at (4.6,2.35) {\SI{50}{W}};
\node[anchor=west,font=\small] at (0.8,1.65) {Vitesse $N$};
\node[anchor=west,font=\small\bfseries,text=popblue] at (4.6,1.65) {\SI{2800}{tr/min}};
\node[anchor=west,font=\small] at (0.8,0.95) {Protection};
\node[anchor=west,font=\small\bfseries,text=popblue] at (4.6,0.95) {IP 44};
\draw[->,thick,poporange] (9.4,3.05) -- (8.15,3.05);
\node[anchor=west,font=\footnotesize,text=poporange] at (9.5,3.05) {à brancher sur le secteur};
\draw[->,thick,poporange] (9.4,2.35) -- (8.15,2.35);
\node[anchor=west,font=\footnotesize,text=poporange] at (9.5,2.35) {puissance absorbée};
\draw[->,thick,poporange] (9.4,1.65) -- (8.15,1.65);
\node[anchor=west,font=\footnotesize,text=poporange] at (9.5,1.65) {tours par minute};
\end{tikzpicture}
\legende{Exemple de plaque signalétique d'un moteur : tension d'alimentation \SI{220}{V} (courant alternatif du réseau ENEO), puissance \SI{50}{W}, vitesse de rotation \SI{2800}{tr/min}.}
\end{popfigure}

\begin{methode}[Lire et exploiter une plaque signalétique]
Pour savoir si un moteur convient à ton installation, procède en quatre étapes :
\begin{enumerate}
\item \textbf{Repérer la tension $U$} (en volts) : c'est la tension que la source doit fournir.
\item \textbf{Repérer la puissance $P$} (en watts) : elle indique la « force » du moteur et permet de calculer l'intensité absorbée par $I = \dfrac{P}{U}$.
\item \textbf{Repérer la vitesse de rotation} (en tr/min) : elle doit convenir à l'appareil entraîné (ventilateur, pompe, meule).
\item \textbf{Vérifier la compatibilité avec la source} : la tension de la source doit être \emph{égale} à la tension nominale du moteur. Un moteur de \SI{12}{V} se branche sur une batterie de \SI{12}{V} ; un moteur de \SI{220}{V} se branche sur le secteur ENEO.
\end{enumerate}
\end{methode}

\begin{savaistu}[Que signifie « tr/min » ?]
La mention « tr/min » indique la \cle{vitesse de rotation} de l'arbre du moteur, c'est-à-dire le nombre de tours complets effectués en une minute. Un ventilateur de plafond tourne à environ \num{1000} à \SI{1500}{tr/min}, le moteur d'une meuleuse d'atelier peut atteindre \SI{2800}{tr/min}, tandis que le moteur d'une moto-taxi dépasse \SI{5000}{tr/min} à plein régime ! Plus la vitesse est élevée, plus les risques en cas de contact avec les parties tournantes sont grands.
\end{savaistu}

\courssub{Choisir un moteur adapté à la source}

Un moteur n'est pas un appareil « universel » : il est construit pour fonctionner sous une tension précise. Le choix d'un moteur se fait donc en fonction de la source d'énergie disponible.

\begin{propriete}[Règle de compatibilité]
La tension nominale du moteur doit être \textbf{égale} à la tension de la source d'alimentation :
\begin{itemize}
\item moteur \SI{12}{V} (courant continu) $\rightarrow$ batterie d'accumulateurs de voiture ou de moto (\SI{12}{V}) ;
\item moteur \SI{220}{V} (courant alternatif) $\rightarrow$ prise de courant du réseau ENEO ou groupe électrogène (\SI{220}{V}).
\end{itemize}
\end{propriete}

\begin{exemplebox}[Choisir une moto-pompe pour un forage]
Un planteur de la région de l'Ouest veut installer une pompe électrique pour remonter l'eau d'un forage. Dans son village, il n'y a pas de réseau ENEO, mais il possède un groupe électrogène qui fournit du \SI{220}{V} alternatif. Il choisit donc une moto-pompe dont la plaque indique « \SI{220}{V} $\sim$ », et non une pompe de \SI{12}{V} prévue pour batterie. S'il avait voulu utiliser des panneaux solaires avec des batteries, il aurait au contraire choisi une pompe \SI{12}{V} ou \SI{24}{V} en courant continu.
\end{exemplebox}

\begin{attention}[Erreur fréquente : brancher un moteur 12 V sur le secteur 220 V]
Brancher un moteur prévu pour \SI{12}{V} sur une prise de \SI{220}{V} provoque sa \textbf{destruction quasi immédiate} : le courant devient presque 20 fois trop intense, les enroulements chauffent brutalement, l'isolant fond et le moteur peut fumer, prendre feu ou exploser ses balais. Il y a aussi un grave \textbf{danger d'électrocution et d'incendie} pour l'utilisateur. Inversement, un moteur \SI{220}{V} branché sur \SI{12}{V} ne tournera pas (ou très faiblement). Vérifie \emph{toujours} la plaque signalétique avant de brancher.
\end{attention}

\begin{exoresolu}[Calculer l'intensité absorbée par un moteur]
Énoncé. Sur la plaque signalétique d'un moteur de moulin, on lit : « \SI{220}{V} $\sim$ \SI{110}{W} ». Calculer l'intensité $I$ du courant qui traverse ce moteur en fonctionnement normal.
\tcblower
Solution détaillée. On connaît la puissance $P = \SI{110}{W}$ et la tension $U = \SI{220}{V}$. La relation entre puissance, tension et intensité est :
\[ P = U \cdot I \quad \text{donc} \quad I = \frac{P}{U} \]
Application numérique :
\[ I = \frac{110}{220} = \SI{0,5}{A} \]
Le moteur est traversé par un courant d'intensité $I = \SI{0,5}{A}$ en fonctionnement normal. Ce calcul permet de choisir un fusible et des fils de section adaptés.
\end{exoresolu}

\courssub{Les consignes de sécurité}

Un moteur électrique présente \emph{deux} types de dangers : le danger \textbf{électrique} (le courant) et le danger \textbf{mécanique} (les parties tournantes). Les consignes de sécurité visent à éviter les deux.

\begin{popfigure}
\begin{tikzpicture}
% Danger électrique
\draw[fill=popgold,draw=popdark,very thick] (0,1.9) -- (1.65,-1.0) -- (-1.65,-1.0) -- cycle;
\draw[fill=popdark] (-0.12,1.1) -- (0.35,1.1) -- (-0.05,0.15) -- (0.35,0.15) -- (-0.45,-0.75) -- (-0.08,-0.18) -- (-0.45,-0.18) -- cycle;
\node[font=\footnotesize\bfseries,text=popdark] at (0,-1.5) {Danger électrique};
% Parties tournantes
\begin{scope}[xshift=5.2cm]
\draw[fill=popgold,draw=popdark,very thick] (0,1.9) -- (1.65,-1.0) -- (-1.65,-1.0) -- cycle;
\foreach \a in {0,60,120,180,240,300} \draw[fill=popdark,rotate=\a] (0,0.18) ellipse (0.16 and 0.55);
\fill[popgold] (0,0.18) circle (0.22);
\fill[popdark] (0,0.18) circle (0.1);
\node[font=\footnotesize\bfseries,text=popdark] at (0,-1.5) {Parties tournantes};
\end{scope}
% Interdiction de toucher
\begin{scope}[xshift=10.4cm]
\draw[draw=popink,very thick,fill=white] (0,0.45) circle (1.45);
\draw[popink,very thick] (-1.02,1.47) -- (1.02,-0.57);
\draw[fill=popdark] (-0.5,-0.35) rectangle (0.15,0.55);
\draw[popdark,very thick] (0.15,0.35) -- (0.75,0.95);
\draw[popdark,very thick] (0.15,0.15) -- (0.8,0.55);
\draw[popdark,very thick] (0.15,-0.05) -- (0.75,0.15);
\node[font=\footnotesize\bfseries,text=popink] at (0,-1.5) {Ne pas toucher};
\end{scope}
\end{tikzpicture}
\legende{Pictogrammes de sécurité rencontrés sur les moteurs et machines : danger électrique (éclair), danger lié aux parties tournantes (engrenage), interdiction de toucher (panneau rond barré).}
\end{popfigure}

\begin{aretenir}[Consignes de sécurité mécanique]
\begin{itemize}
\item \textbf{Ne jamais bloquer l'arbre en marche} : un moteur bloqué n'arrête pas de « forcer » ; le courant augmente fortement, le moteur chauffe et peut brûler.
\item \textbf{Ne pas toucher un moteur en fonctionnement} : ni l'arbre, ni la poulie, ni la courroie ; un doigt ou un pagne peut être happé.
\item \textbf{Protéger les parties tournantes par un carter} : le carter (couvercle de protection) doit toujours être en place avant la mise en marche.
\end{itemize}
\end{aretenir}

\begin{aretenir}[Consignes de sécurité électrique]
\begin{itemize}
\item \textbf{Ne jamais alimenter un moteur sous une tension supérieure à sa tension nominale} (voir la plaque signalétique).
\item \textbf{Éviter l'humidité} : ne pas installer le moteur dans un lieu mouillé, ne pas le manipuler avec les mains mouillées ; l'eau favorise les courts-circuits et l'électrocution.
\item \textbf{Débrancher avant toute intervention} : nettoyage, graissage, changement de courroie se font toujours moteur arrêté et fiche retirée de la prise.
\end{itemize}
\end{aretenir}

\begin{attention}[Pièges à éviter]
\begin{itemize}
\item Croire qu'un moteur arrêté est sans danger alors qu'il est encore branché : il peut redémarrer brusquement. \textbf{Débranche d'abord.}
\item Retirer le carter « pour voir tourner » : c'est la cause de nombreux accidents dans les ateliers.
\item Utiliser une rallonge abîmée ou des fils dénudés pour alimenter le moulin du quartier : risque de court-circuit et d'incendie.
\end{itemize}
\end{attention}

\courssub{Entretien simple d'un moteur}

Un moteur bien entretenu dure des années ; un moteur négligé grille en quelques mois. L'entretien courant, toujours réalisé \textbf{moteur débranché}, comprend trois gestes simples :

\begin{enumerate}
\item \textbf{Le dépoussiérage} : la poussière (farine dans un moulin, sciure dans un atelier de menuiserie) bouche les ouvertures de ventilation et fait chauffer le moteur. On la retire avec un pinceau sec ou un chiffon, jamais avec de l'eau.
\item \textbf{La vérification des balais} : dans un moteur à courant continu (ou universel), les balais (petits blocs de charbon) frottent sur le collecteur pour amener le courant au rotor. Ils s'usent : quand ils deviennent trop courts, le moteur fait des étincelles et tourne mal. Il faut alors les remplacer.
\item \textbf{Le graissage des paliers} : les paliers (ou roulements) guident la rotation de l'arbre. Quelques gouttes d'huile ou un peu de graisse réduisent le frottement, le bruit et l'usure.
\end{enumerate}

\begin{exemplebox}[L'entretien du moulin du quartier]
Le moulin à maïs du quartier fonctionne chaque jour. Son propriétaire applique une routine simple : chaque soir, il débranche le moteur et enlève la farine accumulée autour des ouïes de ventilation ; chaque semaine, il vérifie que la courroie n'est pas détendue et que le carter est bien fixé ; chaque mois, il graisse les paliers et fait contrôler les balais par le technicien. Résultat : moins de pannes, moins de frais, et aucun accident.
\end{exemplebox}

\begin{exercice}[À toi de jouer]
Sur la plaque d'un moteur de ventilateur, on lit : « \SI{220}{V} $\sim$ \SI{60}{W} \SI{1200}{tr/min} ».
\begin{enumerate}
\item Peut-on brancher ce ventilateur sur une batterie de \SI{12}{V} ? Justifier.
\item Calculer l'intensité du courant absorbé en fonctionnement normal.
\item Citer deux consignes de sécurité à respecter avant de nettoyer ce ventilateur.
\end{enumerate}
\end{exercice}

\begin{corrige}[Exercice : moteur de ventilateur]
\begin{enumerate}
\item Non. La tension nominale du moteur (\SI{220}{V} alternatif) doit être égale à la tension de la source. Une batterie de \SI{12}{V} continu est très inférieure : le moteur ne tournera pas normalement. Il faut le brancher sur le secteur \SI{220}{V}.
\item $I = \dfrac{P}{U} = \dfrac{60}{220} \approx \SI{0,27}{A}$.
\item Débrancher la fiche de la prise avant toute intervention ; attendre l'arrêt complet des pales et ne jamais toucher les parties tournantes en marche ; utiliser un chiffon sec (éviter l'humidité).
\end{enumerate}
\end{corrige}

\begin{aretenir}[L'essentiel de la section]
\begin{itemize}
\item La \cle{plaque signalétique} indique la tension $U$ (V), la puissance $P$ (W) et la vitesse de rotation (tr/min) : c'est la carte d'identité du moteur.
\item La tension de la source doit être \textbf{égale} à la tension nominale du moteur : \SI{12}{V} sur batterie, \SI{220}{V} sur le secteur.
\item L'intensité absorbée se calcule par $I = \dfrac{P}{U}$.
\item Sécurité : ne jamais bloquer l'arbre, ne pas toucher un moteur en marche, garder le carter en place, éviter l'humidité, débrancher avant toute intervention.
\item Entretien : dépoussiérer, vérifier les balais, graisser les paliers, toujours moteur débranché.
\end{itemize}
\end{aretenir}

\newpage

\coursec{Activité d'intégration}

\coursec{Le moteur électrique}

\courssub{Objectifs de la leçon}

À la fin de cette leçon, tu seras capable de :
\begin{itemize}
\item identifier sur un schéma les principaux constituants d'un moteur électrique à courant continu (\cle{stator}, \cle{rotor}, \cle{balais}, \cle{collecteur}) et expliquer le rôle de chacun ;
\item exploiter les indications d'une \cle{plaque signalétique} (tension $U$, puissance $P$) pour calculer l'intensité nominale $I = \dfrac{P}{U}$ ;
\item utiliser la \cle{chaîne énergétique} pour expliquer l'échauffement et la destruction d'un moteur dont l'arbre est bloqué ;
\item rédiger des consignes de sécurité et d'entretien adaptées à une situation concrète (pompe, ventilateur, moto-taxi, groupe électrogène).
\end{itemize}

\courssub{Situation déclenchante : la pompe de M. Essono}

M. Essono habite Bafoussam, dans la région de l'Ouest. Pour arroser son champ de tomates pendant la saison sèche, il a acheté une moto-pompe électrique qu'il branche sur le réseau ENEO de son quartier ($220$~V, $50$~Hz). La pompe aspire l'eau d'un forage et la refoule dans un tuyau d'arrosage.

Sur la carcasse du moteur, on lit la plaque signalétique suivante :

\begin{center}
\begin{tabularx}{\linewidth}{|l|X|}
\hline
\rowcolor{popblueL} \textbf{Plaque signalétique} & \textbf{Valeur} \\
\hline
Tension d'alimentation $U$ & \SI{220}{V} \\
\hline
Puissance utile indiquée $P$ & \SI{750}{W} \\
\hline
Fréquence & \SI{50}{Hz} \\
\hline
Indice de protection & IP 44 (protégé contre les corps solides $> 1$~mm et les projections d'eau) \\
\hline
\end{tabularx}
\end{center}

Un après-midi, un amas de feuilles et de sable bouche complètement le tuyau de sortie. L'eau ne sort plus, mais M. Essono, occupé au fond du champ, entend le moteur qui continue de ronronner. Au bout de quelques minutes, une odeur de brûlé se dégage, puis le moteur s'arrête définitivement. Le réparateur ouvre le moteur : le vernis isolant des fils de cuivre des bobinages a fondu, le moteur a \og grillé \fg{}.

\textbf{Problème :} pourquoi le moteur a-t-il grillé et comment l'éviter à l'avenir ?

\courssub{Étude technique : constitution d'un moteur électrique à courant continu}

\begin{experience}[Observation d'un moteur à courant continu (schéma simplifié)]
\textbf{Matériel :} schéma annoté ci-dessous (représentation conventionnelle d'un moteur à courant continu à aimants permanents).
\textbf{Observations :} on distingue deux parties principales : une partie fixe et une partie mobile.
\begin{popfigure}
\begin{center}
\begin{tikzpicture}[scale=0.95, every node/.style={font=\small}]
% Stator (aimants permanents)
\fill[popblue!30] (-3.3,-1.7) rectangle (-2.3,1.7);
\fill[popblue!30] (2.3,-1.7) rectangle (3.3,1.7);
\node[popblue] at (-2.8,0) {\textbf{N}};
\node[popblue] at (2.8,0) {\textbf{S}};
\node[below, popblue] at (-2.8,-1.9) {stator (pôle Nord)};
\node[below, popblue] at (2.8,-1.9) {stator (pôle Sud)};
% Champ magnétique (flèches)
\foreach \y in {-1.2,-0.6,0,0.6,1.2} {
  \draw[popblue, ->, thick] (-2.2,\y) -- (2.2,\y);
}
% Rotor (bobine)
\draw[very thick, poporange] (-1.3,-1.0) rectangle (1.3,1.0);
% Courant dans les côtés actifs
\draw[very thick, poporange, ->] (-1.3,-1.0) -- (-1.3,0.3) node[midway, left=2pt] {$I$};
\draw[very thick, poporange, ->] (1.3,1.0) -- (1.3,-0.3) node[midway, right=2pt] {$I$};
% Forces de Laplace (F)
\draw[popgreen, ->, very thick] (-1.3,0) -- (-2.0,0) node[left] {$\vec{F}_1$};
\draw[popgreen, ->, very thick] (1.3,0) -- (2.0,0) node[right] {$\vec{F}_2$};
\node[poporange, above] at (0,1.3) {rotor (bobine mobile)};
% Axe de rotation
\draw[thick, popdark] (0,-2.5) -- (0,2.5);
\fill[popdark] (0,0) circle (0.1);
\node[right, popdark] at (0,2.5) {arbre de sortie};
% Balais et collecteur (anneaux coupés)
\fill[popdark] (-0.6,-1.6) rectangle (-0.3,-1.2);
\fill[popdark] (0.3,-1.6) rectangle (0.6,-1.2);
\node[below, popdark] at (0,-1.8) {balais (charbon) + collecteur (anneaux isolés)};
% Alimentation
\draw[thick] (-0.45,-1.4) -- (-0.45,-3.2) -- (-1.8,-3.2);
\draw[thick] (0.45,-1.4) -- (0.45,-3.2) -- (1.8,-3.2);
\draw[thick] (-1.8,-3.2) -- (-1.8,-4.0);
\draw[thick] (1.8,-3.2) -- (1.8,-4.0);
\draw[thick] (-1.8,-4.0) circle (0.4);
\draw[thick] (1.8,-4.0) circle (0.4);
\node at (-1.8,-4.0) {$+$};
\node at (1.8,-4.0) {$-$};
\node[below] at (0,-4.5) {alimentation continue (pile ou redresseur)};
\end{tikzpicture}
\end{center}
\legende{Schéma de principe d'un moteur à courant continu : le stator crée un champ magnétique radial $\vec{B}$ ; la bobine (rotor) parcourue par un courant $I$ subit des forces de Laplace $\vec{F}_1$, $\vec{F}_2$ qui créent un couple moteur. Les balais et le collecteur inversent le courant dans la bobine à chaque demi-tour pour maintenir le sens de rotation.}
\end{popfigure}
\end{experience}

\begin{definition}[Les quatre constituants essentiels]
\begin{itemize}
\item \textbf{Stator (partie fixe) :} aimants permanents (ou électro-aimants) qui créent un champ magnétique $\vec{B}$ constant dans l'entrefer. Sur le schéma : pôles N et S bleus.
\item \textbf{Rotor ou armature (partie tournante) :} bobine de fil de cuivre enroulée sur un noyau en fer doux ; c'est le récepteur d'énergie électrique. Sur le schéma : rectangle orange.
\item \textbf{Collecteur (ou collecteur à bagues) :} anneau métallique fendu en deux demi-anneaux isolés l'un de l'autre, solidaire de l'axe. Il inverse le sens du courant dans la bobine à chaque demi-tour.
\item \textbf{Balais (ou charbons) :} pièces en graphite (charbon) fixes, pressées sur le collecteur ; elles assurent le contact électrique entre le circuit extérieur (alimentation) et le collecteur en rotation.
\end{itemize}
\end{definition}

\courssub{Étude théorique : action d'un aimant sur une bobine parcourue par un courant}

\begin{propriete}[Force de Laplace et couple moteur]
Lorsqu'un conducteur rectiligne de longueur $l$, parcouru par un courant $I$, est placé dans un champ magnétique $\vec{B}$, il subit une force $\vec{F}$ (force de Laplace) :
\begin{itemize}
\item \textbf{Direction :} perpendiculaire au plan formé par le conducteur et $\vec{B}$.
\item \textbf{Sens :} donné par la \textbf{règle des trois doigts de la main droite} (pouce = courant $I$, index = champ $\vec{B}$, majeur = force $\vec{F}$).
\item \textbf{Valeur :} $F = B \cdot I \cdot l$ (si le conducteur est perpendiculaire à $\vec{B}$).
\end{itemize}
Dans le moteur, les deux côtés actifs de la bobine sont parcourus par des courants de sens opposés et subissent des forces opposées : elles forment un \textbf{couple} qui fait tourner le rotor.
\end{propriete}

\begin{propriete}[Rôle du collecteur : inversion du courant]
Sans collecteur, la bobine s'arrêterait à la position verticale (couple nul). Le collecteur (demi-anneaux tournants + balais fixes) \textbf{inverse le sens du courant dans la bobine} exactement quand elle passe la position verticale. Les forces changent de sens, le couple reste dans le même sens de rotation : le mouvement devient continu.
\end{propriete}

\begin{aretenir}[Chaîne énergétique du moteur électrique]
\[
\boxed{\text{Énergie électrique absorbée} \;\longrightarrow\; \text{Énergie mécanique utile (rotation)} \;+\; \text{Énergie thermique (pertes par effet Joule)}}
\]
\begin{itemize}
\item L'énergie électrique fournie par la source ($E_{elec} = P \cdot t = U \cdot I \cdot t$) est convertie.
\item En fonctionnement normal : majeure partie $\rightarrow$ mécanique, faible partie $\rightarrow$ chaleur (échauffement normal).
\item Si l'arbre est bloqué (rotation nulle) : énergie mécanique $\approx 0$ ; \textbf{presque toute l'énergie électrique devient chaleur} dans les bobinages (effet Joule $P_{Joule} = R \cdot I^2$).
\end{itemize}
\end{aretenir}

\courssub{Utilisation : plaque signalétique, calcul du courant et sécurité}

\begin{methode}[Lire une plaque signalétique et calculer le courant nominal]
\begin{enumerate}
\item Repérer la tension d'alimentation $U$ (en volts, V) et la puissance utile $P$ (en watts, W).
\item Appliquer la relation de la puissance électrique : $P = U \times I$.
\item En déduire l'intensité nominale : $I = \dfrac{P}{U}$ (en ampères, A).
\item Choisir les protections (fusible, disjoncteur, section de câble) pour un courant \textbf{supérieur ou égal} à cette valeur nominale.
\end{enumerate}
\end{methode}

\begin{exemplebox}[Application à la pompe de M. Essono]
Données : $U = \SI{220}{V}$ ; $P = \SI{750}{W}$.
\[
I = \frac{P}{U} = \frac{750}{220} \approx \SI{3,4}{A}
\]
Le moteur est conçu pour absorber environ \SI{3,4}{A} en fonctionnement normal. Un fusible \SI{4}{A} ou \SI{6}{A} (courbe adaptée au démarrage) et un câble de section \SI{1,5}{mm^2} minimum sont adaptés.
\end{exemplebox}

\begin{attention}[Sécurité électrique et mécanique]
\begin{itemize}
\item \textbf{Risque électrique :} tout appareil près de l'eau (pompe, lave-linge) doit être branché sur une prise avec \textbf{prise de terre} et protégée par un \textbf{disjoncteur différentiel $30$~mA} (norme ENEO). Ne jamais manipuler branché avec les mains mouillées.
\item \textbf{Risque thermique :} un moteur qui ronronne sans tourner (bloqué, surchargé, ventilateur obstrué) grille en quelques minutes. \textbf{Couper l'alimentation immédiatement} si odeur de brûlé, fumée, ou arrêt brutal.
\item \textbf{Entretien :} nettoyer les entrées d'air (ventilation), vérifier l'état des balais (usure), graisser les roulements si nécessaire, contrôler le serrage des cosses.
\end{itemize}
\end{attention}

\courssub{Activité d'application : le cas de la pompe de M. Essono}

\begin{enumerate}
\item \textbf{Identifier.} Sur le schéma de la figure (page précédente), nomme les quatre constituants repérés (stator, rotor, balais, collecteur) et précise le rôle de chacun dans le fonctionnement du moteur.
\item \textbf{Calculer.} À l'aide de la plaque signalétique, calcule l'intensité $I$ du courant absorbé par le moteur en fonctionnement normal. Donne le résultat arrondi au dixième d'ampère.
\item \textbf{Expliquer.} Quand le tuyau est bouché, l'arbre du moteur se bloque : la rotation devient presque nulle alors que le moteur reste alimenté. En t'appuyant sur la chaîne énergétique, explique en quelques phrases pourquoi l'énergie électrique fournie se transforme alors presque entièrement en chaleur, et pourquoi cet échauffement excessif détruit les bobinages.
\item \textbf{Proposer.} Rédige pour M. Essono une liste de \textbf{trois consignes} de sécurité et d'entretien qu'il devra respecter avec sa prochaine pompe. Chaque consigne doit être justifiée par une notion de la leçon.
\end{enumerate}

\courssub{Démarche et corrigé commenté}

\begin{exoresolu}[Pourquoi le moteur de la pompe a-t-il grillé ?]
\textbf{Énoncé.} La pompe de M. Essono porte les indications $U = \SI{220}{V}$ et $P = \SI{750}{W}$. Son tuyau de sortie s'est bouché, l'arbre du moteur s'est bloqué, et le moteur a fini par griller. 1) Identifier les constituants du moteur sur le schéma. 2) Calculer le courant normal de fonctionnement. 3) Expliquer la destruction du moteur à l'aide de la chaîne énergétique. 4) Proposer trois consignes de sécurité et d'entretien.
\tcblower
\textbf{1) Identification des constituants.}

Sur le schéma, on reconnaît les quatre organes essentiels du moteur à courant continu :
\begin{itemize}
\item le \textbf{stator} (les aimants fixes N et S, en bleu) : il crée le champ magnétique $\vec{B}$ radial dans lequel la bobine est plongée ;
\item le \textbf{rotor} (la bobine rectangulaire orange) : parcouru par le courant $I$, il subit les forces de Laplace $\vec{F}_1$ et $\vec{F}_2$ qui créent un couple moteur ;
\item les \textbf{balais} (charbons fixes) et le \textbf{collecteur} (demi-anneaux tournants) : ils assurent le passage du courant entre l'alimentation fixe et la bobine en rotation, et \textbf{inversent le sens du courant à chaque demi-tour} pour que le couple moteur conserve toujours le même sens ;
\item l'\textbf{arbre} (axe central) : il transmet le mouvement de rotation à la pompe (ou à la charge).
\end{itemize}

\textbf{2) Calcul du courant normal de fonctionnement.}

On applique la relation de la puissance électrique absorbée :
\begin{align*}
P &= U \times I \\
I &= \frac{P}{U} \\
I &= \frac{750}{220} \\
I &\approx \SI{3,4}{A}
\end{align*}
En fonctionnement normal, le moteur absorbe un courant d'environ $\SI{3,4}{A}$. C'est la valeur pour laquelle les bobinages (section de fil, vernis isolant) et les protections (fusible, disjoncteur) ont été dimensionnés.

\textbf{3) Explication de la panne par la chaîne énergétique.}

En fonctionnement normal (eau qui coule), la chaîne énergétique s'écrit :
\[ \text{Énergie électrique} \;\rightarrow\; \text{Énergie mécanique (rotation pompe)} \;+\; \text{Énergie thermique (faible échauffement normal)}. \]

Quand le tuyau se bouche, l'eau ne s'écoule plus $\rightarrow$ la pompe exerce un couple résistant énorme $\rightarrow$ l'arbre du moteur se bloque (vitesse $\approx 0$). La production d'énergie mécanique devient \textbf{quasiment nulle}. Pourtant, le moteur reste branché sur le secteur : il continue d'absorber de l'énergie électrique (on entend le \og ronronnement \fg{} du courant dans les bobines).

Or l'énergie ne disparaît pas (principe de conservation) : \textbf{presque toute l'énergie électrique absorbée se transforme désormais en énergie thermique} par effet Joule dans la résistance des bobinages ($P_{Joule} = R \cdot I^2$). Comme il n'y a plus de rotation, il n'y a plus de force contre-électromotrice pour limiter le courant ; le courant peut même dépasser la valeur nominale, aggravant l'échauffement.

La température des fils de cuivre augmente alors très fortement (bien au-delà de \SI{150}{\celsius}), le vernis isolant qui recouvre les spires fond, les spires nues se touchent $\rightarrow$ \textbf{court-circuit interne} $\rightarrow$ emballement thermique $\rightarrow$ odeur de brûlé, puis arrêt définitif. Le moteur est \og grillé \fg{} : destruction irréversible des bobinages.

\textbf{4) Trois consignes pour la prochaine pompe.}

\begin{enumerate}
\item \textbf{Surveiller le fonctionnement et arrêter immédiatement la pompe dès que l'eau ne sort plus, que le bruit change (ronronnement anormal) ou qu'une odeur de brûlé apparaît.} \\
Justification : un moteur bloqué transforme toute l'énergie électrique en chaleur (chaîne énergétique) ; seul le débranchement rapide (coupure de l'énergie électrique) peut sauver les bobinages avant fusion de l'isolant.
\item \textbf{Nettoyer régulièrement la crépine d'aspiration (filtre) et vérifier que le tuyau de sortie n'est pas bouché (plié, obstrué par des débris) avant chaque mise en marche.} \\
Justification : c'est l'obstruction du circuit hydraulique qui a provoqué le blocage mécanique de l'arbre (couple résistant $>$ couple moteur) ; un entretien préventif supprime la cause racine de la surchauffe.
\item \textbf{Brancher la pompe sur une installation conforme : prise avec terre en bon état, protégée par un disjoncteur différentiel $30$~mA et un disjoncteur divisionnaire adapté (ex. \SI{10}{A} ou \SI{16}{A} courbe C pour absorber le pic de démarrage), sans rallonge dénudée, mouillée ou sous-dimensionnée.} \\
Justification : la pompe fonctionne au voisinage de l'eau (risque d'électrisation) ; le disjoncteur différentiel protège la personne, le disjoncteur divisionnaire coupe en cas de surintensité prolongée (moteur bloqué) avant que les câbles ne chauffent, et une rallonge adaptée évite la chute de tension et l'échauffement des fils.
\end{enumerate}
\end{exoresolu}

\courssub{Bilan de la leçon}

Cette leçon t'a permis de mobiliser ensemble tous les savoirs du module \og Technologie et mécanique \fg{} dans une situation réelle camerounaise :
\begin{itemize}
\item tu as \textbf{identifié} sur un schéma les constituants du moteur à courant continu (stator, rotor, balais, collecteur) et relié chacun à son rôle précis (champ magnétique, forces de Laplace, inversion du courant, transmission) ;
\item tu as \textbf{exploité} une plaque signalétique réelle ($U = \SI{220}{V}$, $P = \SI{750}{W}$) et calculé le courant nominal $I = \dfrac{P}{U} \approx \SI{3,4}{A}$, compétence indispensable pour choisir fusible, disjoncteur et section de câble ;
\item tu as \textbf{expliqué} une panne classique (moteur bloqué) grâce à la chaîne énergétique : suppression de la sortie mécanique $\rightarrow$ conversion totale de l'énergie électrique en chaleur $\rightarrow$ destruction thermique des isolants ;
\item tu as \textbf{formulé} des consignes de sécurité (différentiel $30$~mA, terre, arrêt d'urgence) et d'entretien (filtres, vérification tuyaux) directement applicables à la maison, au champ ou à l'atelier.
\end{itemize}

\begin{attention}[L'essentiel à retenir pour le BEPC]
\begin{itemize}
\item \textbf{Constitution :} stator (aimant fixe), rotor (bobine mobile), balais + collecteur (contact tournant + inversion courant).
\item \textbf{Principe :} forces de Laplace sur les côtés actifs de la bobine $\rightarrow$ couple moteur. Collecteur = inverseur de courant pour rotation continue.
\item \textbf{Formule clé :} $P = U \times I$ $\Rightarrow$ $I = \dfrac{P}{U}$ (exigible : lecture plaque + calcul).
\item \textbf{Chaîne énergétique :} $E_{elec} \rightarrow E_{mec} + E_{therm}$. Arbre bloqué $\Rightarrow E_{mec} \approx 0 \Rightarrow E_{elec} \approx E_{therm}$ $\Rightarrow$ surchauffe $\Rightarrow$ destruction.
\item \textbf{Sécurité :} différentiel $30$~mA + terre (proximité eau), disjoncteur calibré, arrêt immédiat si anomalie (bruit, odeur, fumée).
\end{itemize}
\end{attention}

\begin{savaistu}[Le moteur dans le quotidien camerounais]
\begin{itemize}
\item \textbf{Moteurs à courant continu (CC) :} démarreurs de moto-taxi (bendix), essuie-glaces, vitres électriques, perceuses portatives (batterie), petits jouets. Simples à commander (variation tension = variation vitesse).
\item \textbf{Moteurs à courant alternatif (CA) :} la grande majorité des pompes ENEO (\SI{220}{V} monophasé ou \SI{380}{V} triphasé), ventilateurs de plafond, compresseurs de frigo, groupes électrogènes (alternateur). Le principe de base (champ tournant + rotor) est voisin, mais pas de balais/collecteur pour les asynchrones (plus robustes, sans entretien).
\item \textbf{SONARA \& Énergie :} l'essence/gasoil alimente les groupes électrogènes (moteur thermique $\rightarrow$ alternateur $\rightarrow$ courant alternatif) et les moto-pompes thermiques (pas de câble ENEO, autonomie totale). Le rendement global chaîne (thermique $\rightarrow$ mécanique $\rightarrow$ électrique $\rightarrow$ mécanique pompe) est plus faible qu'une pompe directe sur réseau ENEO, mais indispensable en zone non électrifiée.
\item \textbf{Déchets :} un moteur grillé = cuivre (recyclable, valeur marchande à la ferraille) + fer + isolants (polluants). Ne pas jeter dans la nature : apporter chez le ferrailleur ou en déchetterie.
\end{itemize}
\end{savaistu}

\newpage

\coursec{Méthodes et exercices résolus}

\courssub{Méthode 1 : Identifier les organes d'un moteur électrique à courant continu}

\begin{methode}[Identifier les organes d'un moteur à courant continu]
Pour reconnaître et nommer les organes d'un moteur électrique à courant continu sur un schéma ou sur un objet réel (moteur de jouet, ventilateur démonté) :
\begin{enumerate}
\item Repérer d'abord les deux parties principales : la partie \cle{fixe} (le \cle{stator}) et la partie \cle{tournante} (le \cle{rotor}).
\item Sur le stator, identifier l'\cle{aimant} (ou l'électroaimant inducteur) qui crée le champ magnétique.
\item Sur le rotor, identifier la \cle{bobine} (enroulement de fil de cuivre) montée sur l'\cle{arbre} de rotation.
\item Repérer le \cle{collecteur} : cylindre fendu en deux demi-anneaux métalliques, solidaire du rotor.
\item Repérer les \cle{balais} (ou charbons) : petits blocs de graphite frottant sur le collecteur et reliés à l'alimentation.
\item Conclure en citant le rôle de chaque organe : le couple collecteur-balais assure le passage du courant dans la bobine et inverse le sens du courant à chaque demi-tour.
\end{enumerate}
\textbf{Réflexe :} mémoriser la chaîne : alimentation $\rightarrow$ balais $\rightarrow$ collecteur $\rightarrow$ bobine du rotor.
\end{methode}

\begin{exoresolu}[Les organes du moteur d'un ventilateur]
En démontant le moteur à courant continu d'un petit ventilateur, un élève de 3ème observe les éléments suivants : un aimant fixé à la carcasse, un enroulement de fil de cuivre monté sur un axe, un cylindre métallique fendu en deux parties solidaire de l'axe, et deux petits blocs de graphite reliés aux fils d'alimentation.
\begin{enumerate}
\item Nommer chaque élément observé.
\item Préciser quels éléments appartiennent au stator et lesquels appartiennent au rotor.
\item Quel est le rôle du couple collecteur-balais ?
\end{enumerate}
\tcblower
\textbf{1. Nom de chaque élément.}
\begin{itemize}
\item L'aimant fixé à la carcasse est l'\cle{inducteur} (aimant du stator) : il crée le champ magnétique.
\item L'enroulement de fil de cuivre sur l'axe est la \cle{bobine} (l'induit).
\item Le cylindre métallique fendu en deux parties est le \cle{collecteur}.
\item Les deux blocs de graphite sont les \cle{balais} (ou charbons).
\end{itemize}
\textbf{2. Stator et rotor.}
\begin{itemize}
\item Le \cle{stator} (partie fixe) comprend l'aimant inducteur et les balais (fixés au capot).
\item Le \cle{rotor} (partie tournante) comprend la bobine, l'arbre de rotation et le collecteur.
\end{itemize}
\textbf{3. Rôle du couple collecteur-balais.} Les balais frottent sur le collecteur et conduisent le courant de l'alimentation vers la bobine. Grâce à la fente du collecteur, le sens du courant dans la bobine s'inverse à chaque demi-tour : les forces magnétiques changent alors de sens au bon moment, ce qui maintient le rotor en rotation continue dans le même sens.
\textbf{Conclusion :} le moteur du ventilateur est un moteur à courant continu constitué d'un stator (aimant + balais) et d'un rotor (bobine + collecteur + arbre) ; le couple collecteur-balais assure l'alimentation de la bobine et l'inversion du courant à chaque demi-tour.
\end{exoresolu}

\courssub{Méthode 2 : Interpréter l'expérience de la spire mobile dans un champ magnétique}

\begin{methode}[Interpréter l'action d'un aimant sur une bobine parcourue par un courant]
Pour interpréter une expérience où une spire (ou une bobine) mobile est placée entre les pôles d'un aimant en U :
\begin{enumerate}
\item Vérifier les conditions : la spire est \cle{mobile} autour d'un axe et placée dans le \cle{champ magnétique} de l'aimant.
\item Observer ce qui se passe \emph{quand on ferme l'interrupteur} : si la spire tourne ou bascule, c'est qu'elle subit des \cle{forces électromagnétiques}.
\item Faire les expériences témoins : couper le courant (plus de mouvement), inverser le sens du courant (le sens de rotation s'inverse), retourner l'aimant (le sens de rotation s'inverse aussi).
\item Conclure : un conducteur (ou une bobine) parcouru par un courant et placé dans un champ magnétique subit des \cle{forces électromagnétiques} ; leur sens dépend du sens du courant et du sens du champ magnétique.
\item Faire le lien avec le moteur : c'est cette action qui fait tourner le rotor.
\end{enumerate}
\textbf{Réflexe :} pas de courant $\Rightarrow$ pas de force ; on inverse le courant ou l'aimant $\Rightarrow$ on inverse le mouvement.
\end{methode}

\begin{exoresolu}[La spire qui bascule]
Au laboratoire, on place une spire rectangulaire de fil de cuivre, mobile autour d'un axe vertical, entre les branches d'un aimant en U. La spire est reliée à un générateur par l'intermédiaire d'un interrupteur.
\begin{enumerate}
\item On ferme l'interrupteur : que observe-t-on ? Justifier.
\item On inverse les bornes du générateur : que se passe-t-il ?
\item On ouvre l'interrupteur : que devient la spire ? Conclure.
\end{enumerate}
\tcblower
\textbf{1. Fermeture de l'interrupteur.} À la fermeture, un courant circule dans la spire. On observe que la spire \cle{bascule} (tourne) autour de son axe. Justification : la spire est parcourue par un courant électrique et placée dans le champ magnétique de l'aimant en U ; elle subit donc des \cle{forces électromagnétiques} qui créent un couple de rotation.
\textbf{2. Inversion des bornes du générateur.} Le sens du courant dans la spire est inversé. On observe que la spire bascule \cle{dans le sens contraire}. En effet, le sens des forces électromagnétiques dépend du sens du courant : inverser le courant inverse les forces.
\textbf{3. Ouverture de l'interrupteur.} Le courant ne circule plus : la spire ne subit plus aucune force électromagnétique et s'immobilise (elle revient à sa position d'équilibre).
\textbf{Conclusion :} une bobine parcourue par un courant et placée dans un champ magnétique subit des forces électromagnétiques dont le sens dépend du sens du courant et du sens du champ magnétique. C'est le \cle{principe de fonctionnement du moteur électrique}, qui convertit l'énergie électrique en énergie mécanique (de rotation).
\end{exoresolu}

\courssub{Méthode 3 : Expliquer le fonctionnement d'un moteur électrique}

\begin{methode}[Rédiger l'explication du fonctionnement d'un moteur]
Pour expliquer clairement, en quelques étapes, comment un moteur à courant continu tourne :
\begin{enumerate}
\item Citer le \cle{phénomène physique} mis en jeu : l'action d'un champ magnétique sur un conducteur parcouru par un courant (forces électromagnétiques).
\item Décrire le début du mouvement : le courant arrive dans la bobine du rotor par les balais et le collecteur ; la bobine, placée dans le champ de l'aimant, subit des forces qui la font tourner.
\item Expliquer la \cle{rotation continue} : à chaque demi-tour, le collecteur fendu inverse le sens du courant dans la bobine ; les forces s'inversent au bon moment et entretiennent la rotation dans le même sens.
\item Conclure par la \cle{conversion d'énergie} : le moteur transforme l'énergie électrique en énergie mécanique (rotation de l'arbre).
\end{enumerate}
\textbf{Réflexe de rédaction :} toujours terminer par la phrase bilan énergétique : « énergie électrique $\rightarrow$ énergie mécanique ».
\end{methode}

\begin{exoresolu}[Pourquoi le moteur tourne-t-il sans s'arrêter ?]
Un élève affirme : « La bobine du rotor tourne d'un quart de tour puis devrait s'arrêter, car les forces s'annulent. Pourtant le moteur tourne en continu. » Expliquer, en quatre étapes rédigées, pourquoi la rotation est continue, en précisant le rôle du collecteur et la conversion d'énergie réalisée.
\tcblower
\textbf{Étape 1 : le phénomène mis en jeu.} Tout conducteur parcouru par un courant et placé dans un champ magnétique subit des \cle{forces électromagnétiques}. C'est le phénomène à la base du moteur électrique.
\textbf{Étape 2 : la mise en rotation.} Lorsqu'on alimente le moteur, le courant circule de l'alimentation vers la bobine du rotor en passant par les balais et le collecteur. La bobine, baignant dans le champ magnétique de l'aimant du stator, subit des forces qui forment un couple : elle commence à tourner.
\textbf{Étape 3 : la rotation continue grâce au collecteur.} Sans dispositif particulier, la bobine s'arrêterait en position d'équilibre. Mais le \cle{collecteur}, fendu en deux demi-anneaux, inverse automatiquement le sens du courant dans la bobine à chaque demi-tour. Les forces électromagnétiques changent alors de sens au bon moment : le couple reste toujours moteur et la rotation se poursuit dans le même sens, tant que le courant circule.
\textbf{Étape 4 : bilan énergétique.} Le moteur électrique est un convertisseur : il transforme l'\cle{énergie électrique} fournie par l'alimentation en \cle{énergie mécanique} de rotation disponible sur l'arbre (avec une petite partie perdue en chaleur par effet Joule).
\textbf{Conclusion :} l'affirmation de l'élève serait exacte sans le collecteur ; c'est l'inversion du courant à chaque demi-tour par le couple collecteur-balais qui garantit la rotation continue du moteur.
\end{exoresolu}

\courssub{Méthode 4 : Exploiter la plaque signalétique d'un moteur}

\begin{methode}[Lire et utiliser une plaque signalétique]
La \cle{plaque signalétique} est l'étiquette rivetée sur la carcasse du moteur. Pour l'exploiter :
\begin{enumerate}
\item Relever la \cle{tension nominale} $U$ (en volts) : c'est la tension d'alimentation à respecter.
\item Relever le \cle{courant nominal} $I$ (en ampères) ou la \cle{puissance} $P$ (en watts).
\item Si l'une des trois grandeurs manque, utiliser la relation de la puissance électrique :
\[ P = U \times I \qquad \text{avec } P \text{ en W, } U \text{ en V, } I \text{ en A.} \]
\item En déduire les grandeurs cherchées : $I = \dfrac{P}{U}$ ou $U = \dfrac{P}{I}$.
\item Vérifier la cohérence : un moteur prévu pour \SI{12}{V} ne doit jamais être branché sur le secteur \SI{220}{V}.
\end{enumerate}
\textbf{Réflexe :} avant tout calcul, convertir si nécessaire (mA en A : diviser par 1000).
\end{methode}

\begin{exoresolu}[La plaque signalétique du moteur d'une pompe]
Sur la plaque signalétique du moteur électrique d'une pompe à eau, on lit : « $U = \SI{220}{V}$ ; $P = \SI{550}{W}$ ».
\begin{enumerate}
\item Que signifient ces deux indications ?
\item Calculer l'intensité nominale du courant qui traverse ce moteur en fonctionnement normal.
\item Le fusible de protection du circuit est calibré à \SI{2}{A}. Convient-il ? Justifier.
\end{enumerate}
\tcblower
\textbf{1. Signification des indications.} $U = \SI{220}{V}$ est la \cle{tension nominale} : le moteur doit être alimenté sous \SI{220}{V} (tension du réseau ENEO). $P = \SI{550}{W}$ est la \cle{puissance nominale} consommée par le moteur en fonctionnement normal.
\textbf{2. Calcul de l'intensité nominale.} On utilise la relation $P = U \times I$, d'où :
\begin{align*}
I &= \frac{P}{U} \\
I &= \frac{550}{220} \\
I &= \SI{2,5}{A}
\end{align*}
L'intensité nominale du courant est $I = \SI{2,5}{A}$.
\textbf{3. Choix du fusible.} Le fusible de \SI{2}{A} est inférieur à l'intensité nominale du moteur ($\SI{2}{A} < \SI{2,5}{A}$). Il fondra donc dès la mise en marche normale de la pompe : il \cle{ne convient pas}. Il faut un fusible de calibre légèrement supérieur à \SI{2,5}{A} (par exemple \SI{3}{A}) pour protéger le circuit sans couper intempestivement.
\textbf{Conclusion :} le moteur fonctionne sous \SI{220}{V} avec une intensité de \SI{2,5}{A} ; le fusible de \SI{2}{A} est sous-calibré et doit être remplacé.
\end{exoresolu}

\courssub{Méthode 5 : Appliquer les consignes de sécurité liées aux moteurs électriques}

\begin{methode}[Identifier les dangers et les règles de sécurité]
Pour répondre à une question de sécurité concernant un appareil à moteur (ventilateur, pompe de forage, démarreur de groupe électrogène, meuleuse) :
\begin{enumerate}
\item Identifier les \cle{dangers} : électrisation (contact avec des fils dénudés), court-circuit, échauffement du moteur, entraînement mécanique (doigts, cheveux, vêtements dans les parties tournantes).
\item Citer les règles de prévention \emph{avant} l'utilisation : vérifier que la tension de l'appareil correspond à celle du réseau, contrôler l'état des câbles et de la prise, ne jamais utiliser un appareil électrique les mains mouillées ou dans un milieu humide.
\item Citer les règles \emph{pendant} l'utilisation : ne pas approcher les mains des parties tournantes, ne pas bloquer l'arbre du moteur (surchauffe), débrancher avant toute intervention.
\item Citer les protections : fusible ou disjoncteur adapté au calibre, prise de terre, carter de protection des parties tournantes.
\item Rédiger la réponse en reliant chaque règle au danger qu'elle prévient.
\end{enumerate}
\textbf{Réflexe :} une règle de sécurité se justifie toujours par le danger évité.
\end{methode}

\begin{exoresolu}[Le ventilateur de la salle de classe]
Le ventilateur de la salle de 3ème présente un câble d'alimentation partiellement dénudé. Un élève propose de l'utiliser quand même « parce qu'il tourne encore ».
\begin{enumerate}
\item Citer deux dangers encourus.
\item Donner deux consignes de sécurité à respecter avant toute nouvelle utilisation.
\item Expliquer pourquoi il est dangereux de bloquer les pales du ventilateur pendant qu'il fonctionne.
\end{enumerate}
\tcblower
\textbf{1. Dangers encourus.}
\begin{itemize}
\item L'\cle{électrisation} (voire l'électrocution) : en touchant la partie dénudée du câble, le corps humain peut être traversé par le courant sous \SI{220}{V}.
\item Le \cle{court-circuit} : si les deux fils dénudés se touchent, une surintensité peut provoquer des étincelles, l'échauffement du câble et un début d'incendie.
\end{itemize}
\textbf{2. Consignes de sécurité.}
\begin{itemize}
\item \cle{Débrancher} immédiatement l'appareil et ne plus l'utiliser.
\item Faire \cle{réparer ou remplacer} le câble par une personne qualifiée avant toute nouvelle utilisation ; vérifier aussi que la prise est en bon état et reliée à la terre.
\end{itemize}
\textbf{3. Blocage des pales.} Si l'on bloque les pales, le rotor du moteur est empêché de tourner alors qu'il reste alimenté : le moteur n'assure plus la conversion en énergie mécanique et l'énergie électrique reçue se transforme presque entièrement en \cle{chaleur} (effet Joule). Le moteur \cle{surchauffe}, ce qui peut détruire les enroulements, dégager une odeur de brûlé et provoquer un incendie.
\textbf{Conclusion :} on ne doit jamais utiliser un appareil au câble dénudé ni bloquer l'arbre d'un moteur en marche : on risque l'électrisation, le court-circuit et la destruction du moteur par échauffement.
\end{exoresolu}

\courssub{Méthode 6 : Réaliser le schéma du circuit d'alimentation d'un moteur}

\begin{methode}[Schématiser un circuit simple avec moteur]
Pour schématiser le circuit électrique d'alimentation d'un moteur à courant continu (expérience de cours) :
\begin{enumerate}
\item Dresser la liste des dipôles : \cle{générateur} (pile), \cle{interrupteur}, \cle{moteur}, fils de connexion.
\item Dessiner le circuit en \cle{série} : une seule boucle fermée reliant la borne $+$ de la pile à la borne $-$ en passant par l'interrupteur et le moteur.
\item Respecter les symboles normalisés : long trait fin (borne $+$) et petit trait épais (borne $-$) pour la pile, cercle avec la lettre M pour le moteur.
\item Indiquer le sens du courant (de la borne $+$ vers la borne $-$ à l'extérieur du générateur) par une flèche.
\item Prévoir les variantes demandées : inversion du sens de rotation $\Rightarrow$ inverser les bornes de la pile ; arrêt du moteur $\Rightarrow$ ouvrir l'interrupteur.
\end{enumerate}
\textbf{Réflexe :} un circuit ouvert (interrupteur ouvert) $\Rightarrow$ pas de courant $\Rightarrow$ le moteur ne tourne pas.
\end{methode}

\begin{exoresolu}[Faire tourner le moteur dans les deux sens]
On dispose d'une pile de \SI{4,5}{V}, d'un interrupteur, d'un moteur à courant continu et de fils de connexion.
\begin{enumerate}
\item Dessiner le schéma du circuit permettant de faire tourner le moteur, en indiquant le sens du courant.
\item Le moteur tourne dans le sens des aiguilles d'une montre. Que faire pour qu'il tourne en sens inverse ? Justifier.
\item Que se passe-t-il si l'on remplace la pile de \SI{4,5}{V} par une pile de \SI{9}{V} alors que la plaque du moteur indique \SI{4,5}{V} ?
\end{enumerate}
\tcblower
\textbf{1. Schéma du circuit.} Le circuit est un circuit série à une seule boucle :
\begin{center}
\begin{tikzpicture}[scale=1, every node/.style={font=\small}]
% Pile : petit trait epais = borne -, long trait fin = borne +
\draw[thick] (0,0) -- (0.6,0);
\draw[line width=2pt] (0.6,-0.2) -- (0.6,0.2);
\draw[thick] (0.8,-0.45) -- (0.8,0.45);
\node[below] at (0.7,-0.5) {Pile \SI{4,5}{V}};
\node[above] at (0.6,0.25) {$-$};
\node[above] at (0.8,0.5) {$+$};
\draw[thick] (0.8,0) -- (2.2,0);
% Interrupteur (ferme)
\draw[thick] (2.2,0) -- (2.85,0);
\fill (2.2,0) circle (0.04);
\fill (2.85,0) circle (0.04);
\node[below] at (2.5,-0.05) {Interrupteur};
\draw[thick] (2.85,0) -- (4.5,0) -- (4.5,1.2);
% Moteur
\draw[thick] (4.5,1.2) -- (3.9,1.2);
\draw[thick,fill=popblueL] (3.45,1.2) circle (0.45);
\node at (3.45,1.2) {\textbf{M}};
\draw[thick] (3.0,1.2) -- (0,1.2) -- (0,0);
% Sens du courant
\draw[-{Stealth},poporange,very thick] (1.6,0) -- (2.0,0);
\node[poporange,below] at (1.8,-0.1) {$I$};
\end{tikzpicture}
\end{center}
Le courant sort de la borne $+$ de la pile, traverse l'interrupteur puis le moteur, et revient à la borne $-$.
\textbf{2. Inversion du sens de rotation.} Il suffit d'\cle{inverser les bornes de la pile} (ou les deux fils aux bornes du moteur). Justification : le sens des forces électromagnétiques qui s'exercent sur la bobine du rotor dépend du sens du courant ; en inversant le courant, on inverse le couple moteur, donc le sens de rotation.
\textbf{3. Branchement sous \SI{9}{V}.} La tension appliquée ($\SI{9}{V}$) est le double de la tension nominale ($\SI{4,5}{V}$) : le moteur est en \cle{surtension}. L'intensité devient trop grande, le moteur tourne trop vite, chauffe fortement par effet Joule et risque d'être \cle{détruit} (enroulements brûlés).
\textbf{Conclusion :} le circuit série pile-interrupteur-moteur permet de commander la rotation ; on inverse le sens de rotation en inversant le sens du courant, et l'on doit toujours respecter la tension nominale indiquée sur la plaque signalétique.
\end{exoresolu}

\begin{attention}[Les 5 erreurs qui coûtent des points au BEPC]
\begin{enumerate}
\item \textbf{Confondre stator et rotor.} Le stator est la partie \emph{fixe} (aimant, balais), le rotor la partie \emph{tournante} (bobine, collecteur, arbre). Inverser les deux dans une définition fait perdre tous les points de la question.
\item \textbf{Oublier le rôle du collecteur.} Beaucoup d'élèves écrivent « la bobine tourne parce qu'elle subit des forces » sans préciser que le collecteur \emph{inverse le sens du courant à chaque demi-tour}. Sans cette phrase, l'explication de la rotation continue est incomplète.
\item \textbf{Oublier la conversion d'énergie.} La question « quel est le principe du moteur ? » exige la conclusion : transformation d'\emph{énergie électrique} en \emph{énergie mécanique}. Ne jamais écrire « le moteur produit de l'électricité » (c'est le rôle du générateur !).
\item \textbf{Se tromper dans la formule de puissance.} Écrire $P = \dfrac{U}{I}$ ou $P = U + I$ au lieu de $P = U \times I$. Vérifier la cohérence : $I = \dfrac{P}{U}$ donne des ampères si $P$ est en watts et $U$ en volts. Ne pas oublier de convertir les mA en A.
\item \textbf{Négliger les unités et la justification en sécurité.} Donner un résultat sans unité (2,5 au lieu de \SI{2,5}{A}) coûte des points. De même, une consigne de sécurité non justifiée (« il faut débrancher » sans dire pourquoi) est à moitié notée : toujours relier la règle au danger (électrisation, court-circuit, surchauffe).
\end{enumerate}
\end{attention}

\newpage

\coursec{Exercices}

\courssub{Je vérifie mes connaissances}

\begin{exercice}[QCM : le moteur électrique]
Pour chaque question, recopie la (ou les) bonne(s) réponse(s).
\begin{enumerate}
\item Un moteur électrique convertit principalement :
\begin{itemize}
\item[a)] l'énergie mécanique en énergie électrique ;
\item[b)] l'énergie électrique en énergie mécanique ;
\item[c)] l'énergie chimique en énergie électrique.
\end{itemize}
\item Dans un moteur à courant continu, la partie fixe s'appelle :
\begin{itemize}
\item[a)] le rotor ; \item[b)] le stator ; \item[c)] le collecteur.
\end{itemize}
\item Le rôle du collecteur est de :
\begin{itemize}
\item[a)] fixer le moteur sur son support ;
\item[b)] inverser le sens du courant dans la bobine à chaque demi-tour pour maintenir la rotation ;
\item[c)] refroidir le moteur.
\end{itemize}
\item Les balais (charbons) assurent :
\begin{itemize}
\item[a)] le contact électrique entre la partie fixe et le collecteur en rotation ;
\item[b)] l'étanchéité du moteur ;
\item[c)] la fixation de l'aimant.
\end{itemize}
\end{enumerate}
\end{exercice}

\begin{exercice}[Vrai ou faux]
Recopie chaque affirmation, indique si elle est vraie ou fausse, puis justifie en une phrase.
\begin{enumerate}
\item Une spire parcourue par un courant, placée entre les pôles d'un aimant, ne subit aucune force.
\item Si on inverse le sens du courant dans la bobine, le sens de rotation du moteur s'inverse.
\item Le rotor est la partie mobile du moteur électrique.
\item Un moteur électrique ne chauffe jamais, même après plusieurs heures de fonctionnement.
\end{enumerate}
\end{exercice}

\begin{exercice}[Définitions]
Définis en une ou deux phrases chacun des termes suivants : \cle{moteur électrique}, \cle{stator}, \cle{rotor}, \cle{collecteur}, \cle{balai}.
\end{exercice}

\begin{exercice}[Calcul direct]
Sur la plaque signalétique d'un moteur électrique, on lit : $220~\text{V}$ -- $440~\text{W}$.
\begin{enumerate}
\item Rappelle la relation liant la puissance $P$, la tension $U$ et l'intensité $I$.
\item Calcule l'intensité du courant qui traverse ce moteur en fonctionnement normal.
\end{enumerate}
\end{exercice}

\courssub{Je m'entraîne}

\begin{exercice}[L'expérience de la spire mobile]
Au laboratoire, on place une spire rectangulaire, parcourue par un courant continu, entre les pôles d'un aimant en U. La spire est mobile autour d'un axe vertical.
\begin{enumerate}
\item Décris ce que l'on observe lorsqu'on ferme l'interrupteur.
\item Que se passe-t-il si l'on inverse le sens du courant dans la spire ? Justifie ta réponse.
\item Que se passe-t-il si l'on inverse les pôles de l'aimant ?
\item Conclus : quel est le phénomène physique mis en jeu dans le moteur électrique ?
\end{enumerate}
\end{exercice}

\begin{exercice}[Le ventilateur de la maison]
Pendant la saison sèche à Douala, la famille de Ngo Bell utilise un ventilateur électrique plusieurs heures par jour.
\begin{enumerate}
\item Établis la chaîne énergétique du ventilateur en marche (énergie reçue $\rightarrow$ énergie utile $\rightarrow$ énergie perdue).
\item Après trois heures de fonctionnement, le moteur du ventilateur est chaud. Explique pourquoi.
\item Cite deux précautions à prendre pour utiliser ce ventilateur en toute sécurité.
\end{enumerate}
\end{exercice}

\begin{exercice}[Plaque signalétique d'une pompe à eau]
Une pompe à eau utilisée dans un quartier de Yaoundé porte les indications suivantes : $220~\text{V}$ -- $1\,100~\text{W}$.
\begin{enumerate}
\item Que signifie chacune de ces deux indications ?
\item Calcule l'intensité du courant absorbé par la pompe en fonctionnement normal.
\item L'installation est protégée par un fusible de $6~\text{A}$. Ce fusible convient-il ? Justifie.
\end{enumerate}
\end{exercice}

\begin{exercice}[Sécurité : les bons gestes]
Voici une liste de comportements observés dans un atelier de réparation de motos à Bafoussam. Pour chacun, indique s'il est correct ou dangereux, et justifie.
\begin{enumerate}
\item Débrancher le moteur avant toute intervention mécanique.
\item Bloquer l'arbre du moteur avec la main pendant qu'il tourne.
\item Brancher un petit moteur prévu pour $12~\text{V}$ directement sur le secteur $220~\text{V}$.
\item Remettre en place le carter de protection avant la mise en marche.
\item Travailler les mains mouillées sur le moteur branché.
\end{enumerate}
\end{exercice}

\courssub{Je me prépare au Bac}

\begin{exercice}[Type BEPC : étude d'un moteur à courant continu]
Au cours d'une séance de technologie, l'enseignant présente aux élèves de 3ème le schéma d'un moteur électrique à courant continu. On y distingue cinq éléments repérés par les numéros 1 à 5 : la partie fixe portant un aimant, la partie mobile bobinée, le dispositif à lamelles solidaire de l'axe, les contacts frottants en graphite et l'aimant.
\begin{enumerate}
\item Nomme chacun des cinq éléments repérés sur le schéma.
\item Donne le rôle de chacun de ces éléments dans le fonctionnement du moteur.
\item Explique, en quelques lignes, pourquoi le collecteur est indispensable au maintien de la rotation continue du rotor.
\item Le moteur étudié porte l'indication $12~\text{V}$ -- $60~\text{W}$. Calcule l'intensité du courant qui le traverse en fonctionnement normal.
\end{enumerate}
\end{exercice}

\begin{exercice}[Type BEPC : le groupe électrogène et le moteur du moulin]
Dans un village de la région de l'Adamaoua, une coopérative utilise un moulin à maïs entraîné par un moteur électrique alimenté par un groupe électrogène. Le moteur porte les indications : $220~\text{V}$ -- $2\,200~\text{W}$.
\begin{enumerate}
\item Quelle conversion d'énergie le moteur électrique réalise-t-il ?
\item Établis la chaîne énergétique complète de l'installation, depuis le carburant du groupe électrogène jusqu'à la mouture du maïs.
\item Calcule l'intensité du courant absorbé par le moteur en fonctionnement normal.
\item Le moteur fonctionne $4~\text{h}$ par jour. Calcule l'énergie électrique consommée par jour, en kilowattheures (\si{kWh}).
\item Après une longue journée de travail, le moteur est très chaud. Explique ce phénomène et propose une mesure pour éviter sa détérioration.
\end{enumerate}
\end{exercice}

\begin{exercice}[Type BEPC : investigation sur le sens de rotation]
Lors d'un TP, les élèves disposent d'un petit moteur à courant continu, d'un générateur de $6~\text{V}$, de fils de connexion et d'un interrupteur. Ils réalisent trois essais : essai 1, le moteur est branché normalement ; essai 2, on inverse les deux bornes du moteur ; essai 3, on bloque l'axe du moteur pendant la mise sous tension.
\begin{enumerate}
\item Schématise le circuit électrique réalisé (générateur, interrupteur, moteur).
\item Précise ce que l'on observe dans l'essai 1 et dans l'essai 2. Quelle loi qualitative ces observations mettent-elles en évidence ?
\item Dans l'essai 3, le moteur chauffe dangereusement. Explique pourquoi le blocage de l'arbre est dangereux pour le moteur.
\item Rédige une conclusion générale sur le principe de fonctionnement du moteur électrique à courant continu, en citant les grandeurs qui déterminent le sens de rotation.
\item Cite deux appareils de la vie courante camerounaise qui utilisent un moteur électrique, en précisant pour chacun l'énergie utile produite.
\end{enumerate}
\end{exercice}

\newpage

\coursec{Corrigés des exercices}

\courssub{Je vérifie mes connaissances}

\begin{corrige}[Exercice 1 — QCM : le moteur électrique]
\begin{enumerate}
\item \textbf{Réponse b).} Un moteur électrique convertit principalement l'énergie électrique en énergie mécanique (de rotation). C'est l'inverse d'un alternateur ou d'une dynamo, qui convertissent l'énergie mécanique en énergie électrique.
\item \textbf{Réponse b).} La partie fixe d'un moteur à courant continu est le \cle{stator} (inducteur, souvent un aimant). Le rotor est la partie mobile ; le collecteur n'est qu'un organe de liaison électrique.
\item \textbf{Réponse b).} Le collecteur, grâce à ses lamelles, inverse le sens du courant dans la bobine à chaque demi-tour : les forces exercées sur la bobine changent alors de sens au bon moment, ce qui maintient la rotation toujours dans le même sens.
\item \textbf{Réponse a).} Les balais (ou charbons), en graphite, frottent sur le collecteur en rotation et assurent le passage du courant entre la partie fixe (alimentation) et la partie mobile (bobines du rotor).
\end{enumerate}
\end{corrige}

\begin{corrige}[Exercice 2 — Vrai ou faux]
\begin{enumerate}
\item \textbf{Fausse.} Une spire parcourue par un courant et placée dans le champ magnétique d'un aimant subit des forces électromagnétiques qui créent un couple : la spire tourne (c'est le principe du moteur).
\item \textbf{Vraie.} Le sens des forces électromagnétiques dépend du sens du courant : inverser le courant inverse le sens de rotation.
\item \textbf{Vraie.} Par définition, le rotor est la partie mobile (tournante) du moteur ; le stator est la partie fixe.
\item \textbf{Fausse.} Une partie de l'énergie électrique reçue est toujours convertie en chaleur (effet Joule dans les bobinages, frottements) : tout moteur réel chauffe en fonctionnement.
\end{enumerate}
\end{corrige}

\begin{corrige}[Exercice 3 — Définitions]
\begin{itemize}
\item \cle{Moteur électrique} : appareil qui convertit l'énergie électrique en énergie mécanique de rotation, grâce aux forces électromagnétiques exercées sur des conducteurs parcourus par un courant et placés dans un champ magnétique.
\item \cle{Stator} : partie fixe du moteur électrique ; il porte l'aimant (ou l'électroaimant inducteur) qui crée le champ magnétique.
\item \cle{Rotor} : partie mobile du moteur électrique ; il porte les bobines parcourues par le courant et tourne autour d'un axe.
\item \cle{Collecteur} : cylindre formé de lamelles conductrices isolées entre elles, solidaire de l'axe du rotor ; il inverse le sens du courant dans les bobines à chaque demi-tour pour maintenir la rotation.
\item \cle{Balai} (ou charbon) : pièce en graphite fixée au stator, qui frotte sur le collecteur et assure le contact électrique entre l'alimentation fixe et les bobines en rotation.
\end{itemize}
\end{corrige}

\begin{corrige}[Exercice 4 — Calcul direct]
\begin{enumerate}
\item La relation est : $P = U \times I$, avec $P$ en watts (\si{W}), $U$ en volts (\si{V}) et $I$ en ampères (\si{A}).
\item On en déduit $I = \dfrac{P}{U}$. Application numérique :
\[ I = \frac{440}{220} = 2~\text{A}. \]
L'intensité du courant qui traverse le moteur en fonctionnement normal est $\boxed{I = 2~\text{A}}$.
\end{enumerate}
\end{corrige}

\courssub{Je m'entraîne}

\begin{corrige}[Exercice 5 — L'expérience de la spire mobile]
\begin{enumerate}
\item Lorsqu'on ferme l'interrupteur, la spire, parcourue par le courant, se met à \textbf{tourner} autour de son axe vertical : elle subit un couple de forces électromagnétiques.
\item Si l'on inverse le sens du courant, la spire tourne \textbf{en sens inverse}. En effet, le sens des forces électromagnétiques dépend du sens du courant : inverser le courant inverse le sens des forces, donc le sens de rotation.
\item Si l'on inverse les pôles de l'aimant (donc le sens du champ magnétique), la spire tourne également \textbf{en sens inverse} : le sens des forces dépend aussi du sens du champ magnétique.
\item \textbf{Conclusion :} le phénomène mis en jeu est l'\cle{action électromagnétique} : un conducteur parcouru par un courant et placé dans un champ magnétique subit une force. Le moteur électrique exploite ces forces pour produire une rotation continue.
\end{enumerate}
\end{corrige}

\begin{corrige}[Exercice 6 — Le ventilateur de la maison]
\begin{enumerate}
\item Chaîne énergétique du ventilateur en marche :
\[ \text{Énergie électrique} \;\longrightarrow\; \text{Énergie mécanique (rotation des pales)} \;+\; \text{Énergie thermique (perdue)}. \]
L'énergie reçue est l'énergie électrique du réseau ; l'énergie utile est l'énergie mécanique qui met l'air en mouvement ; l'énergie perdue est la chaleur dégagée par le moteur.
\item Après trois heures de fonctionnement, le moteur est chaud car une partie de l'énergie électrique est convertie en chaleur : par \textbf{effet Joule} dans les bobinages traversés par le courant, et par les \textbf{frottements} des pièces en rotation (balais, paliers).
\item Deux précautions (au choix) :
\begin{itemize}
\item ne pas introduire les doigts ou des objets dans la grille de protection ;
\item ne pas utiliser le ventilateur les mains mouillées ou près de l'eau ;
\item débrancher l'appareil avant tout nettoyage ;
\item ne pas couvrir le moteur (laisser l'air circuler pour le refroidissement) ;
\item vérifier que la tension du réseau correspond à celle de la plaque signalétique.
\end{itemize}
\end{enumerate}
\end{corrige}

\begin{corrige}[Exercice 7 — Plaque signalétique d'une pompe à eau]
\begin{enumerate}
\item $220~\text{V}$ : tension nominale d'alimentation, c'est-à-dire la tension sous laquelle la pompe doit être branchée pour fonctionner normalement. $1\,100~\text{W}$ : puissance électrique nominale absorbée par la pompe en fonctionnement normal.
\item $I = \dfrac{P}{U} = \dfrac{1\,100}{220} = 5~\text{A}$. La pompe absorbe un courant d'intensité $\boxed{I = 5~\text{A}}$.
\item Le fusible de $6~\text{A}$ convient : il supporte un courant maximal de $6~\text{A}$, supérieur au courant normal de $5~\text{A}$, donc il ne fond pas en fonctionnement normal ; en revanche, il fondra en cas de surintensité (court-circuit, blocage), ce qui protège la pompe et l'installation. Un fusible de calibre inférieur à $5~\text{A}$ (par exemple $4~\text{A}$) fondrait dès la mise en marche : il ne conviendrait pas.
\end{enumerate}
\end{corrige}

\begin{corrige}[Exercice 8 — Sécurité : les bons gestes]
\begin{enumerate}
\item \textbf{Correct.} Débrancher le moteur avant toute intervention supprime tout risque de mise en marche intempestive et d'électrisation.
\item \textbf{Dangereux.} Bloquer l'arbre en rotation avec la main expose à des blessures graves (écrasement, coupure) ; de plus, un moteur bloqué sous tension chauffe fortement et peut être détruit.
\item \textbf{Dangereux.} Brancher un moteur $12~\text{V}$ sur le secteur $220~\text{V}$ lui applique une tension très supérieure à sa tension nominale : le courant devient excessif, le moteur chauffe, grille et peut provoquer un incendie.
\item \textbf{Correct.} Le carter de protection empêche le contact avec les pièces en rotation (arbre, ventilateur) : il doit toujours être remis en place avant la mise en marche.
\item \textbf{Dangereux.} L'eau rend le corps conducteur : travailler les mains mouillées sur un appareil branché expose à l'électrisation, voire à l'électrocution.
\end{enumerate}
\end{corrige}

\courssub{Je me prépare au Bac}

\begin{corrige}[Exercice 9 — Type BEPC : étude d'un moteur à courant continu]
\begin{enumerate}
\item Identification des éléments : 1 — stator (partie fixe portant l'aimant) ; 2 — rotor (partie mobile bobinée) ; 3 — collecteur (dispositif à lamelles solidaire de l'axe) ; 4 — balais ou charbons (contacts frottants en graphite) ; 5 — aimant (inducteur).
\item Rôles :
\begin{itemize}
\item \textbf{Stator} : supporte l'aimant et fixe le moteur ; il crée, avec l'aimant, le champ magnétique.
\item \textbf{Rotor} : porte les bobines parcourues par le courant ; soumis aux forces électromagnétiques, il tourne et transmet le mouvement par son axe.
\item \textbf{Collecteur} : inverse le sens du courant dans les bobines à chaque demi-tour.
\item \textbf{Balais} : assurent le contact électrique entre l'alimentation fixe et le collecteur en rotation.
\item \textbf{Aimant} : produit le champ magnétique dans lequel baignent les bobines du rotor.
\end{itemize}
\item Sans collecteur, la bobine ne ferait qu'un quart de tour puis s'arrêterait en position d'équilibre (ou oscillerait). En inversant le sens du courant à chaque demi-tour, le collecteur inverse le sens des forces au bon instant : le couple exercé sur le rotor garde toujours le même sens, ce qui entretient une rotation continue.
\item $I = \dfrac{P}{U} = \dfrac{60}{12} = 5~\text{A}$. Le moteur est traversé par un courant d'intensité $\boxed{I = 5~\text{A}}$.
\end{enumerate}
\textbf{Barème indicatif (sur 10) :} question 1 : 2,5 pts (0,5 par élément) ; question 2 : 2,5 pts ; question 3 : 2 pts ; question 4 : 3 pts (relation 1 pt, application numérique 1 pt, résultat avec unité 1 pt).

\textbf{Points de vigilance :} ne pas confondre stator et rotor ; pour la question 3, la justification attendue doit mentionner explicitement l'\emph{inversion du sens du courant à chaque demi-tour} ; pour le calcul, citer la relation $P = U \times I$ avant l'application numérique et ne pas oublier l'unité (\si{A}).
\end{corrige}

\begin{corrige}[Exercice 10 — Type BEPC : le groupe électrogène et le moteur du moulin]
\begin{enumerate}
\item Le moteur électrique convertit l'\textbf{énergie électrique} en \textbf{énergie mécanique} de rotation (avec une partie perdue en chaleur).
\item Chaîne énergétique complète :
\[ \text{Énergie chimique (carburant)} \rightarrow \text{Énergie mécanique (moteur thermique du groupe)} \rightarrow \text{Énergie électrique (alternateur)} \rightarrow \text{Énergie mécanique (moteur du moulin)} \rightarrow \text{mouture du maïs}. \]
À chaque étape, une partie de l'énergie est perdue sous forme de chaleur.
\item $I = \dfrac{P}{U} = \dfrac{2\,200}{220} = 10~\text{A}$. Le moteur absorbe $\boxed{I = 10~\text{A}}$.
\item L'énergie consommée est $E = P \times t$, avec $P = 2\,200~\text{W} = 2,2~\text{kW}$ et $t = 4~\text{h}$ :
\[ E = 2,2 \times 4 = 8,8~\text{kWh}. \]
Le moteur consomme $\boxed{E = 8,8~\text{kWh}}$ par jour.
\item Le moteur chauffe car une partie de l'énergie électrique est convertie en chaleur par effet Joule dans les bobinages et par frottements ; après une longue journée, cette chaleur s'accumule. Pour éviter sa détérioration : ménager des pauses de refroidissement, assurer une bonne ventilation du moteur (ne pas obstruer les ouïes), ou choisir un moteur de puissance adaptée à la charge pour éviter la surcharge.
\end{enumerate}
\textbf{Barème indicatif (sur 10) :} question 1 : 1 pt ; question 2 : 2,5 pts ; question 3 : 2 pts ; question 4 : 2,5 pts (conversion en kW exigée) ; question 5 : 2 pts.

\textbf{Points de vigilance :} dans la chaîne énergétique, préciser la nature de l'énergie à chaque étape (chimique, mécanique, électrique) ; pour l'énergie, convertir la puissance en \si{kW} avant de multiplier par le temps en heures ; citer l'effet Joule dans l'explication de l'échauffement.
\end{corrige}

\begin{corrige}[Exercice 11 — Type BEPC : investigation sur le sens de rotation]
\begin{enumerate}
\item Schéma du circuit (générateur, interrupteur et moteur en série) :
\begin{popfigure}
\begin{tikzpicture}[scale=1, every node/.style={font=\small}]
% fils du circuit
\draw[thick, popdark] (0,0) rectangle (6,3);
% générateur
\draw[thick, popblue] (0,1.2) -- (0,1.8);
\draw[thick, popblue] (-0.15,1.5) circle (0.3);
\node[popblue] at (-0.75,1.5) {$+$};
\node[popblue] at (-0.75,1.05) {$-$};
\node[popblue, below] at (0,0.9) {Générateur \SI{6}{V}};
% interrupteur
\draw[thick, poporange] (2.4,3) -- (3.0,3);
\fill[poporange] (2.4,3) circle (0.06);
\fill[poporange] (3.6,3) circle (0.06);
\draw[thick, poporange] (3.6,3) -- (3.0,3.5);
\node[poporange, above] at (3,3.5) {Interrupteur};
% moteur
\draw[thick, popgreen] (6,1.5) circle (0.5);
\node[popgreen] at (6,1.5) {\textbf{M}};
\node[popgreen, below] at (6,0.85) {Moteur};
\end{tikzpicture}
\legende{Circuit en série : générateur de \SI{6}{V}, interrupteur et moteur à courant continu M.}
\end{popfigure}
\item \textbf{Essai 1 :} le moteur tourne dans un sens (par exemple sens horaire). \textbf{Essai 2 :} en inversant les deux bornes, on inverse le sens du courant dans le moteur : le moteur tourne \textbf{en sens inverse}. Ces observations mettent en évidence la loi qualitative suivante : \emph{le sens de rotation d'un moteur à courant continu dépend du sens du courant qui le traverse} (et, de façon générale, du sens du champ magnétique de l'inducteur).
\item Quand l'arbre est bloqué, le moteur ne convertit plus l'énergie électrique en énergie mécanique : presque toute l'énergie reçue est transformée en chaleur par effet Joule dans les bobinages, et le courant absorbé devient très supérieur au courant normal. L'échauffement excessif peut détruire l'isolation des bobines et griller le moteur, voire provoquer un incendie.
\item \textbf{Conclusion :} le moteur électrique à courant continu fonctionne grâce aux forces électromagnétiques qu'exerce le champ magnétique du stator (aimant) sur les bobines du rotor parcourues par un courant ; le collecteur inverse le courant à chaque demi-tour pour entretenir la rotation. Le sens de rotation est déterminé par le \textbf{sens du courant} et le \textbf{sens du champ magnétique}.
\item Exemples : le \textbf{ventilateur} (énergie utile : énergie mécanique de rotation des pales, qui brassent l'air) ; la \textbf{pompe à eau} (énergie mécanique pour refouler l'eau) ; on peut aussi citer le moulin à maïs, le fer à repasser... non, celui-ci est thermique : préférer le \textbf{mixeur} (rotation des lames) ou la \textbf{machine à coudre} électrique.
\end{enumerate}
\textbf{Barème indicatif (sur 10) :} question 1 : 2 pts ; question 2 : 2 pts ; question 3 : 2 pts ; question 4 : 2,5 pts ; question 5 : 1,5 pt.

\textbf{Points de vigilance :} le schéma doit comporter les trois dipôles en série avec leurs symboles normalisés ; pour l'essai 3, la justification attendue doit citer l'\emph{effet Joule} et l'absence de conversion mécanique ; dans la conclusion, citer explicitement les deux grandeurs qui fixent le sens de rotation : sens du courant et sens du champ magnétique.
\end{corrige}

\newpage

\coursec{Fiche bilan}

\begin{popfigure}
\begin{center}\tcbox[colback=popdark,colframe=popdark,coltext=white,arc=9pt,boxsep=4pt,left=6pt,right=6pt]{\bfseries\small Le moteur électrique}\end{center}
\vspace{-2pt}
\begin{tcbraster}[raster columns=3,raster equal height=rows,raster column skip=6pt,raster row skip=6pt,enhanced,blankest]
\begin{tcolorbox}[enhanced,colback=white,colframe=popblue,boxrule=0.9pt,arc=3pt,title={Vie courante},fonttitle=\bfseries\scriptsize,coltitle=white,colbacktitle=popblue,left=4pt,right=4pt,top=2pt,bottom=2pt,boxsep=1pt,toptitle=1.5pt,bottomtitle=1.5pt,halign=left]
\begin{itemize}[nosep,leftmargin=*,itemsep=1pt,font=\scriptsize]
\item ventilateur, perceuse
\item pompe, mixeur
\end{itemize}
\end{tcolorbox}
\begin{tcolorbox}[enhanced,colback=white,colframe=poporange,boxrule=0.9pt,arc=3pt,title={Constitution},fonttitle=\bfseries\scriptsize,coltitle=white,colbacktitle=poporange,left=4pt,right=4pt,top=2pt,bottom=2pt,boxsep=1pt,toptitle=1.5pt,bottomtitle=1.5pt,halign=left]
\begin{itemize}[nosep,leftmargin=*,itemsep=1pt,font=\scriptsize]
\item stator (inducteur)
\item rotor (induit)
\item collecteur + balais
\end{itemize}
\end{tcolorbox}
\begin{tcolorbox}[enhanced,colback=white,colframe=popgreen,boxrule=0.9pt,arc=3pt,title={Action magnétique},fonttitle=\bfseries\scriptsize,coltitle=white,colbacktitle=popgreen,left=4pt,right=4pt,top=2pt,bottom=2pt,boxsep=1pt,toptitle=1.5pt,bottomtitle=1.5pt,halign=left]
\begin{itemize}[nosep,leftmargin=*,itemsep=1pt,font=\scriptsize]
\item force de Laplace
\item bobine dans un champ
\end{itemize}
\end{tcolorbox}
\begin{tcolorbox}[enhanced,colback=white,colframe=poppurple,boxrule=0.9pt,arc=3pt,title={Principe},fonttitle=\bfseries\scriptsize,coltitle=white,colbacktitle=poppurple,left=4pt,right=4pt,top=2pt,bottom=2pt,boxsep=1pt,toptitle=1.5pt,bottomtitle=1.5pt,halign=left]
\begin{itemize}[nosep,leftmargin=*,itemsep=1pt,font=\scriptsize]
\item conversion E élec. $\to$ méca.
\item rotation continue
\end{itemize}
\end{tcolorbox}
\begin{tcolorbox}[enhanced,colback=white,colframe=popgold,boxrule=0.9pt,arc=3pt,title={Plaque signalétique},fonttitle=\bfseries\scriptsize,coltitle=white,colbacktitle=popgold,left=4pt,right=4pt,top=2pt,bottom=2pt,boxsep=1pt,toptitle=1.5pt,bottomtitle=1.5pt,halign=left]
\begin{itemize}[nosep,leftmargin=*,itemsep=1pt,font=\scriptsize]
\item $P$, $U$, $I$
\item vitesse (tr/min)
\end{itemize}
\end{tcolorbox}
\begin{tcolorbox}[enhanced,colback=white,colframe=poppink,boxrule=0.9pt,arc=3pt,title={Sécurité},fonttitle=\bfseries\scriptsize,coltitle=white,colbacktitle=poppink,left=4pt,right=4pt,top=2pt,bottom=2pt,boxsep=1pt,toptitle=1.5pt,bottomtitle=1.5pt,halign=left]
\begin{itemize}[nosep,leftmargin=*,itemsep=1pt,font=\scriptsize]
\item terre, fusible
\item pas de surcharge
\end{itemize}
\end{tcolorbox}
\end{tcbraster}

\legende{Carte mentale de la leçon : le moteur électrique (constitution, principe, utilisation et sécurité).}
\end{popfigure}

\begin{bilan}
\textbf{Définitions clés}
\begin{itemize}
\item Un \cle{moteur électrique} est une machine qui convertit l'\cle{énergie électrique} en \cle{énergie mécanique} (mouvement de rotation).
\item Le \cle{stator} (ou inducteur) est la partie fixe : il crée le champ magnétique (aimant ou électroaimant).
\item Le \cle{rotor} (ou induit) est la partie mobile : bobine(s) de fil conducteur montée(s) sur l'arbre.
\item Le \cle{collecteur} et les \cle{balais} assurent le contact électrique avec le rotor et inversent le sens du courant à chaque demi-tour.
\end{itemize}

\textbf{Formules et grandeurs à connaître par cœur}
\begin{center}
\begin{tabularx}{\linewidth}{|l|X|l|}\hline
\rowcolor{popblueL} \textbf{Grandeur} & \textbf{Signification} & \textbf{Unité} \\ \hline
$U$ & tension d'alimentation du moteur & volt (V) \\ \hline
$I$ & intensité du courant absorbé & ampère (A) \\ \hline
$P = U \times I$ & puissance électrique absorbée & watt (W) \\ \hline
$n$ & vitesse de rotation de l'arbre & tr/min \\ \hline
\end{tabularx}
\end{center}

\textbf{Principe de fonctionnement (à savoir rédiger)}
\begin{enumerate}
\item Une bobine parcourue par un courant et placée dans le champ magnétique d'un aimant subit des \cle{forces électromagnétiques} (forces de Laplace).
\item Ces forces forment un \cle{couple} qui fait tourner la bobine (le rotor).
\item Le collecteur inverse le sens du courant à chaque demi-tour : la rotation devient \cle{continue}.
\end{enumerate}

\textbf{Méthodes express}
\begin{itemize}
\item \textbf{Changer le sens de rotation} : inverser le sens du courant (ou celui du champ magnétique).
\item \textbf{Changer la vitesse} : augmenter la tension $U$ augmente la vitesse de rotation.
\item \textbf{Lire une plaque signalétique} : repérer $U$ (V), $I$ (A), $P$ (W), $n$ (tr/min) ; vérifier la compatibilité avec le réseau (\SI{220}{V} chez ENEO).
\item \textbf{Calcul de puissance} : $P = U \times I$. Exemple : $U = \SI{12}{V}$, $I = \SI{2}{A}$ donne $P = \SI{24}{W}$.
\end{itemize}
\end{bilan}

\begin{aretenir}[Checklist avant le BEPC]
\begin{itemize}
\item $\square$ Je sais citer trois appareils courants utilisant un moteur électrique.
\item $\square$ Je sais nommer et situer le stator, le rotor, le collecteur et les balais sur un schéma.
\item $\square$ Je sais que le moteur convertit l'énergie électrique en énergie mécanique.
\item $\square$ Je sais expliquer pourquoi une bobine parcourue par un courant tourne dans un champ magnétique.
\item $\square$ Je sais le rôle du collecteur : inverser le courant pour obtenir une rotation continue.
\item $\square$ Je sais comment inverser le sens de rotation d'un moteur à courant continu.
\item $\square$ Je sais lire et exploiter une plaque signalétique ($U$, $I$, $P$, $n$).
\item $\square$ Je sais calculer la puissance absorbée avec $P = U \times I$ et donner le résultat en watts.
\item $\square$ Je sais citer deux règles de sécurité : mise à la terre, protection par fusible, ne pas surcharger.
\item $\square$ Je sais légender un schéma de moteur électrique à courant continu.
\end{itemize}
\end{aretenir}

\findecours

\end{document}
